\documentclass[11pt,a4paper]{article}
\usepackage[margin=29mm]{geometry}
\usepackage{fontspec}
\setmainfont{TeX Gyre Pagella}
\setmonofont[Scale=MatchLowercase]{DejaVu Sans Mono}
\usepackage{amsmath,amssymb,amsthm,mathtools}
\usepackage{unicode-math}
\setmathfont{TeX Gyre Pagella Math}
\setmathfont[version=bold]{TeX Gyre Pagella Math}
\usepackage{microtype}
\usepackage{longtable,booktabs,array,calc}
\usepackage{needspace}
\usepackage{enumitem}
\usepackage{fvextra}
\DefineVerbatimEnvironment{Highlighting}{Verbatim}{breaklines,fontsize=\small}
\RecustomVerbatimEnvironment{verbatim}{Verbatim}{breaklines,fontsize=\small}
\usepackage{xurl}
\usepackage[unicode,colorlinks=true,linkcolor=blue!45!black,urlcolor=blue!45!black,
  citecolor=blue!45!black]{hyperref}
\usepackage{xcolor}
\hypersetup{pdftitle={Odd prime-factor parity in shifted primes},pdfauthor={Igor Inozemtsev}}
\setlength{\emergencystretch}{2em}
\setlength{\parindent}{1.2em}
\setlength{\parskip}{0.25em}
\setlist{itemsep=0.2em,topsep=0.4em}
\setcounter{tocdepth}{2}
\allowdisplaybreaks[1]
\providecommand{\tightlist}{\setlength{\itemsep}{0pt}\setlength{\parskip}{0pt}}
\providecommand{\passthrough}[1]{#1}
\newcommand{\lean}[1]{{\small\nolinkurl{#1}}}
\Urlmuskip=0mu plus 1mu
\title{Odd prime-factor parity in shifted primes\\[0.3em]
  \large A proof of Green's Problem 64}
\author{Igor Inozemtsev}
\date{12 November 2026}

\begin{document}
\maketitle
\begin{abstract}
We prove that infinitely many primes $p$ satisfy $\Omega(p-2)\equiv1\pmod2$,
where prime factors are counted with multiplicity. The proof constructs a
nonnegative weight of prescribed parity and positive mass. Correlation
estimates justify a weighted sieve, while a bound uniform in the width of a
balanced factor range controls the remaining composites. Subtracting these
composites leaves positive prime mass and supplies a prime above every bound.
Detailed proofs of the weight, correlation, congruence and rectangle estimates are included;
appendices identify their formal counterparts.
\end{abstract}

\noindent\textit{Mathematical sources.}
During the final stage of the Lean formalization, the agent found Jizhou
Guo's manuscript~\cite{guo} on GitHub and used its exposition as a guide.
Guo did not participate in the project or provide personal assistance.
The final development follows the specialization
of its affine-prime-form argument to $N(u)=u+2$, including the choice
$r_0=20$, $\sigma=1/44$, the arithmetic side conditions and the order of
analytic parameter choices. The parity-weight, correlation and sieve
strategy is based on the two OpenAI manuscripts~\cite{openai-e,openai-s}.
Guo's broader statements concerning other shifts, both parities and
squarefree values are outside the scope of the two public Lean theorems here.

The OpenAI collection's
\href{https://github.com/openai/math/blob/adc7f1241b42e322a6451854ab7e4b4c146bf78a/README.md}{README at the cited revision}
describes results at different stages of verification and warns that some
unformalized results may contain errors. No independent human verification
of these source manuscripts is claimed here. These citations record the
origin of the strategy. The supporting inputs needed for the two Green64
statements are proved in the included Lean development. Its machine checks
do not constitute independent human review of the source manuscripts or
establish all of their claims.

\noindent\textit{Formalization and production.}
The proof and formalization were developed with gpt-6 astra in Codex.
This article is derived from the mathematical
Markdown exposition of the Lean development. The statement is the fixed-shift
question in Green's collection~\cite[Problem 64]{green}. No claim of priority
or a comprehensive literature survey is made. The accompanying verification
records specify the checked Lean version, statements and axiom audit.
Those checks and the mathematical exposition are separate parts of the
verification; this article does not assert an independent human review.

\clearpage
\begingroup
\setlength{\parskip}{0pt}
\tableofcontents
\endgroup

\clearpage
\hypertarget{proof}{%
\section{Green's Problem 64}\label{proof}}

\textbf{Theorem.} There are infinitely many primes \(p\) such that
\(\Omega(p-2)\) is odd, where \(\Omega(n)\) counts prime factors with
multiplicity.

This answers Problem 64 in
\href{https://people.maths.ox.ac.uk/greenbj/papers/open-problems.pdf}{Ben
Green's list}, p.~31. We prove the assertion for the fixed shift
\(p-2\). Throughout the argument, \(\lambda(n)=(-1)^{\Omega(n)}\) for
positive integers \(n\).

The proof constructs a nonnegative weight supported on integers \(u\)
with \(\lambda(u)=-1\), and shows that its sum over prime values of
\(u+2\) is positive in every sufficiently large multiplicative interval.
The central analytic input is the shifted-correlation estimate (B18),
which gives the Type II estimate (B3) uniformly in its bounded companion
coefficient. Its proof transfers cancellation on independent prime
coordinates to arithmetic histories, retaining the signed comparison and
accounting for repeated labels. A weighted sieve then removes most
composites; a separate estimate bounds those whose two prime factors are
close in size. The constants in the bound for the remaining composites
are chosen before the width of the balanced range; this permits that
last contribution to be subtracted.

The argument is divided into \protect\hyperlink{weights}{weights and
their mass}, \protect\hyperlink{correlation}{correlation and bilinear
estimates}, and the \protect\hyperlink{sieve}{sieve estimates}. The
present file defines the objects used in the conclusion and derives the
theorem from the two quantitative sieve estimates.

\Needspace{6\baselineskip}

\hypertarget{proof-weights-with-a-prescribed-parity}{%
\subsection{Weights with a prescribed
parity}\label{proof-weights-with-a-prescribed-parity}}

Write \(x=e^L\), \(T=\log L\), and suppose \(L\) is sufficiently large.
Set

\[
r=20,\qquad \sigma=\frac1{44},\qquad s_j=\sigma 2^{-j}L,
\qquad d=\max\{j\ge0:s_j\ge20T\}.
\]

For \(0\le j\le d\) let

\[
Q_j=\{q\text{ prime}:s_j\le\log q\le\tfrac65s_j\}.
\]

Also take every prime interval

\[
P_j=\{p\text{ prime}:\tfrac32s_j\le\log p\le\tfrac95s_j\}
\]

whose logarithmic interval is contained in either \([L^{1/10},L^{1/5}]\)
or \([L^{3/10},L^{2/5}]\). Let \(\mathcal G\) be this set of indices and
\(\mathcal P=\bigcup_{g\in\mathcal G}P_g\). All the \(Q_j\) and \(P_g\)
are pairwise disjoint. Put

\[
V_j=\sum_{q\in Q_j}\frac1q,\qquad M_g=\sum_{p\in P_g}\frac1p.
\]

The bounds used below hold uniformly for large \(L\):

\[
\frac1{12}\le V_j\le4,\qquad
\nu\le M_g\le4,\qquad \nu=\frac14\log\frac{12}{11}>0.
\]

For a positive integer \(u\), separate its prime factors as
\(u=u_*u_{\mathcal P}\), where \(u_*\) has no prime factor in
\(\mathcal P\) and \(u_{\mathcal P}\) has no prime factor outside
\(\mathcal P\). Define

\[
\omega_g(u)=\#\{p\in P_g:p\mid u\},\qquad
W(u)=\prod_{g\in\mathcal G}
       \frac{\omega_g(u)2^{-(\omega_g(u)-1)}}{M_g},
\]

with a zero factor when \(\omega_g(u)=0\). Notice the distinction
between \(\omega_g\), which counts distinct primes, and \(\Omega\),
which counts multiplicities.

Let \(a(v)\) be zero unless all prime factors of \(v\) lie in the
\(Q_j\) and \(v\) has exactly \(20\) factors, with multiplicity, in each
\(Q_j\). On that support set

\[
a(v)=\prod_{q\mid v}\frac1{v_q(v)!}.
\]

Equivalently, \(a(v)\) counts ordered choices of twenty primes in each
\(Q_j\) with product \(v\), divided by \((20!)^{d+1}\). Repeated primes
are allowed. Define

\[
E(u)=(-1)^{\sum_{p\in\mathcal P}v_p(u)},\qquad
A_\rho(u)=W(u)a(u_*)\frac{1+\rho E(u)}2
\quad(\rho\in\{-1,1\}).
\]

We write \(A_0(u)=W(u)a(u_*)\) for the unsigned weight; the subscript
zero is a separate notation, outside the two-sign definition above.

Since \(20(d+1)\) is even, \(a(u_*)>0\) implies \(\lambda(u_*)=1\).
Consequently

\[
A_\rho(u)>0\quad\Longrightarrow\quad\lambda(u)=\rho. \tag{1}
\]

This identity uses prime multiplicities throughout; it does not require
\(u\) to be squarefree.

Fix a smooth function \(\Psi:\mathbb R\to[0,1]\), supported in \([1,2]\)
and positive at some point of \((1,2)\). Such a function is obtained by
translating and rescaling a smooth bump function. Set

\[
w_\rho(u)=A_\rho(u)\Psi(u/x),\qquad
X_\rho=\sum_{u\ge1}w_\rho(u).
\]

Set \(w_\rho(0)=0\), consistently with \(\Psi(0)=0\). All these sums are
finite because \(w_\rho(u)\) vanishes outside \([x,2x]\). The mass
estimate proved in
\protect\hyperlink{weights-localization}{Section~\ref{weights-localization}}
gives a constant \(c_0>0\), depending on \(\Psi\) but not on \(L\) or
\(\rho\), such that

\[
X_\rho\ge c_0xL^{-221}>0 \tag{2}
\]

for both signs and all sufficiently large \(L\).

\Needspace{6\baselineskip}

\hypertarget{proof-the-two-sieve-estimates}{%
\subsection{The two sieve
estimates}\label{proof-the-two-sieve-estimates}}

Let

\[
\mathfrak S=2\prod_{p\ge3,\ p\text{ prime}}
                  \left(1-\frac1{(p-1)^2}\right).
\]

The product converges, and \(\mathfrak S\ge1\); the elementary product
bound is proved in
\protect\hyperlink{sieve-local-densities}{Section~\ref{sieve-local-densities}}.
Write

\[
\mathcal M_\rho(L)=\sum_{u+2\text{ prime}}w_\rho(u)
\]

and, for \(0<\kappa<1/50\), write

\[
\mathcal C_\rho(L,\kappa)
 =\sum_{\substack{u+2\text{ composite}\\
                   P^-(u+2)>x^{1/2-\kappa}}}w_\rho(u),
\]

where \(P^-(n)\) is the least prime factor of \(n>1\).

The weighted sieve and the subtraction of composites outside the
balanced range give, for every fixed \(0<\kappa<1/50\) and all
sufficiently large \(L\),

\[
\mathcal M_{-1}(L)+\mathcal C_{-1}(L,\kappa)
 \ge\frac{\mathfrak S}{2}\frac{X_{-1}}L. \tag{3}
\]

The negative sign is the one needed here. The complete subtraction,
including the order of the approximation parameters, is given in
\protect\hyperlink{sieve-removing-the-unbalanced-composites}{Section~\ref{sieve-removing-the-unbalanced-composites}}.

The bound for the remaining composites has the following uniform form:
there exists \(C_{\mathrm{bal}}>0\) such that, for every
\(0<\kappa<1/50\) and every \(\varepsilon>0\), both signs satisfy

\[
\mathcal C_\rho(L,\kappa)
 \le (C_{\mathrm{bal}}\kappa+\varepsilon)\frac{X_\rho}L
 \tag{4}
\]

once \(L\) is sufficiently large. The threshold may depend on \(\kappa\)
and \(\varepsilon\); \(C_{\mathrm{bal}}\) does not. Its proof in
\protect\hyperlink{sieve-balanced-composites}{Section~\ref{sieve-balanced-composites}}
uses an upper bound for positive factor pairs in a rectangle, followed
by a count of compatible dyadic rectangles. It includes pairs with equal
prime factors.

\Needspace{6\baselineskip}

\hypertarget{proof-positive-prime-mass-and-infinitude}{%
\subsection{Positive prime mass and
infinitude}\label{proof-positive-prime-mass-and-infinitude}}

Choose

\[
\kappa=\min\left\{\frac1{100},
                  \frac{\mathfrak S}{8C_{\mathrm{bal}}}\right\},
\qquad \varepsilon=\frac{\mathfrak S}{8}.
\]

Then \(0<\kappa<1/50\) and
\(C_{\mathrm{bal}}\kappa+\varepsilon\le\mathfrak S/4\). For this fixed
choice, intersect the eventual ranges in (2), (3), and (4). Subtracting
(4) from (3) gives

\[
\mathcal M_{-1}(L)
 \ge\frac{\mathfrak S}{4}\frac{X_{-1}}L>0. \tag{5}
\]

Every summand in \(\mathcal M_{-1}\) is nonnegative. Therefore some
positive integer \(u\) satisfies

\[
u+2\text{ is prime},\qquad A_{-1}(u)>0,\qquad\Psi(u/x)>0.
\]

The support of \(\Psi\) gives \(x\le u\le2x\). By (1),
\(\lambda(u)=-1\), so \(\Omega(u)\) is odd. For \(p=u+2\) we have
\(p-2=u>0\) and hence \(\Omega(p-2)\) odd.

Given any natural number \(M\), choose \(L\) in the eventual range of
(5) with \(e^L>M\). The resulting prime satisfies \(p=u+2>M\). Thus the
set of primes in the theorem is unbounded, and therefore infinite.
\(\square\)

\clearpage
\hypertarget{weights}{%
\section{The weights and their mass}\label{weights}}

We use the notation and fixed prime bands from
\protect\hyperlink{proof}{the main proof}. This section proves
nonnegativity, the parity implication, and the uniform lower bound
\(X_\rho\gg xL^{-221}\).

\Needspace{6\baselineskip}

\hypertarget{weights-geometry-and-exact-factorization}{%
\subsection{Geometry and exact
factorization}\label{weights-geometry-and-exact-factorization}}

The inequalities

\[
1<\frac65<\frac32<\frac95<2,\qquad\frac9{10}<1
\]

separate the two intervals at a given dyadic scale from each other and
from the intervals at the next scale. Thus no prime occurs in two bands.
Every \(Q_j\) prime is at least \(e^{20T}=L^{20}\). Every \(P_g\) prime
is at least \(\exp(L^{1/10})\), which exceeds \(L^{20}\) eventually.

The logarithm of the ratio of the endpoints of any one band, measured in
the logarithmic variable, is \(\log(6/5)\). The prime number theorem and
partial summation therefore give

\[
\sum_{e^a\le p\le e^{6a/5}}\frac1p
  =\log(6/5)+o(1)\qquad(a\to\infty).
\]

The error is uniform when \(a\ge20T\). We use the weaker uniform bounds

\[
\frac1{12}\le V_j\le4,\qquad \nu\le M_g\le4,
\qquad \nu=\frac14\log\frac{12}{11}>0.
\]

The containment conditions give \(O(T)\) marked groups and at least one
group for all sufficiently large \(L\). The moving ordinary depth
satisfies \(d+1\le2T\). The containment conditions retain both endpoints
of each whole prime interval.

For \(u>0\) the decomposition \(u=u_*u_{\mathcal P}\) is unique. Its
factors are coprime because their prime supports are disjoint. If
\(a(u_*)>0\), then \(\Omega(u_*)=20(d+1)\). Complete additivity of
\(\Omega\) on positive integers gives

\[
\lambda(u)=(-1)^{20(d+1)}
 (-1)^{\sum_{p\in\mathcal P}v_p(u)}=E(u).
\]

The projection \((1+\rho E(u))/2\) is either zero or one. This proves
\(A_\rho\ge0\) and the implication
\(A_\rho(u)>0\Rightarrow\lambda(u)=\rho\). In particular, repeated
ordinary primes contribute their full multiplicity; replacing \(\Omega\)
by the number of distinct prime divisors would invalidate this step.

For \(k\ge1\), \(k2^{-(k-1)}\le1\). Since \(a(v)\le1\) and
\(M_g\ge\nu\),

\[
0\le w_\rho(u)\le A_\rho(u)\le W(u)a(u_*)
 \le\nu^{-|\mathcal G|}\le L^{C_0}. \tag{W1}
\]

\Needspace{6\baselineskip}

\hypertarget{weights-harmonic-probability-laws}{%
\subsection{Harmonic probability
laws}\label{weights-harmonic-probability-laws}}

Put

\[
H_Q=\prod_{j=0}^d\frac{V_j^{20}}{20!}.
\]

Expanding the product over the ordered ordinary slots gives the exact
identity

\[
\sum_{v\ge1}\frac{a(v)}v=H_Q. \tag{W2}
\]

For one marked group define

\[
Z_g(q)=\prod_{p\in P_g}\left(1+\frac q{p-1}\right),\qquad
H_g=\frac{Z'_g(1/2)}{M_g},\qquad H_P=\prod_gH_g.
\]

A nonempty set of distinct group primes of size \(k\) receives the
factor \(k(1/2)^{k-1}/M_g\). Summing all positive multiplicities of a
prime gives \(\sum_{a\ge1}p^{-a}=1/(p-1)\). Differentiating the finite
product \(Z_g\) therefore proves

\[
H_g=\sum_{\substack{n\ge1\\p\mid n\Rightarrow p\in P_g}}
 \frac{\omega_g(n)2^{-(\omega_g(n)-1)}}{M_gn}
 =\frac{Z_g(1/2)}{M_g}\sum_{p\in P_g}\frac1{p-1/2}.
 \tag{W3}
\]

The zero term for \(n=1\) is understood. All series here are absolutely
convergent. Since the prime sets are disjoint, multiplying (W3) gives

\[
\sum_{n:n_*=1}\frac{W(n)}n=H_P. \tag{W4}
\]

The harmonic bounds and \(\log(1+t)\le t\) give \(Z_g(1/2)\le C\) and
\(\sum_{p\in P_g}(p-1/2)^{-1}\le C M_g\) for an absolute \(C\). The
reverse inequality \((p-1/2)^{-1}\ge p^{-1}\) gives

\[
1\le Z_g(1/2)\le H_g\le C,
\qquad 1\le H_P\le L^{C_1}. \tag{W5}
\]

Define independent random positive integers \(V\) and \(R\) by

\[
\Pr(V=v)=\frac{a(v)}{vH_Q},\qquad
\Pr(R=n)=\frac{W(n)}{nH_P}\quad(n_*=1).
\]

The law of \(V\) can be realized by independent choices in the twenty
slots of each \(Q_j\), each prime \(q\in Q_j\) having probability
\(1/(qV_j)\). The prime-group components of \(R\) are independent as
well.

For a single group, total multiplicity one has unnormalized harmonic
mass \(M_g/M_g=1\). Total multiplicity two arising from two distinct
primes has mass

\[
\frac1{M_g}\sum_{p<q\in P_g}\frac1{pq}
 =\frac{M_g^2-\sum_{p\in P_g}p^{-2}}{2M_g}.
\]

The second expression is bounded below by a positive absolute constant:
\(M_g\ge\nu\) and \(\sum p^{-2}=o(1)\) uniformly. By (W5), division by
\(H_g\) leaves a positive uniform probability for both parities.
Condition on every group except one. Whatever the parity of their
product, the remaining group can supply either total parity. Hence there
is \(\delta>0\) such that

\[
\Pr(E(R)=\rho)\ge\delta\qquad(\rho=\pm1). \tag{W6}
\]

We also need a tail estimate. Take \(t=L^{-2/5}\). Replace each factor
\(1/(p-1)\) in the Euler-product calculation by
\(p^{-1+t}/(1-p^{-1+t})\). Since \(\log p\le L^{2/5}\), we have
\(p^t\le e\). The harmonic sums of these replacement factors are bounded
uniformly, as are the corresponding differentiated finite products.
There are \(O(T)\) groups, so

\[
\sum_{n:n_*=1}\frac{W(n)}{n^{1-t}}\le L^{C_2}.
\]

For every fixed \(\eta>0\), multiplication by \(x^{-\eta t}\) then gives

\[
\sum_{\substack{n>x^\eta\\n_*=1}}\frac{W(n)}n
 \le L^{C_2}\exp(-\eta L^{3/5}). \tag{W7}
\]

Because \(H_P\ge1\), (W7) also bounds \(\Pr(R>x^\eta)\). This estimate
retains all positive prime multiplicities in \(R\).

\Needspace{6\baselineskip}

\hypertarget{weights-localization}{%
\subsection{Localization}\label{weights-localization}}

Set \(U=VR\). The disjoint prime supports make the factorization into
\(V\) and \(R\) unique. Combining the two probability laws gives

\[
X_\rho=H_QH_P\,
 \mathbb E\big[U\Psi(U/x)\,1_{E(R)=\rho}\big]. \tag{W8}
\]

We prove a lower bound of order \(x/L\) for the expectation, uniformly
in \(\rho\).

Choose \(1<\alpha<\beta<2\) and \(c_\Psi>0\) with \(\Psi\ge c_\Psi\) on
\([\alpha,\beta]\). Write

\[
m_j=\frac{11}{10}s_j,\qquad
\Delta=\frac15\sigma L.
\]

The midpoint sum over all ordinary scales is exactly
\(20\sum_{j\ge0}m_j=L\). Keep one \(Q_0\) slot unconditioned. Set
\(k=13\) and \(\epsilon=1/62400\). For each nonreserved slot in a band
\(j<k\), restrict the selected prime by

\[
(11/10-\epsilon)s_j<\log q\le(11/10+\epsilon)s_j.
\]

The numerical conditions for this corridor are

\[
288r=5760\le2^{13},\qquad kr\epsilon=\frac1{240}\le\frac1{120},
\qquad0<\epsilon\le\frac1{10}.
\]

There are finitely many restrictions, and the prime reciprocal-sum
estimate gives each a positive harmonic probability bounded below
independently of \(L\). Their joint probability has the same property.

For large \(L\) we have \(d\ge12\). The largest possible sum of
logarithms in the remaining bands is at most

\[
20\frac65\sigma L\sum_{j\ge13}2^{-j}
 =\frac6{11}2^{-12}L<\frac\Delta{24}.
\]

The infinite midpoint sum over those bands is also less than
\(\Delta/24\). Their difference consequently has absolute value less
than \(\Delta/24\).

Independently require \(E(R)=\rho\) and \(\log R\le\Delta/12\). By (W6)
and (W7), this event has probability at least \(\delta/2\) for large
\(L\). On the intersection, if \(B\) is the logarithm of the product of
all variables except the reserved prime, then

\[
|B-(L-m_0)|\le\frac\Delta6.
\]

The interval

\[
[L+\log\alpha-B,\ L+\log\beta-B]
\]

lies inside the logarithmic interval of \(Q_0\), at distance comparable
to \(L\) from its endpoints. Its length is the fixed positive number
\(\log(\beta/\alpha)\). The prime number theorem and partial summation
give harmonic probability at least \(c/L\) for the reserved prime to lie
there, uniformly in the conditioned variables. On this event
\(U\Psi(U/x)\ge\alpha c_\Psi x\).

The probabilities of the restrictions on the other variables have a
positive product independent of \(L\). Applying (W8) proves

\[
X_\rho\ge c\frac{xH_QH_P}L
\qquad(\rho=\pm1). \tag{W9}
\]

Finally, use \(200^{-20}\le (1/12)^{20}/20!\), together with
\(\log200\le11/2\). Thus \(V_j\ge1/12\) and \(d+1\le2T\) give

\[
H_Q\ge\left(\frac{(1/12)^{20}}{20!}\right)^{d+1}
 \ge200^{-20(d+1)}\ge L^{-220}.
\]

Together with \(H_P\ge1\) and (W9), this yields
\(X_\rho\ge c_0xL^{-221}\), as required.

\clearpage
\hypertarget{correlation}{%
\section{Replacing prime and roughness tests}\label{correlation}}

We replace prime and roughness indicators by their complete-cell
averages. The weights \(A_\rho\), \(x=e^L\), and \(X_\rho\) are those of
\protect\hyperlink{proof}{the main proof}. All constants are fixed
before taking \(L\to\infty\). In particular, a constant attached to a
finite partition is not allowed to vary with \(L\).

The estimates concern the particular prime and roughness discrepancies
defined below. The other coefficient in a bilinear sum may be an
arbitrary bounded complex function. This is the property needed when
replacing the second of two factors.

\Needspace{6\baselineskip}

\hypertarget{correlation-complete-logarithmic-cells}{%
\subsection{Complete logarithmic
cells}\label{correlation-complete-logarithmic-cells}}

Put \(W=\exp(\sqrt L)\) and \(\Delta=L^{-S}\), where \(S>0\) will be
chosen. For \(j\ge0\) define

\[
C_j=(e^{j\Delta},e^{(j+1)\Delta}]\cap\mathbb N.
\]

Every integer \(n\ge2\) belongs to the unique cell with
\(j=\lceil\log n/\Delta\rceil-1\). Define the roughness indicator

\[
R_y(n)=1_{n\ge2,\ P^-(n)>y}.
\]

For either \(I(n)=R_y(n)\), with \(y\ge W\), or
\(I(n)=1_{n\text{ prime}}\), define its approximation by

\[
B_I(n)=R_W(n)\,
 \frac{\sum_{m\in C_j}I(m)}{\sum_{m\in C_j}R_W(m)}
 \quad(n\in C_j),\qquad \alpha_I(n)=I(n)-B_I(n). \tag{C1}
\]

A zero denominator is assigned the value zero. On a cell where \(I\) is
supported on the \(W\)-rough integers, a zero denominator also forces
the numerator to vanish. Consequently

\[
\sum_{n\in C_j}\alpha_I(n)=0. \tag{C2}
\]

Restricting a coefficient to a factor interval means multiplying it by
that interval's indicator. The ratio in (C1) continues to use the whole
cell, including when the interval cuts one or both endpoint cells. This
convention makes restriction to a smaller interval an exact identity.

For the principal character, (C2) handles complete cells. On each such
cell, subtract the value of \(n^{-1+it}\) at one endpoint. The
derivative bound

\[
\left|\frac{d}{dn}n^{-1+it}\right|\le\frac{1+|t|}{n^2}
\]

bounds the variation. At most two endpoint cells remain. Since
\(q\le W\), the principal character is one on every supported integer.
The resulting bound is

\[
\left(8(e^\Delta-1)+\frac2a\right)(2+|t|).
\]

For a nonprincipal character, write each \(y\)-rough integer uniquely as
\(n=pd\), where \(p=P^-(n)\). For fixed \(d\), the prime \(p\) lies in
the interval obtained from \(a\le pd\le b\), \(p>y\), and
\(p\le P^-(d)\); when \(d=1\) the last condition is omitted. This is a
single clipped interval, even when \(p\) divides \(d\). Siegel--Walfisz
and partial summation bound its prime character sum with weight
\(p^{-1+it}\) by \(K(2+|t|)(\log y)^{-D'}\) for any requested \(D'\),
uniformly when the modulus is at most a fixed power of \(\log y\). The
outer cofactor costs only

\[
\sum_{d\le b}\frac1d\le1+\log b.
\]

Because \(\log y\ge\sqrt L\), \(b\le e^{2L}\), and \(q,|t|\le L^B\),
choosing \(D'\) sufficiently large gives any fixed inverse power of
\(L\) for the rough character sum. On each cell slice, the proxy is a
scalar in \([0,1]\) times the \(W\)-rough indicator, so the same
estimate applies. There are at most \(2L^{S+1}+2\) slices. Requesting
saving \(L^{-(B+S+4)}\) for the rough sums makes the total
test-minus-proxy sum at most \(L^{-B}\).

Combining the two character cases gives the following precise statement:
for fixed \(B>0\), \(S>2B\), and \(\beta>0\), eventually, uniformly for

\[
2\le a\le b\le2a,\quad e^{\beta L}\le a,\quad b\le e^{2L},\quad
1\le q\le L^B,\quad |t|\le L^B,\quad y\ge W,
\]

and every Dirichlet character \(\chi\) modulo \(q\),

\[
\left|\sum_{a\le n\le b}\alpha_{R_y}(n)\chi(n)n^{-1+it}\right|
\le L^{-B}. \tag{C3}
\]

For the prime discrepancy, use \(y=2\sqrt a\): every complete cell
meeting \([a,b]\) lies in a range where \(y\)-roughness is equivalent to
primality. The two proxies therefore agree on the complete cells,
including the two cells cut by the interval endpoints.

\Needspace{6\baselineskip}

\hypertarget{correlation-the-bilinear-estimate-used-in-the-proof}{%
\subsection{The bilinear estimate used in the
proof}\label{correlation-the-bilinear-estimate-used-in-the-proof}}

For finite sets of positive integers \(K,J\), set

\[
\mathcal T_\rho(K,J;\alpha,\beta)
 =\sum_{0\le u\le Y}w_\rho(u)
   \sum_{\substack{m\in K,\ n\in J\\mn=u+2}}\alpha(m)\beta(n),
\qquad Y=\lfloor2x\rfloor. \tag{C4}
\]

All pairs are ordered, and \(m=n\) is retained. In applications
\(\alpha\) is a restricted discrepancy (C1), while \(|\beta|\le1\). The
\protect\hyperlink{bilinear-the-statement-on-one-factor-rectangle}{bilinear
chapter} states the estimate with all factor ranges and quantifiers: on
an admissible dyadic rectangle it bounds (C4) by \(CxL^{-D}\) for any
prescribed \(D\), with the cell mesh and \(C\) chosen before the
interval endpoints and the bounded companion coefficient.

Here is the parameter choice used to apply it. Fix \(0<b_1<1/100\) and
\(0<\kappa<1/50\), and put

\[
\beta_P=b_1/12,\qquad \beta=b_1/3,\qquad b_*=b_1/3.
\]

Reserve one ordinary prime slot in a fixed band \(Q_{j_Q}\) so small
that

\[
\frac65\sigma2^{-j_Q}\le\frac{\beta_P}{4}.
\]

Such an index exists because \(2^{-j}\to0\); the moving depth eventually
contains it. Given an error exponent \(D\), choose the analytic
precision \(A\) and the localization exponent \(K_0\) with

\[
D+2+C_E+C_F\le A,\qquad 2A+2C_W+344\le K_0. \tag{C5}
\]

Here \(C_W,C_E,C_F\) are the fixed polynomial exponents respectively
bounding the marked weight, the selected-factor coefficients, and the
Fourier integrals of the ordinary-slot cutoffs; their values are
specified in \protect\hyperlink{bilinear}{the bilinear reduction}.
Choose \(A\) first, then \(K_0\).

The derivation in \protect\hyperlink{bilinear}{Section~\ref{bilinear}}
expands the ordinary slots, allocates a selected divisor
\(e\in[EL^{-88},E]\) with \(E=UL^{-K_0}\), and applies Cauchy--Schwarz.
For the explicitly defined sum \(\mathcal B'\) in (B8), with
\(G=L^{C_E}\prod_i|\widehat f_{\delta_i}(\vartheta_i)|\), it gives

\[
|\mathcal B'|^2\le6EVG^2
 \left(2xL^{2C_W+344}+(C_{\rm phys}+24C_1)xL^{-A}\right). \tag{C6}
\]

The first term is the diagonal divisor-moment bound. The second follows
from the \protect\hyperlink{marked-correlation}{marked correlation} and
its Fourier tail. Since \(EV\le C_{\rm box}xL^{-K_0}\), (C5) implies
\(|\mathcal B'|\le CGxL^{-A}\). The remaining slot integrals cost
\(L^{C_E+C_F}\). Fourier inversion of \(\psi(t)=\Psi(e^t)\) restores the
localization; its zeroth and \(H\)th absolute Fourier moments are
finite, so the resulting constant is independent of \(L\) and of the
bounded companion coefficient.

Apply the unsigned and signed estimates to the exact projection
\(A_\rho=(A_0+\rho A_0E)/2\). The triangle inequality and \(|\rho|=1\)
give the same bound for \(\mathcal T_\rho\).

\Needspace{6\baselineskip}

\hypertarget{correlation-two-successive-replacements}{%
\subsection{Two successive
replacements}\label{correlation-two-successive-replacements}}

Write \(\mathcal T_{ij}\) for (C4) with the first coefficient sharp when
\(i=0\) and approximated when \(i=1\), and similarly for the second
coefficient. Then

\[
\mathcal T_{00}-\mathcal T_{11}
=(\mathcal T_{00}-\mathcal T_{10})
 +(\mathcal T_{10}-\mathcal T_{11}). \tag{C7}
\]

On the cells under consideration, (C1) gives \(0\le B_I\le1\). Both
companions are therefore bounded, so (B3) applies in either order. We
use bounds \(Q\) and \(2Q\), respectively, for the two differences.

A rectangle is partitioned into its attained binary dyadic rectangles.
The finite carrier consists of the scales of actual positive factor
pairs \(mn=u+2\) with \(\Psi(u/x)\ne0\). There are at most

\[
D_Y=(\lfloor\log_2(Y+2)\rfloor+1)^2 \tag{C8}
\]

such pairs of scales. All full-cell ratios remain those of (C1). The
internal subdivision cost is absorbed by the extra \(2\) in (C5). The
external rectangle count (C8) is then paid explicitly. In the special
balanced positive estimate, a sharper linear count is available; that
separate count is used in
\protect\hyperlink{sieve-counting-balanced-rectangles}{Section~\ref{sieve-counting-balanced-rectangles}}.

\Needspace{6\baselineskip}

\hypertarget{correlation-the-finite-error-sum-and-its-quantifiers}{%
\subsection{The finite error sum and its
quantifiers}\label{correlation-the-finite-error-sum-and-its-quantifiers}}

Fix an integer \(N\) and a nondecreasing sequence

\[
t_0=b_1,\qquad t_N=1/2-\kappa.
\]

For \(i<N\), the unbalanced region consists of positive pairs \(mn=u+2\)
with \(x^{t_i}<n\le x^{t_{i+1}}\). Its first sharp coefficient is
\(R_{x^{t_i}}(m)\) and its second is the prime indicator at \(n\). For
the balanced region both coefficients are prime indicators, and both
factors lie in \([x^{1/2-2\kappa},x^{1/2+2\kappa}]\).

Let \(U^{(i)}_{ab,\rho}\) and \(B_{ab,\rho}\) be the corresponding sums
with replacement flags \(a,b\in\{0,1\}\). The total error actually used
is

\[
\mathcal E(L)=\sum_{\rho=\pm1}\left(
 \sum_{i<N}|U^{(i)}_{00,\rho}-U^{(i)}_{10,\rho}|
 +|B_{00,\rho}-B_{10,\rho}|
 +|B_{10,\rho}-B_{11,\rho}|\right). \tag{C9}
\]

A second quantity \(\mathcal E_{\rm all}\) sums all four successive
differences, in both replacement orders, for each of the \(N+1\) regions
and both signs.

The bounds have the following quantifier order:

\[
\begin{gathered}
\text{fix }N,b_1,\kappa,D\text{ and the cutoff }\Psi;\\
\text{choose }H,J,B,S,C\text{ with }S>2(B+1)+3, C>0;\\
\text{then for every admissible sequence }(t_i),
\text{ eventually in }L,\\
\mathcal E(L)\le2(N+3)Q(L),\qquad
\mathcal E_{\rm all}(L)\le12(N+1)Q(L),\\
Q(L)=C D_Y xL^{-D}. \tag{C10}
\end{gathered}
\]

The threshold may depend on the partition. The approximation parameters
and \(C\) precede its endpoints. To check the factors in (C10), for one
sign there are \(N\) unbalanced first steps and two balanced steps with
bounds \(Q\) and \(2Q\). This gives \((N+3)Q\). For each region and sign
the four steps in both orders cost \(Q+2Q+Q+2Q=6Q\).

The conclusion follows from the rectangle estimates by finite partition
and the triangle inequality.

\Needspace{6\baselineskip}

\hypertarget{correlation-normalization-by-the-actual-mass}{%
\subsection{Normalization by the actual
mass}\label{correlation-normalization-by-the-actual-mass}}

Take \(D=230\). Since \(Y=\lfloor2e^L\rfloor\), (C8) gives
\(D_Y=O(L^2)\). For fixed \(N\) and the fixed parameters in (C10),

\[
\mathcal E(L),\mathcal E_{\rm all}(L)
 =O(e^LL^{-228})=o(e^LL^{-222}). \tag{C11}
\]

The lower bound from \protect\hyperlink{weights}{Section~\ref{weights}}
is uniform in the sign:

\[
\frac{X_\rho}{L}\ge c_0e^LL^{-222}.
\]

Given \(\varepsilon>0\), apply (C11) with the absolute error tolerance
\(\varepsilon c_0\) and intersect its eventual range with the mass lower
bound. This proves, for either sign,

\[
\mathcal E(L)=o(X_\rho/L),\qquad
\mathcal E_{\rm all}(L)=o(X_\rho/L). \tag{C12}
\]

In particular, (C7) gives the balanced two-step error required in the
final sieve argument.

\clearpage
\hypertarget{bilinear}{%
\section{The bilinear reduction}\label{bilinear}}

This chapter gives the finite sums and estimates behind the bilinear
replacement in
\protect\hyperlink{correlation}{Section~\ref{correlation}}. We use its
complete-cell discrepancies and the weights from
\protect\hyperlink{proof}{Section~\ref{proof}}. To distinguish parity
from Fourier frequency, \(\rho\in\{\pm1\}\) denotes parity and
\(\tau\in\mathbb R\) denotes a Mellin frequency.

\Needspace{6\baselineskip}

\hypertarget{bilinear-the-statement-on-one-factor-rectangle}{%
\subsection{The statement on one factor
rectangle}\label{bilinear-the-statement-on-one-factor-rectangle}}

Fix \(D>0\), \(\eta>0\), and \(C_{\rm box}\ge1\). Reserve a slot in a
fixed ordinary band \(Q_{j_Q}\) with

\[
\frac65\sigma2^{-j_Q}\le\eta/4.
\]

Choose fixed exponents \(A,K_0\) satisfying

\[
A\ge D+C_E+C_F,\qquad K_0\ge2A+2C_W+344. \tag{B1}
\]

Here \(C_W\ge0\) bounds \(W(u)\) by \(L^{C_W}\), \(C_E=2\log(20!)\), and
\(C_F\) is the fixed Fourier cost defined below. There are integers
\(H,J,B\), a mesh exponent \(S>2(B+1)+3\), and \(C>0\) with the
following property. For each fixed \(\beta,\gamma>0\), eventually in
\(L\), the conclusion holds uniformly in all the data below:

\[
\begin{gathered}
U,V>0,\quad x\le4UV,\quad UV\le C_{\rm box}x,\quad V\ge1,\\
E=UL^{-K_0},\quad1\le E\le x^{1-\eta},\\
2\le a\le b\le2a,\quad x^\beta\le a,\quad b\le x^2,\\
K=[a,b]\cap\mathbb N\subset[U,2U),\qquad
N\subset[V,2V]\cap\mathbb N,\qquad |b_n|\le1.
\end{gathered} \tag{B2}
\]

Let \(\alpha\) be the restriction to \(K\) of either the prime
discrepancy or the \(x^\gamma\)-rough discrepancy on the complete cells
of width \(L^{-S}\). Then

\[
|\mathcal T_\rho(K,N;\alpha,b)|\le CxL^{-D}
\qquad(\rho=\pm1). \tag{B3}
\]

The parameters \(H,J,B,S,C\) are fixed before \(\beta,\gamma\) and the
rectangle endpoints; the eventual threshold may depend on
\(\beta,\gamma\). There is no regularity hypothesis on \(b_n\). An empty
first interval gives zero and may be omitted. Here \(\mathcal T_\rho\)
is (C4), with \(w_\rho(u)= \Psi(u/x)A_\rho(u)\) and \(x=e^L\); thus the
shift is \(mn=u+2\). The constants also depend on the fixed smooth
cutoff \(\Psi\) and the reserved band \(j_Q\). The integer \(H\) is a
Fourier-tail order, \(J\) is the number of additional shared marks,
\(B\) is the requested complete-cell discrepancy order, and \(S\) is its
mesh exponent.

The range allows a companion factor of size \(x^\eta\): if
\(U\asymp x^{1-\eta}\) and \(V\asymp x^\eta\), then
\(E=UL^{-K_0}\le x^{1-\eta}\) eventually. More generally (B2) implies
\(V\ge\tfrac14x^\eta L^{-K_0}\). In the sieve application one fixes
\(\eta=b_1/12\) with \(0<b_1<1/100\), so \(0<\eta<1/1200\). There is no
lower bound on this fixed positive exponent independent of \(b_1\), but
it cannot tend to zero with \(L\) in this statement. Uniformity in
\(b_n\) concerns the entire sum on scale \(x\). For a singleton
\(N=\{n\}\), (B3) still has right side \(CxL^{-D}\), not
\(C(x/n)L^{-D}\); it does not assert a prime-correlation saving at the
individual dilation's natural scale \(x/n\).

The proof first establishes the same conclusion for each of the weights
\(W(u)a(u_*)\varepsilon(u)\), where \(\varepsilon=1\) or
\(\varepsilon=E(u)\). The identity \(A_\rho=(A_0+\rho A_0E)/2\) then
proves (B3).

\Needspace{6\baselineskip}

\hypertarget{bilinear-allocating-ordinary-slots}{%
\subsection{Allocating ordinary
slots}\label{bilinear-allocating-ordinary-slots}}

Order the nonreserved ordinary slots by decreasing logarithmic scale,
breaking ties by their slot labels. Write their number as
\(n=20(d+1)-1\), their scales as \(s_i\), and their prime choices as
\(q_i\). Repeated prime values retain their separate labels. An
allocation \(\delta\in\{0,1\}^n\) assigns the slots with \(\delta_i=1\)
to a product \(e\) and the remaining slots to a product \(r\).

Use two fixed smooth functions \(f_0,f_1\), with values in \([0,1]\),
supported in \([-85,85]\). On \([-84,84]\) they are \(f_1(t)=h(t)\) and
\(f_0(t)=1-h(t)\), where \(h\) is a smooth transition equal to zero on
\((-\infty,0]\) and one on \([1,\infty)\). At slot \(i\), set

\[
R_i=\log E-\sum_{j<i}\delta_j\log q_j,
\qquad
\prod_{i=1}^n f_{\delta_i}\left(\frac{R_i-\log q_i}{s_i}\right)
\tag{B4}
\]

for the allocation weight. The slot geometry makes these weights sum to
one on the relevant tuples and places every contributing selected
product in the interval

\[
EL^{-88}\le e\le E. \tag{B5}
\]

Here is the numerical reason for the lower endpoint. A skipped slot has
\(R_i<\log q_i+s_i\). In every nonterminal band, the next band contains
at least nineteen nonreserved slots; their total logarithmic mass is at
least \((19/2)s_i>(11/5)s_i\ge\log q_i+s_i\). In the last band,
\(s_d<40\log L\), so \(\log q_i+s_i\le88\log L\). Induction on the
remaining slots bounds the final residual by \(88\log L\). A selected
slot has \(R_i>\log q_i\) and therefore keeps the residual nonnegative.
The same prefix bounds keep each used argument inside the region where
\(f_0+f_1=1\); summing the two choices successively proves the partition
identity in (B4).

The initial supply condition is ensured after excluding a large marked
part. On \(u\in[x,2x]\), a marked part at most \(x^{\eta/4}\) and the
reserved prime at most \(x^{\eta/4}\) leave at least \(x^{1-\eta/2}\) in
the remaining ordinary slots, whereas \(E\le x^{1-\eta}\). The Rankin
estimate (W7) bounds the excluded tuples by any prescribed logarithmic
saving. The finite source comparison also retains the Fourier truncation
error; its total, with coefficients of modulus at most one, is bounded
by \(4xL^{-D}\) after the required precision is fixed.

For clarity, the Fourier-separated coefficients can be written
explicitly. Use \(\mathrm e(t)=e^{2\pi it}\) and
\(\widehat f(\xi)=\int f(t)\mathrm e(-t\xi)\,dt\). For fixed
\(\delta,\boldsymbol\vartheta=(\vartheta_i)\), put

\[
\omega_i=-\frac{\vartheta_i}{s_i}
 -\delta_i\sum_{j>i}\frac{\vartheta_j}{s_j}.
\]

The exponent in Fourier inversion of (B4) is
\(\log E\sum_i\vartheta_i/s_i+\sum_i\omega_i\log q_i\). Define the two
finite coefficient sums

\[
\begin{aligned}
b_e&=\sum_{\substack{(q_i)_{\delta_i=1}\\\prod q_i=e}}
        \prod_{\delta_i=1}\mathrm e(\omega_i\log q_i),\\
c_r&=(20!)^{-(d+1)}
      \sum_{\substack{(q_i)_{\delta_i=0}\\\prod q_i=r}}
        \prod_{\delta_i=0}\mathrm e(\omega_i\log q_i),\\
H(v_*)&=\sum_{\substack{p\in Q_{j_Q}\\p\mid v_*}}c_{v_*/p}.
\end{aligned} \tag{B6}
\]

Each prime variable lies in its assigned band. In particular \(H\)
retains the reserved prime and all repeated-prime assignments. For
\(v>0\), define

\[
a_v(v)=v^{i\tau}W(v)\varepsilon(v)H(v_*),
\]

and

\[
a_e(e)=1_{EL^{-88}\le e\le E}\,
\exp(i\tau\log e)\,\mathrm e\!\left(\log E\sum_i\vartheta_i/s_i\right)
\prod_i\widehat f_{\delta_i}(\vartheta_i)\,b_e. \tag{B7}
\]

Here \(e\) in \(\log e\) is the selected positive integer. With this
notation the finite sum to be estimated is

\[
\mathcal B'=\sum_e a_e(e)\sum_{n\in N}b_n
 \sum_{\substack{m\in K\\mn>2,\ e\mid mn-2}}
 \alpha(m)a_v((mn-2)/e). \tag{B8}
\]

This formula uses the actual integer \(u=mn-2\). When a containing
interval is needed after Cauchy--Schwarz, \(u\le\lceil6UV\rceil\)
suffices.

At most \((20!)^{d+1}\) selected assignments can have a prescribed
product. Since \(d+1\le2\log L\),

\[
|a_e(e)|\le G,
\qquad
G=L^{C_E}\prod_i|\widehat f_{\delta_i}(\vartheta_i)|. \tag{B9}
\]

For the complementary sum, the full factorial normalization in (B6)
gives \(|c_r|\le1\). Thus \(|H(v_*)|\le d_2(v_*)\), where \(d_2\) is the
divisor-counting function, and \(|a_v(v)|\le L^{C_W}d_2(v)\).

\Needspace{6\baselineskip}

\hypertarget{bilinear-the-positive-square-and-the-diagonal}{%
\subsection{The positive square and the
diagonal}\label{bilinear-the-positive-square-and-the-diagonal}}

Cauchy--Schwarz in \((e,n)\) gives

\[
|\mathcal B'|^2\le
\left(\sum_{e,n}|a_e(e)b_n|^2\right)
\sum_{e,n}\left|
 \sum_m\alpha(m)1_{e\mid mn-2}a_v((mn-2)/e)\right|^2.
\tag{B10}
\]

The first factor is at most \(6EVG^2\). Enlarge the second, nonnegative
sum to all odd \(e\le2E\) and all \(n\le3V\), using a nonnegative smooth
weight \(\eta_e(e/E)\eta_n(n/V)\) that is at least one on (B5) and
\(n\in[V,2V]\). Its supremum is at most \(2L^{88}\). Here the supports
are part of the construction:

\[
\eta_e(t)\ne0\Longrightarrow \tfrac12L^{-88}<t<2,
\qquad
\eta_n(t)\ne0\Longrightarrow \tfrac12<t<3. \tag{B10a}
\]

For completeness, put \(a=L^{-88}\) and fix a smooth function \(b\) with
\(0\le b\le1\), supported in \((1/2,3/2)\) and equal to one on
\([2/3,4/3]\). The actual majorant is
\(\eta_e(t)=b(t)+\sum_{j\ge0,\,2^ja\le1}b(t/(2^ja))\). The plateaux of
successive scales meet, so this is at least one on \([a,1]\). Every
scale lies in \([a,1]\), which proves (B10a). The number of summands is
at most \(2L^{88}\) for \(L\ge8\), and, for \(K_j=\|b^{(j)}\|_\infty\),
\(\|\eta_e^{(j)}\|_\infty\le2K_jL^{88(j+1)}\). The fixed bump
\(\eta_n\), equal to one on \([1,2]\) and supported in \((1/2,3)\), has
derivative bounds independent of \(L\). Only after this enlargement do
we expand the square into its diagonal and off-diagonal parts.

Write

\[
F_{e,n}(m)=\alpha(m)1_{mn>2,\ e\mid mn-2}a_v((mn-2)/e).
\]

The off-diagonal is the concrete finite sum

\[
\mathcal O=\sum_{\substack{e\le2E\\e\text{ odd}}}\sum_{n\le3V}
 \eta_e(e/E)\eta_n(n/V)
 \sum_{\substack{m,m'\in K\\m\ne m'}}
 F_{e,n}(m)\overline{F_{e,n}(m')}. \tag{B11}
\]

No positivity is used for (B11). Positivity was used only for the
unexpanded square in (B10).

For the diagonal, fix \(u=mn-2\). There are at most \(d_2(u+2)\) choices
of \((m,n)\) and at most \(d_2(u)\) divisors \(e\). The bound after (B9)
therefore gives an unweighted diagonal at most

\[
L^{2C_W}\sum_{1\le u\le\lceil6UV\rceil}d_2(u)^3d_2(u+2).
\]

The elementary divisor-moment estimate is

\[
\sum_{u\le Y}d_2(u)^3d_2(u+2)
 \le(Y+2)(1+\log(Y+2))^{255}. \tag{B12}
\]

Indeed, \(d_2(u)^3d_2(u+2)\le d_2(u)^4d_2(u+2)^4\). Cauchy--Schwarz
bounds the sum by the eighth divisor moment up to \(Y+2\). The pointwise
comparison \(d_2(n)^8\le d_{256}(n)\) and successive harmonic summation
give \(\sum_{n\le N}d_{256}(n)\le N(1+\log N)^{255}\). Since
\(UV\le C_{\rm box}x\), the right side of (B12) is at most \(xL^{256}\)
eventually. Including the smoothing weight gives

\[
\mathcal D\le2xL^{2C_W+344}. \tag{B13}
\]

Thus the remaining estimate needed for (B3) is an arbitrary logarithmic
saving for (B11), with all signed marked factors retained.

\Needspace{6\baselineskip}

\hypertarget{bilinear-the-shifted-correlation}{%
\subsection{The shifted
correlation}\label{bilinear-the-shifted-correlation}}

For a nonzero summand of (B11), both \(mn-2\) and \(m'n-2\) are
divisible by the odd integer \(e\). The congruence \(mn\equiv2\pmod e\)
makes \(n\) invertible modulo \(e\), so \(e\mid m'-m\). Put

\[
k=\frac{m'-m}{e}\ne0,\qquad
z=m'\frac{mn-2}{e},\qquad
z+2k=m\frac{m'n-2}{e}. \tag{B14}
\]

Both physical positions are positive. The support of the smoothing gives
\(|k|\le L^{K_0+89}\) and places the dyadic scales of \(z\) between
\(xL^{-(K_0+90)}\) and \(xL^{K_0+90}\), with a fixed enlargement for the
smooth dyadic partition. These bounds follow by substituting
\(E=UL^{-K_0}\), \(m,m'\asymp U\), and the support bounds on \(e,n\).
Indeed, (B10a) gives \(e>EL^{-88}/2\) and \(n>V/2\), while
\(|m'-m|\le U\). Thus \(|k|<2L^{K_0+88}\le L^{K_0+89}\). Also
\(m,m'\in[U,2U]\), \(e<2E\) and \(n<3V\) give
\(cU^2V/E\le z\le CU^2VL^{88}/E\) for fixed \(c,C>0\); the lower bound
uses \(mn-2\ge mn/2\) eventually. Together with \(UV\asymp x\), these
are the stated polynomial logarithmic bounds, also for negative \(k\).

The coefficient \(\alpha\) is supported on integers with no prime factor
at most \(W=e^{\sqrt L}\). Hence \(m,m'\) are units modulo \(2|k|\)
eventually. The odd quotient condition on \(e\) is the following exact
residue mask, with \(d_k=|k|\):

\[
T_{d_k}=
\begin{cases}\{d_k+1\}\pmod{2d_k},&d_k\text{ even},\\
\varnothing,&d_k\text{ odd}.
\end{cases}
\]

For characters modulo \(2d_k\), define

\[
c_\chi=\frac1{\varphi(2d_k)}\sum_{r\in T_{d_k}}\overline{\chi(r)}.
\]

Character orthogonality gives

\[
1_{m'm^{-1}\in T_{d_k}}
 =\sum_\chi c_\chi\chi(m')\overline{\chi(m)},
\qquad\sum_\chi|c_\chi|\le1. \tag{B15}
\]

The smoothing factors have explicit expressions in the three logarithmic
coordinates \(t_1=\log(m/U)\), \(t_2=\log(m'/U)\) and \(t_3=\log(z/Z)\):

\[
\begin{aligned}
g_e(t)&=\frac{U}{kE}(e^{t_2}-e^{t_1}),\\
g_n(t)&=\frac{Z}{kUV}e^{t_3}(e^{-t_1}-e^{-t_2})
           +\frac{2}{UV}e^{-t_1}.
\end{aligned} \tag{B15a}
\]

They equal \(e/E\) and \(n/V\) respectively by (B14). Multiply
\(\eta_e(g_e(t))\eta_n(g_n(t))\) by fixed smooth buffers in the three
coordinates, each equal to one on the relevant containing interval. This
gives a compactly supported function on a fixed logarithmic box. Since
\(|k|\ge1\), \(U/E=L^{K_0}\) and \(Z/(UV)\) is bounded by a fixed power
of \(L\), the chain rule and the derivative bounds after (B10a) bound
each of its fixed-order derivatives by a fixed power of \(L\). Fourier
inversion of this three-variable function separates \(m,m',z\); the
quantitative truncation is given in
\protect\hyperlink{bilinear-summing-frequencies-and-returning-to-the-bilinear-sum}{the
frequency summation}.

This reduces (B11) to correlations at the shift \(2k\). Here are their
actual coefficient functions. For the residual \(c\) in (B6), write

\[
H_\sigma(v)=\sum_{\substack{p\in Q_{j_Q}\\p\mid v}}p^{i\sigma}c_{v/p},
\]

and, for \(n>0\), put

\[
\begin{aligned}
F(n)&=\sum_{m\mid n_*}\overline{\alpha(m)}\chi(m)
       m^{i(\xi_2-\tau)}H_\sigma(n_*/m),\\
G(n)&=\sum_{m\mid n_*}\alpha(m)\overline{\chi(m)}
       m^{i(\xi_1+\tau)}\overline{H_\sigma(n_*/m)}.
\end{aligned} \tag{B16}
\]

The same original \(\alpha\) occurs at both endpoints, with the
displayed conjugation. If \(m\) supports \(\alpha\), dividing by \(m\)
removes no marked prime. Thus \(W(n/m)=W(n)\) and
\(\varepsilon(n/m)=\varepsilon(n)\) whenever \(m\mid n_*\).
Multiplicativity of imaginary powers then turns the full cofactor
expansion exactly into

\[
\begin{aligned}
\mathcal C_{k,Z}=
\sum_{\substack{z\ge1\\z+2k>0}}
&\psi(z/Z)z^{i(\tau+\xi_3)}(z+2k)^{-i\tau}
 F(z)G(z+2k)\\
&\hspace{4mm}\cdot W(z)W(z+2k)\varepsilon(z)\varepsilon(z+2k).
\end{aligned} \tag{B17}
\]

The smooth function \(\psi\) is supported in \([1/64,64]\), so the sum
is finite. The marked-correlation estimate says that, for any fixed
\(D_1,C_Z,C_k,C_f>0\), the mesh can be chosen so that

\[
|\mathcal C_{k,Z}|\le ZL^{-D_1} \tag{B18}
\]

uniformly for \(0<|k|\le L^{C_k}\), \(xL^{-C_Z}\le Z\le xL^{C_Z}\),
characters modulo \(2|k|\), both choices of \(\varepsilon\), and
\(|\tau|,|\sigma|,|\xi_i|\le L^{C_f}\). The cutoff is smooth with fixed
support and prescribed bounds on its logarithmic derivatives. The
approximation parameters precede the fine interval's positive lower
exponent and the roughness exponent.

The proof of (B18) is given in
\protect\hyperlink{marked-correlation}{Section~\ref{marked-correlation}}.
It combines alternating coordinate averages, their arithmetic character
approximation, a physical moment estimate, and a separate discrepancy
bound for the comparison term. The scope of this assertion matters.
\(\alpha\) is exactly the restricted discrepancy (C1), and \(c_r\) is
exactly the ordinary-slot coefficient (B6), for an arbitrary allocation
and arbitrary slot Fourier frequencies. Only the external frequencies
\(\tau,\sigma,\xi_1,\xi_2,\xi_3\) have the displayed polynomial
restriction. For every fixed \(j_Q,D_1,C_Z,C_k, C_f\) and cutoff
derivative budget, there are common \(J,B,S\); for each fixed
\(\beta,\gamma>0\) there is then one threshold, uniform in the fine
interval, its further restriction, allocation, slot frequencies,
character, external frequencies, shift and scale. The result is not
claimed for an arbitrary bounded first coefficient \(\alpha\).

\Needspace{6\baselineskip}

\hypertarget{bilinear-summing-frequencies-and-returning-to-the-bilinear-sum}{%
\subsection{Summing frequencies and returning to the bilinear
sum}\label{bilinear-summing-frequencies-and-returning-to-the-bilinear-sum}}

For completeness, the losses in passing from (B18) to (B11) can be paid
before choosing the cell mesh. Fix the desired off-diagonal accuracy
\(A\) and let

\[
D_{\rm tail}=A+K_0+200,\qquad
D_{\rm main}=K_0+315+D_{\rm tail}.
\]

Fix an integer \(P>3\), so \((1+\|t\|)^{-P}\) is integrable on
\(\mathbb R^3\), and set

\[
\begin{aligned}
p_0&=5(K_0+90)(P+1),&p_1&=5(K_0+90)(P+2),\\
D_1&=D_{\rm main}+p_0,&
C_f&=p_1+D_{\rm main}+2C_W+30,\qquad T_f=L^{C_f}.
\end{aligned}\tag{B19}
\]

The five factors defining the cutoff \(f=f_{k,Z,L}\) are the two
compositions after (B15a) and the three fixed buffers. They have common
compact support. The chain and product rules give, for each fixed
derivative order \(j\), a bound
\(\|\partial^j f\|_\infty\le C_jL^{5(K_0+90)(j+1)}\). Integration by
parts through orders \(P\) and \(P+1\), respectively, on this fixed
support, therefore gives the two uniform moment bounds

\[
\int_{\mathbb R^3}|\widehat f(\zeta)|\,d\zeta
 \le C_0'L^{p_0},\qquad
\int_{\mathbb R^3}\|\zeta\||\widehat f(\zeta)|\,d\zeta
 \le C_1L^{p_1}.
\tag{B19a}
\]

The constants \(C_0',C_1\) depend on the fixed buffers, not on \(L,k,Z\)
or the fine cells. In particular the zeroth Fourier moment itself need
not be bounded independently of \(L\).

Split Fourier inversion at the ball of radius \(T_f\). On the ball,
apply (B18) with saving \(D_1\), after shifting \(\xi_i\) by \(2\pi\)
times the Fourier variables and increasing the fixed frequency exponent.
The character coefficients have total modulus at most one by (B15). Thus
the integrated principal term at each \(k,Z\) is at most

\[
ZL^{-D_1}C_0'L^{p_0}=C_0'ZL^{-D_{\rm main}}.
\tag{B19b}
\]

There are at most \(O(L^{K_0+89})\) shifts and \(O(L)\) dyadic scales;
their largest scale is \(O(xL^{K_0+90})\). Hence the sum of (B19b) is
\(O(xL^{2K_0+180-D_{\rm main}})\), which is \(O(xL^{-A})\). The extra
\(p_0\) in the requested correlation precision is essential and is
chosen before \(J,B,S\).

Outside the ball, use

\[
\int_{\|\zeta\|>T_f}|\widehat f(\zeta)|\,d\zeta
\le T_f^{-1}\int\|\zeta\||\widehat f(\zeta)|\,d\zeta. \tag{B20}
\]

Finite sums of absolute coefficient values and the derivative bounds for
the scaled cutoffs give a polynomial factor in \(L\); the additional
powers in (B19) make the total tail \(O(xL^{-D_{\rm tail}})\). In this
comparison, a physical pair where a cutoff vanishes is still included in
the finite Fourier integral. Its full Fourier integral is zero, so its
truncated contribution equals minus its tail and is bounded by (B20).
This accounts for the two endpoint cutoff errors without discarding a
signed subfamily.

The resulting bounds are

\[
|\mathcal O_{\rm ball}|\le C_{\rm phys}xL^{-A},\qquad
|\mathcal O-\mathcal O_{\rm ball}|\le24C_1xL^{-D_{\rm tail}},
\]

where \(C_1\) is the fixed coefficient in (B19a), and
\(C_{\rm phys}=250C_{\rm box}(1+2/\log2)(C_0'+132C_1)\) with a fixed
zeroth-moment coefficient \(C_0'\). Hence

\[
|\mathcal O|\le(C_{\rm phys}+24C_1)xL^{-A}. \tag{B21}
\]

Together, (B10), (B13), and (B21) give

\[
|\mathcal B'|^2\le6EVG^2
 \left(2xL^{2C_W+344}+(C_{\rm phys}+24C_1)xL^{-A}\right).
\]

Since \(EV\le C_{\rm box}xL^{-K_0}\), (B1) implies
\(|\mathcal B'|\le C'GxL^{-A}\) for a fixed \(C'\).

Let \(M_f\ge1\) bound the \(L^1\) norms of
\(\widehat f_0,\widehat f_1\). There are at most \(40\log L\)
nonreserved slots. Integrating all slot frequencies and summing the two
choices at each slot costs at most

\[
(2M_f)^n\le L^{C_F},\qquad C_F=40\log(2M_f). \tag{B22}
\]

Thus (B9) and (B1) absorb the total loss \(L^{C_E+C_F}\). Finally write
\(\psi_0(t)=\Psi(e^t)\) and restore the localization by Fourier
inversion. The moments

\[
M_0=\int|\widehat\psi_0(s)|\,ds,\qquad
M_H=\int|s|^H|\widehat\psi_0(s)|\,ds
\]

are finite. The first controls bounded Mellin frequencies; for the tail
choose \(H\) large enough and use the second. With the source-comparison
error from \protect\hyperlink{bilinear-allocating-ordinary-slots}{the
slot allocation}, the resulting constant is bounded by
\(4+C'M_0+20(2\pi)^HM_H\). This proves (B3) once (B18) has been
established.

\clearpage
\hypertarget{marked-correlation}{%
\section{Cancellation in the marked
correlation}\label{marked-correlation}}

We prove the estimate (B18) of
\protect\hyperlink{bilinear}{Section~\ref{bilinear}}. All the sums and
operators below are finite. We use prime Siegel--Walfisz, the
complete-cell discrepancy (C3), the high-frequency prime estimate
\protect\hyperlink{local-energy-the-high-frequency-prime-input}{(U1)},
smooth Fourier inversion, and the displayed finite energy estimates. The
bilinear exponential-sum bound comes from the row sums of an additive
Gram matrix and Cauchy--Schwarz.

Write \(L=\log x\) and \(\mathrm e(t)=e^{2\pi it}\). We retain the
actual marked groups and the multiplicity parity from
\protect\hyperlink{proof}{Section~\ref{proof}}. The BIG groups have
\(\log p\ge L^{3/10}\); the SMALL groups used for omissions have
\(\log p\le L^{1/5}\). Their harmonic masses lie in \([v,v_{\max}]\),
where \(v=\tfrac14\log(12/11)\) and \(v_{\max}=4\). We write these
masses as \(V_g\) in this chapter; they are \(M_g\) in the definition of
the weight. The fixed damping parameter is \(q=1/2\). The number of
groups is \(O(\log L)\). Constants implicit in this chapter may depend
on the fixed accuracy, frequency bounds, and mark counts, but not on
\(L\), the shift, or the endpoints of a fine cell.

The proof estimates a signed difference between the original correlation
and an arithmetic resampling term. A separate use of the fine-cell
discrepancy then makes the resampling term small. There are three
different spaces in this argument: a product of prime alphabets in
\protect\hyperlink{marked-correlation-orthogonal-averages-and-the-alternating-identity}{the
orthogonal-average argument},
\protect\hyperlink{marked-correlation-the-arithmetic-state-space-and-the-local-contraction}{the
residue and memory space}, and
\protect\hyperlink{marked-correlation-from-residue-averages-to-the-integer-interval}{the
integer-position space}. The local contraction is on the second space.
Its relation to the third is a moment identity and an interval-counting
estimate, not an identification of their norms. Throughout the operator
argument write \(s=2k\) for the integer shift in (B17). A forward frame
moves by \(sD\), and a reverse frame by \(-sD\).

\Needspace{6\baselineskip}

\hypertarget{marked-correlation-orthogonal-averages-and-the-alternating-identity}{%
\subsection{Orthogonal averages and the alternating
identity}\label{marked-correlation-orthogonal-averages-and-the-alternating-identity}}

In a BIG group \(P_g\), use the probability law

\[
V_g=\sum_{p\in P_g}\frac1p,
\qquad \mu_g(p)=\frac1{V_gp}.
\]

There are \(M=J+\ell\) labelled coordinates in each group. Here \(J\)
will be fixed later and the application has \(\ell=1\). Precisely, put
\(\mathcal X=\prod_{g\in\mathrm{BIG}}P_g^M\), with probability
\(\mu(x)=\prod_{g,i}\mu_g(x_{g,i})\), and use
\(\langle f,h\rangle_\mu=\sum_x\mu(x)\overline{f(x)}h(x)\). The tests
before symmetrization are arbitrary functions
\(F,H:\mathcal X\to\mathbb C\). Set

\[
(\mathcal SF)(x)=\frac1{(M!)^{G}}
 \sum_{\sigma\in\prod_g S_M}F((x_{g,\sigma_g(i)})_{g,i}),
\qquad f=\mathcal SF,\quad h=\mathcal SH,
\tag{M0}
\]

where \(G=\#\mathrm{BIG}\). Relabelling coordinates preserves \(\mu\),
so Jensen gives \(\|\mathcal SF\|_\mu\le\|F\|_\mu\). The actual kernel
independently averages its source and target permutations: changing
variables by their inverses moves these two averages onto \(F\) and
\(H\), exactly. This identity, rather than a symmetry assumption on the
arithmetic coefficients, licenses (M0). The tests can depend jointly on
all groups. The maps (T14)--(T15) identify them as Fourier coefficients
of the two physical endpoint tests.

Let \(E_i\) average coordinate \(i\). These are commuting orthogonal
projections. The Hoeffding components

\[
H_Sf=\prod_{i\in S}(I-E_i)\prod_{i\notin S}E_i f
\quad(S\subseteq[M])
\]

are orthogonal and sum to \(f\). Symmetry gives equal squared norms to
components of the same degree \(|S|\). Choose disjoint sets of \(m\)
probe coordinates and \(\ell\) trailing coordinates. At degree \(j\),
the fraction of components containing every probe and no trailing
coordinate is

\[
\frac{\binom{M-m-\ell}{j-m}}{\binom Mj}
\le\binom{m+\ell}{m}^{-1}.
\]

Consequently

\[
\left\|\prod_{i\in\mathrm{probe}}(I-E_i)
             \prod_{i\in\mathrm{tail}}E_i f\right\|_2^2
\le\binom{m+\ell}{m}^{-1}\|f\|_2^2. \tag{M1}
\]

Apply this successively to different groups, using Fubini. For
\(\ell=1\) and \(s\) groups the factor is \((m+1)^{-s}\). The actual BIG
groups contain two disjoint blocks \(\mathcal A,\mathcal B\), each of
size at least \(c\log L\), where \(c=(120\log2)^{-1}\). Given \(A_0>0\),
choose a fixed \(m\) so that \(A_0\le c\log(m+1)\). Then each block
gives squared-energy saving \(L^{-A_0}\).

The coupled phase cannot be moved through these coordinate averages. To
retain it exactly, put

\[
R(f,h)=\mathbb E[\overline f\,\Phi h],\qquad |\Phi|=1,
\]

For a fixed positive SMALL divisor \(d\), orientation \(o\in\{1,-1\}\),
and real \(\zeta,\theta\), its precise value is

\[
D(x)=d\prod_{g\in\mathrm{BIG}}\prod_{i=0}^{J-1}x_{g,i},
\qquad
\Phi(x)=D(x)^{io\zeta}\mathrm e(o sD(x)\theta).
\tag{M0a}
\]

Thus \(|\Phi|=1\), and \(\Phi\) ignores the \(\ell\) trailing
coordinates. It need not ignore the probes. Let \(U\) average all
trailing coordinates, and set

\[
E_a=U\prod_{i\in a}E_i,
\qquad P_{\mathcal A}=\sum_{a\subseteq\mathcal A\times[m]}
 (-1)^{|a|}E_a
 =U\prod_{i\in\mathcal A\times[m]}(I-E_i).
\]

Define \(P_{\mathcal B}\) in the same way. Expanding the two products
gives the exact identity

\[
\begin{aligned}
R(Uf,Uh)&-
\sum_{\substack{\varnothing\ne a\subseteq\mathcal A\times[m]\\
                 \varnothing\ne b\subseteq\mathcal B\times[m]}}
 (-1)^{|a|+|b|}R(E_af,E_bh)\\
&=R(P_{\mathcal A}f,Uh)+R(Uf,P_{\mathcal B}h)
  -R(P_{\mathcal A}f,P_{\mathcal B}h).
\end{aligned} \tag{M2}
\]

Both subsets in the comparison sum are nonempty. Each term on the right
has a contracted endpoint; since \(|\Phi|=1\), Cauchy--Schwarz and (M1)
give

\[
|\text{right side of (M2)}|^2
 \le9L^{-A_0}\|f\|_2^2\|h\|_2^2. \tag{M3}
\]

The number of comparison patterns is a fixed power of \(L\): \(m\) is
fixed and there are \(O(\log L)\) groups. We will pay this cost in the
arithmetic approximation to each pattern. It is not subtracted from the
exponent supplied by the already summed identity (M2).

\Needspace{6\baselineskip}

\hypertarget{marked-correlation-implementing-the-averages-by-characters}{%
\subsection{Implementing the averages by
characters}\label{marked-correlation-implementing-the-averages-by-characters}}

For a nonempty collection of probes, let \(D\) be their product under
the product probability law. Partition its logarithm into cells
\(I_j=[e^{jh},e^{(j+1)h})\), put \(\nu_j=\Pr(D\in I_j)\), and retain
cells with \(\nu_j\ge\xi\). The arithmetic kernel is

\[
K(D,X)=\sum_{\nu_j\ge\xi}
 \frac{\mathbf1_{D,X\in I_j}}{\nu_j}
 \sum_{\substack{\chi\ \mathrm{primitive}\\\operatorname{cond}\chi\le Q}}
       \overline{\chi(D)}\chi(X). \tag{M4}
\]

The reference operation is exact matching of the product, not the
constant kernel. If \(\pi(X)=\Pr(D=X)\), put

\[
I_\pi(D,X)=\frac{1_{D=X}}{\pi(X)},\qquad
\sum_D\pi(D)I_\pi(D,X)\phi(D)=\phi(X)
\quad(\pi(X)>0).
\tag{M4a}
\]

Set the kernel to zero when \(\pi(X)=0\); such \(X\) do not occur in the
product law. Applying (M4a) separately in the two probe blocks sets the
auxiliary products equal to the target first-probe and source
second-probe products. The other, freed endpoint coordinates are then
averaged with their original laws. This is the ordinary resampling form
in (M2). In particular, it evaluates the phase at the spliced products;
it does not average that phase independently. The estimate below
compares this matching form with (M4), including its phase and endpoint
tests; it does not assert that their kernels are close in operator norm.

All marked primes exceed \(Q\). Characters in (M4) thus have modulus one
on its support. The cell Gram matrix has diagonal \(\nu_j\). For two
different primitive characters, freeze every prime coordinate but one.
The condition on the product clips that remaining prime to an interval.
The quotient of the two characters induces a nonprincipal primitive
character with conductor at most \(Q^2\). Prime Siegel--Walfisz and
partial summation bound the corresponding entry by

\[
|G_j(\chi,\chi')|\le (K_1/V_g)L^{-B_1}.
\]

The other coordinates have total probability one. Summing over them does
not introduce the number of tuples. There are at most \(Q^2\)
characters. Therefore, if \(Q\le L^{C_q}\) and
\(\xi^{-1}\le L^{C_\xi}\), requesting accuracy \(B_1=B+C_\xi+2C_q\)
gives

\[
\|K\|_{2\to2}\le1+(K_1/v)L^{-B}. \tag{M5}
\]

Different cells give orthogonal blocks of this operator, so their number
does not multiply (M5).

To compare (M4) with ordinary resampling, retain the full phase
\(\mathrm e(\theta'XY)(XY)^{i\zeta}\). On a major arc
\(\theta'=a/r+\beta\), \(r\le Q\), the function \(\mathrm e(aXY/r)\) on
units is a sum of primitive-character modes with coefficient \(\ell^1\)
norm at most \(r\). The two kernels reproduce these modes up to

\[
(2\epsilon+\epsilon^2)r,
\qquad \epsilon=(K_1/v)L^{-B}.
\]

Within two matched logarithmic cells, the remaining phase varies by at
most

\[
Q^4\bigl(|\zeta|+2\pi|\beta|R_0\bigr)(h_A+h_B),
\qquad R_0=e^{(j_A+1)h_A+(j_B+1)h_B}. \tag{M6}
\]

Here \(R_0\) bounds the product of the two cell variables. This follows
by differentiating the phase in the two logarithmic coordinates. The
factor \(Q^4\) bounds the two finite character sums.

On minor arcs, move the two kernels onto the endpoint tests and use
(M5), eventually with norm at most \(2\). The finite integer bilinear
exponential-sum inequality on dyads of sizes \(M_0,N_0\) has spectral
factor bounded by

\[
\frac1r+\frac2{N_0}+\frac{\log(2r)}{M_0}
       +\frac{r\log(2r)}{M_0N_0}. \tag{M7}
\]

For example, for windows of integer lengths \(P\le M_0\), \(Q'\le N_0\),
the bound preceding (M7) is \([P+(1+Q'/r)(P+r\log(2r))]/(M_0N_0)\). Thus
(M7) is at most \(5L^{-D}\) provided

\[
r,N_0\ge L^D,\quad M_0\ge L^D\log(2r),\quad
M_0N_0\ge L^Dr\log(2r). \tag{M8}
\]

Here is a denominator choice that verifies this alternative. Put
\(d_*=2/5\), \(S_D=L^{D+d_*+2}\), and apply Dirichlet approximation with
cutoff \(\lfloor M_0N_0/S_D\rfloor\). It gives a reduced \(a/r\) with

\[
r\le M_0N_0/S_D,\qquad
|\beta|M_0N_0\le S_D/r,\qquad |\beta|\le r^{-2}.
\]

Each nonempty BIG probe product is at least \(e^{L^{3/10}}\) and has
logarithm \(O_m(L^{2/5}\log L)\); in particular
\(\log(2M_0N_0)\le L^{d_*+2}\) eventually. If \(r>Q\ge L^D\), these
bounds give all four inequalities in (M8). If \(r\le Q\), the matched
cells satisfy \(R_0\le4e^{h_A+h_B}M_0N_0\), so \(|\beta|R_0\) has a
fixed polynomial bound. This supplies the major-arc variation estimate
with explicit control of its denominator.

Passing from labelled prime tuples to their products costs at most

\[
\prod_g\frac{k_g!}{V_g^{k_g}}, \tag{M9}
\]

because a product in a group has at most \(k_g!\) ordered
representations. Here \(k_g\le m\), so (M9) is another fixed power of
\(L\). After including it in the requested minor-arc precision, the
source bounds for the raw term and the two-kernel term are,
respectively, \(125L^{-B}\|f\|_2^2\|h\|_2^2\) and
\(2000L^{-B}\|f\|_2^2\|h\|_2^2\). Their squared difference is bounded by
\(4250L^{-B}\|f\|_2^2\|h\|_2^2\).

If the two product ranges meet \(n_A,n_B\) cells, the probability lost
on discarded cells is at most \(2(n_A+n_B)\xi\). To bound their number,
each probe product has logarithm at most \(B_L=mC_G L^{2/5}\log L\),
where \(C_G\log L\) bounds the BIG group count. Consequently

\[
n_A,n_B\le2+B_L/h,\qquad
\#\{\text{coarse dyads for either product}\}
 \le2+B_L/\log2.
\tag{M9a}
\]

These are polynomial bounds for the mesh (M10), uniform in the
conditioned shared coordinates. They apply also when the product range
is further clipped. The number of patterns is at most
\(2^{m(|\mathcal A|+|\mathcal B|)}\), also polynomial. Shared
coordinates are conditioned on and then integrated by Cauchy--Schwarz.
For a requested one-rectangle saving \(B_{\rm save}\), let
\(C_{\rm factor}\) pay for (M9). The following choices suffice:

\[
\begin{gathered}
D\ge B_{\rm save}+2C_{\rm factor},\qquad
C_s\ge D+d_*+2,\\
C_4\ge C_s+1,\quad C_4\ge C_{\rm freq},\\
H_m\ge B_{\rm save}+4C_q+C_4,\qquad
C_\xi\ge B_{\rm save}+H_m+d_*+4.
\end{gathered}
\]

Here \(|\zeta|\le L^{C_{\rm freq}}\) and \(C_q\) bounds the conductor
cutoff. Request precision \(B_{\rm save}+C_q\) in the rational-phase
reproduction to pay for \(r\le L^{C_q}\); the prime Siegel--Walfisz
precision is then increased by \(C_\xi+2C_q\) as in (M5). Any fixed
constants can be absorbed by an extra power and a larger threshold.
After paying the pattern and dyad counts in \(B_{\rm save}\), (M3)--(M9)
give the small local signed-operator estimate with

\[
h=\frac{\log2}{\lceil4L^{H_m}\rceil},\qquad
Q=\lceil4L^{D+1}\rceil,\qquad \xi=L^{-C_\xi}. \tag{M10}
\]

Choose \(C_q>D+1\) after \(D\) and before \(H_m,C_\xi\).

\Needspace{6\baselineskip}

\hypertarget{marked-correlation-the-arithmetic-state-space-and-the-local-contraction}{%
\subsection{The arithmetic state space and the local
contraction}\label{marked-correlation-the-arithmetic-state-space-and-the-local-contraction}}

\Needspace{6\baselineskip}

\hypertarget{marked-correlation-residues-active-marks-and-memory}{%
\subsubsection{Residues, active marks and
memory}\label{marked-correlation-residues-active-marks-and-memory}}

Fix \(N=2\lceil L^{13/25}\rceil\), \(M=J+1\), and finite cutoffs \(B_g\)
for the BIG memory lengths. The estimates below are uniform in these
cutoffs. For a SMALL group put

\[
\mathbb T_g=\prod_{p\in P_g}\mathbb Z/p\mathbb Z,
\qquad
\mathscr S_g=\{(y,a):y\in\mathbb T_g,\ a:[M]\hookrightarrow P_g,
                         \ y_{a_i}=0\ (i\in[M])\}.
\tag{T1}
\]

Here and below \([M]=\{0,\ldots,M-1\}\). Give every element of
\(\mathscr S_g\) mass \(|\mathbb T_g|^{-1}V_g^{-M}\), and write
\(\omega_g(y)=\#\{p:y_p=0\}\). This is Haar probability in \(y\) and
counting measure in the valid ordered marks, with the displayed factor
\(V_g^{-M}\); it is not a probability on the marked fibre. Translation
means \(y_p\mapsto y_p+\Delta\pmod p\) simultaneously for every \(p\).

For BIG primes set

\[
q_i=\begin{cases}\sqrt q,&i=0,N,\\q,&0<i<N,\end{cases}
\quad \eta_i=1-q_i,\quad H_N=\sum_{i=0}^N\eta_i,
\quad b_p=1-H_N/p,\quad \widetilde\mu_g(p)=\frac1{V_gp b_p}.
\tag{T2}
\]

Eventually \(b_p\ge1/2\). A BIG state consists of an active array
\(a_g\in P_g^M\) and an ordered list \(((p_1,r_1),\ldots,(p_j,r_j))\),
\(0\le j\le B_g\), where \(r_i\in\mathbb Z/p_i\mathbb Z\). Its mass is

\[
\left(\prod_{i<M}\widetilde\mu_g(a_{g,i})\right)
 \frac1{j!}\prod_{i=1}^j\widetilde\mu_g(p_i).
\tag{T3}
\]

The residues in this list have \textbf{counting} measure: summing the
mass of \((p,r)\) over \(r\) gives \(p\widetilde\mu_g(p)\), while its
zero residue has mass \(\widetilde\mu_g(p)\). Equal primes and equal
atoms are allowed here. The active BIG array also has no distinctness
condition at this stage. Distinctness is imposed on the appropriate
birth lists in the transition and then globally in the forest expansion.

Let \(\widetilde{\mathscr S}\) be the product of (T1) and these BIG
spaces, with mass \(w\). Let \(\pi\) replace each ordered OLD list by
its multiset, leaving all SMALL states and active arrays unchanged. On
\(\mathscr S=\pi(\widetilde{\mathscr S})\) put

\[
\bar w(X)=\sum_{\pi(x)=X}w(x),\qquad
\mathcal E(F)=\sum_X\bar w(X)|F(X)|^2,
\qquad
(\Pi u)(X)=\frac{\sum_{\pi(x)=X}w(x)u(x)}{\bar w(X)}.
\tag{T4}
\]

A zero-mass fibre is assigned zero. In our nonempty fibres the masses
are positive. Pullback \(\pi^*F=F\circ\pi\) is an isometry,
\(\Pi\pi^*=I\), and Jensen gives
\(\mathcal E(\Pi u)\le\sum_xw(x)|u(x)|^2\). For a multiset with atom
multiplicities \(n_z\), the OLD part of its mass is
\(\prod_z\widetilde\mu_g(p(z))^{n_z}/\prod_z n_z!\). This follows by
counting its \(j!/\prod_z n_z!\) ordered representatives in (T3). In
particular, passing to multisets never discards collisions.

An OLD atom records a previously introduced prime and its current
residue. On a step \(\Delta=o sD\) all surviving residues become
\(r+\Delta\pmod p\). A \emph{promotion} supplies an unshared target
active slot from an OLD occurrence whose translated residue is zero;
distinct promoted slots use distinct OLD indices. A \emph{store} puts an
unshared source active occurrence into OLD, at residue
\(\Delta\pmod p\). Stored and promoted primes must not divide \(D\). The
remaining unshared target slots are new births sampled with
\(\widetilde\mu_g\). The output length must stay below \(B_g\). A
transition is \emph{clean} when it has no stores and no promotions.
Clean transitions translate every OLD occurrence, including repetitions.

For definiteness the active operator is the following finite sum. Choose
the raw pattern or a nonempty pair \((a,b)\) of probe subsets, auxiliary
primes with the original harmonic law \(\mu\), and two independent
uniform permutations in each SMALL and BIG group. Copy the shared active
coordinates, use the divisor \(D\) of (M18b), and sum all legal
store/promotion choices just described. In each BIG group the
coefficient is

\[
\rho_0^{z_g+z'_g}V_g^{-t_g}
 \prod_{i\text{ fresh at the target}}\widetilde\mu_g(a'_{g,i}),
\tag{T5}
\]

where \(z_g,z'_g\) count zero-residue OLD occurrences at the two
endpoints and \(t_g\) is the number of promotions. Choose
\(q^{1/4}<\rho_0<1\). In each SMALL group multiply by
\(V_g^{-1}q^{(\omega_g(y)+\omega_g(y')-2M)/2}\), imposing shared marks
and \(y'=y+\Delta\). Finally multiply by the auxiliary mass, the
permutation masses, and \(K_fD^{i\zeta}\mathrm e(sD\theta)\) in the
forward orientation; the reverse orientation conjugates
\(K_fD^{i\zeta}\) and changes \(\theta\) to \(-\theta\) in the return
phase. There is only one such complex scalar for the full frame.

The operator \(E_o\) is this row sum, pulled back and averaged by (T4),
with the two gates \(|\mathrm{OLD}|\le T=\lceil L^2\rceil\) and the
condition that the list of fresh target numerical primes is distinct.
Only distinctness between different birth times is omitted. These are
the operators denoted \(E_i\) in (F7), with alternating \(o\). The
integer-position matrix (M18c) has a different carrier and will be
related to their full path sum in
\protect\hyperlink{marked-correlation-from-residue-averages-to-the-integer-interval}{the
residue-to-interval transfer}.

\Needspace{6\baselineskip}

\hypertarget{marked-correlation-the-frozen-big-form-including-its-errors}{%
\subsubsection{The frozen BIG form, including its
errors}\label{marked-correlation-the-frozen-big-form-including-its-errors}}

We first identify the harmonic form bounded by (M1)--(M10). Freeze the
SMALL shared tuple and its product \(d>0\). For a pattern \(\nu\), let
\(S_g\subset[J]\) be its shared coordinates. For source and target
arrays \(u,v\) and permutations \(\sigma,\tau\), put

\[
B_\mu(u,v)=\prod_g\left[
 \prod_{i\in S_g}1_{u_{g,\sigma_g(i)}=v_{g,\tau_g(i)}}
 \prod_{i\notin S_g}\mu_g(v_{g,\tau_g(i)})\right].
\tag{T6}
\]

In this formula the complement is in \([M]\), so it includes the
trailing coordinate. Multiply (T6) by the auxiliary mass, the two
permutation masses, the two-endpoint auxiliary ban, and the single
oriented scalar from (T5); sum over \(\nu\), auxiliaries and
permutations to obtain a kernel \(\mathsf B_\mu(d,\zeta,\theta,o,s)\).
Its pairing is
\(\sum_{u,v}\mu(u)\overline{F(u)}\mathsf B_\mu(u,v)H(v)\). This defines
the frozen form for arbitrary tests; no factorization of \(F\) or \(H\)
between groups is required.

Here is the identity connecting it to (M2). Temporarily remove the
auxiliary ban. In a comparison pattern replace the two character kernels
by ordinary probe resampling. Use the shared-copy indicators to identify
the unchanged coordinates. At a probe belonging to the first subset use
the target value, at a probe belonging to the second use the source
value, and integrate the freed counterpart. The trailing coordinates at
both endpoints are independently integrated. The divisor in the phase is
the head product of this spliced array, so the resulting pairing is
exactly \(-(-1)^{|a|+|b|}R(E_a\mathcal SF,E_b\mathcal SH)\). The empty
pattern has the opposite sign and gives
\(R(U\mathcal SF,U\mathcal SH)\). Changing variables by the independent
endpoint permutations proves

\[
\mathsf B_\mu(F,H)=\mathsf R(F,H)
              +\mathsf C(F,H)-\mathsf Z(F,H),
\tag{T7}
\]

where \(\mathsf R\) is the right side of (M2), \(\mathsf C\) is the sum
of (character-resampling minus ordinary-resampling) comparison forms
without the ban, and \(\mathsf Z\) is the character form on auxiliaries
meeting either endpoint. This is an identity of signed forms, before
taking absolute values. In particular, the phase is never commuted
through a coordinate average.

Equations (M4)--(M10) bound \(|\mathsf C|^2\) by
\(L^{-A}\|F\|_\mu^2\|H\|_\mu^2\) at any fixed accuracy \(A\). The same
bound holds for \(\mathsf Z\) eventually: for a fixed endpoint array,
the probability that one of at most \(mG\) auxiliaries hits its \(MG\)
marked values is at most \(mMG^2/(v e^{L^{3/10}})\). Split the ban into
its source and target parts and apply weighted Cauchy--Schwarz to each
part. The unrestricted marginal at the other endpoint is the original
product law. The pointwise kernel cap and pattern sum are polynomial in
\(L\), so this exponential atom saving absorbs them. This argument keeps
the tests in their original \(L^2\) norms; it does not bound a test on a
rare set by its supremum. Together with (M3) and
\(|a+b-c|^2\le3(|a|^2+|b|^2+|c|^2)\), (T7) gives the explicit squared
bound \(33L^{-A}\) times the two energies. To obtain exponent \(2E\),
use \(A=2E+2\). The reduction to \(d=s=1\) replaces \(\theta\) by
\(ds\theta\) and introduces the unit scalar \(d^{io\zeta}\); hence this
bound is uniform in \(d>0\), \(s\ne0\) and \(\theta\in\mathbb R\).

The density in (T6) must next be changed to (T2). This is done with one
multiplier common to every pattern and permutation. Put

\[
r_g(p)=b_p^{-1/2},\qquad
(\mathcal AF)(u)=\prod_{g,i}r_g(u_{g,i})F(u).
\tag{T8}
\]

Then \(\|\mathcal AF\|_\mu=\|F\|_{\widetilde\mu}\). The precise ratio of
the tilted joint law to the harmonic joint law, after this change at
both endpoints, is

\[
R_\nu(u,v;\sigma,\tau)=
 \prod_g\prod_{i\notin S_g}
 r_g(u_{g,\sigma_g(i)})r_g(v_{g,\tau_g(i)}).
\tag{T9}
\]

For a shared coordinate the indicator in (T6) makes the two endpoint
square roots supply its full ratio; for a free coordinate an additional
square root is needed at each end. This verifies (T9) coordinate by
coordinate. Auxiliary draws still use \(\mu\), not \(\widetilde\mu\).

Let \(\mathsf B_{\widetilde\mu}^{\rm nd}\) include the fresh-birth list
condition in \(E_o\), and let \(\mathsf B_\mu^{\rm bad}\) be the same
harmonic form restricted to its failure. The exact identity is

\[
\mathsf B_{\widetilde\mu}^{\rm nd}(F,H)
 =\mathsf B_\mu(\mathcal AF,\mathcal AH)
   +\mathsf T(F,H)-\mathsf B_\mu^{\rm bad}(\mathcal AF,\mathcal AH),
\tag{T10}
\]

where \(\mathsf T\) uses the multiplier \((R_\nu-1)\) and retains the
fresh-birth condition. Set \(P=e^{L^{3/10}}\), \(a_0=1-H_N/P\),
\(I=MG\), and

\[
\delta_{\rm tilt}=a_0^{-I}-1,\qquad
\mathcal K=1+(2^{m|\mathcal A|}-1)(2^{m|\mathcal B|}-1)Q^4\xi^{-2}.
\tag{T11}
\]

The raw pattern has cap one and every other pattern has cap
\(Q^4\xi^{-2}\), so \(\mathcal K\) is their total cap. From (T9),
\(1\le R_\nu\le a_0^{-I}\); normalized auxiliary and permutation sums
and the two marginals of (T6) therefore give

\[
\begin{aligned}
|\mathsf T(F,H)|^2
 &\le(\mathcal K\delta_{\rm tilt})^2
       \|F\|_{\widetilde\mu}^2\|H\|_{\widetilde\mu}^2,\\
|\mathsf B_\mu^{\rm bad}(\mathcal AF,\mathcal AH)|^2
 &\le\mathcal K^2\frac{I^2}{vP}
       \|F\|_{\widetilde\mu}^2\|H\|_{\widetilde\mu}^2.
\end{aligned}\tag{T12}
\]

For the second bound, expose the fresh coordinates in (T6). Any pair
coincides with conditional probability at most \(1/(vP)\); there are at
most \(I^2\) pairs. The restricted row has this saving and the full
column has mass one. Weighted Schur gives the displayed squared bound,
valid even for tests concentrated on the bad set. Since \(H_N\le5L\),
\(I=O_J(\log L)\), and \(\mathcal K\le L^{C(m,C_q,C_\xi)}\), both errors
in (T12) are smaller than any inverse power of \(L\). Applying the
three-term square inequality to (T10) gives \(9L^{-2E}\) times the
tilted endpoint energies, after choosing the harmonic precision. This is
the full tilted BIG estimate, including the ban and the actual list
filter.

\Needspace{6\baselineskip}

\hypertarget{marked-correlation-the-small-maps-and-the-translated-old-residues}{%
\subsubsection{The SMALL maps and the translated OLD
residues}\label{marked-correlation-the-small-maps-and-the-translated-old-residues}}

It remains to transfer this bound to (T4). Fix the OLD lengths \(j_g\)
and the ordered OLD primes \(\boldsymbol\pi_g\), retaining their
residues as variables. On this sector use the finite abelian group

\[
\mathbb T=\prod_{g\in\mathrm{SMALL}}\mathbb T_g
       \times\prod_{g\in\mathrm{BIG}}\prod_{i=1}^{j_g}
                                  \mathbb Z/\pi_{g,i}\mathbb Z.
\tag{T13}
\]

Repeated OLD primes give separate coordinates. The translation
\(t\mapsto t+\Delta\) acts diagonally on every occurrence; this is
exactly the clean OLD update. For a SMALL permutation \(\sigma\) and
shared tuple \(d=(d_{g,i})_{g,i<J}\), define

\[
(P_{\sigma,d}U)(a;y,r)=
 \prod_{g\in\mathrm{SMALL}}
    V_g^{-1/2}q^{\max(\omega_g(y)-M,0)/2}
 \sum_{\substack{(y,c)\in\prod_g\mathscr S_g\\
                   c_{g,\sigma_g(i)}=d_{g,i}\ (i<J)}}
          U((y,c),a,r).
\tag{T14}
\]

A fibre with no valid marks contributes zero. These are conditional
sums, not conditional probabilities. In the clean coefficient the SMALL
scalar is exactly the product of the two half factors in (T14). The
shared tuple has the same value at the two endpoints, and their residues
differ by \(\Delta=o sD\). Partitioning both endpoint sums by
\((d,y,r)\) gives precisely the pairing of \(P_{\sigma,d}U\) with the
translated \(P_{\tau,d}F\), with the frozen BIG kernel between them. The
displacement still depends on \(d\), the BIG shared marks and the
auxiliaries. It has not been frozen before this disintegration.

Use normalized Fourier coefficients, so that

\[
\widehat u(\chi)=\frac1{|\mathbb T|}\sum_tu(t)\overline{\chi(t)},
\qquad
\widehat{u(\cdot+\Delta)}(\chi)=\chi(\Delta)\widehat u(\chi).
\]

Parseval replaces translation by multiplication by a character. Write
\(\chi(\Delta)=\mathrm e(\vartheta_\chi\Delta)\); the frozen BIG block
now has frequency \(\theta+\vartheta_\chi\). The bound above is uniform
in this frequency, however large the OLD sector. For this fixed prime
sector put

\[
c_{\boldsymbol\pi}=
 \prod_{g\in\mathrm{SMALL}}V_g^{-M}
 \prod_{g\in\mathrm{BIG}}
       \frac{\prod_{i=1}^{j_g}\pi_{g,i}\widetilde\mu_g(\pi_{g,i})}{j_g!}.
\]

The factors \(\pi_{g,i}\) convert OLD counting measure to normalized
Haar measure on (T13). The exact clean pairing in this sector is

\[
c_{\boldsymbol\pi}\,
 \mathbb E_{\sigma,\tau}\sum_d\sum_{\chi\in\widehat{\mathbb T}}
 \mathsf B_{\widetilde\mu}^{\rm nd}
   \bigl(\widehat{P_{\sigma,d}U}(\chi),
         \widehat{P_{\tau,d}F}(\chi);
         d_S,\zeta,\theta+\vartheta_\chi,o,s\bigr),
\quad d_S=\prod_{g,i<J}d_{g,i}.
\tag{T15}
\]

The two arguments of \(\mathsf B\) are functions of the whole BIG active
array \(a\). After (T8) and (M0) they are the tests \(f,h\) of the
orthogonal-average argument. Formula (T15) is the required operator
identity: it keeps the full joint phase, both endpoint permutations, all
shared values, and every translated OLD occurrence.

There is no loss equal to \(|\mathbb T|\) in its norm estimate. For a
fixed \(y,d\), the number of completions in one SMALL group is
\((\omega_g(y)-J)_{1}\); for \(\ell\) omitted marks it is
\((\omega_g(y)-J)_{\ell}\). Cauchy--Schwarz in that fibre, followed by
summation over \(d\) and Parseval, bounds each lifted endpoint energy by
\(C_{\rm cond}\) times its original sector energy, where

\[
C_{\rm cond}=\prod_{g\in\mathrm{SMALL}}
       \frac{\ell!}{(1-q)^{\ell+1}V_g^\ell},\qquad \ell=1.
\tag{T16}
\]

Indeed, on a nonempty fibre \(\omega_g=J+\ell+z\), and
\(q^z(z+\ell)_{\ell}\le\sum_{z\ge0}q^z(z+\ell)_{\ell} =\ell!/(1-q)^{\ell+1}\).
Every original ordered mark state occurs once when its shared tuple is
summed. Cauchy--Schwarz over \((d,\chi)\) and then over sectors yields

\[
|\langle U,E_{\rm clean}^{(1)}F\rangle|^2
 \le9L^{-2E_0}C_{\rm cond}^2\,\mathcal E(U)\mathcal E(F).
\tag{T17}
\]

Here the superscript means that the OLD endpoint parameter is one. For
general \(\rho_0\le1\), multiply both tests by
\(\rho_0^{\#\text{zero OLD}}\); these multipliers and both OLD gates are
contractions. Pullback and orbit averaging in (T4) are contractions as
well. To pass from a pairing bound to an operator bound, set the first
test equal to the output of the operator and cancel its norm (the zero
output is immediate).

If \(C_G\log L\) bounds the number of SMALL groups and \(V_g\ge v\),
then \(C_{\rm cond}\le L^{c_{\rm cond}}\), with
\(c_{\rm cond}=C_G\log\max(1,\ell!/((1-q)^{\ell+1}v^\ell))\). Choose
\(E_0=E+c_{\rm cond}+1\) before choosing \(m\) or \(J\). For \(L\ge3\),
(T17) is at most \(L^{-2E}\) times the two energies. This proves the
clean part of (F7), with all normalizations specified.

\Needspace{6\baselineskip}

\hypertarget{marked-correlation-promotions-stores-and-the-full-transition}{%
\subsubsection{Promotions, stores and the full
transition}\label{marked-correlation-promotions-stores-and-the-full-transition}}

Fix an OLD BIG prime \(p\), its residue \(r\), the BIG divisor \(D_B\),
and an orientation \(o\in\{\pm1\}\). Let \(F\) be the product of all
SMALL marks and \(U\) the product of the omitted \(\ell\) marks in each
SMALL group. The retained SMALL divisor is \(D_S=F/U\). Promotion
requires

\[
r+osD_SD_B\equiv0\pmod p,
\quad\text{or equivalently}\quad
rU+osD_BF\equiv0\pmod p. \tag{M11}
\]

The constant term is nonzero: \(p\nmid D_BF\) and \(0<|s|<p\).
Furthermore

\[
\log U\le\ell\,\#\mathrm{SMALL}\,L^{1/5}
 =O(\log L\,L^{1/5})<L^{3/10}\le\log p.
\]

Thus \(0<U<p\), and (M11) determines at most one \(U\). Distinct marks
within each group and disjointness of the groups determine the omission
tuple from \(U\). Uniform coordinate permutations induce uniform
omission subsets. Hence

\[
\Pr(\text{an OLD promotion})
\le\frac{\#\mathrm{OLD}}
 {\binom{J+\ell}{\ell}^{\#\mathrm{SMALL}}}. \tag{M12}
\]

In the application \(\ell=1\), \(\#\mathrm{OLD}\le\lceil L^2\rceil\),
and \(\#\mathrm{SMALL}\ge c_0\log L\) with \(c_0=(40\log2)^{-1}\). The
right side is at most \(2L^{2-c_0\log(J+1)}\). If \(L^\kappa\) bounds
the complete row cost, choose \(J\ge m\) so that

\[
2H_0+2\kappa+3\le c_0\log(J+1). \tag{M13}
\]

This makes the promotion row rare enough for the squared-energy saving
\(L^{-2H_0}\). Reverse the transition orientation to obtain the same
estimate for storing an active label: it is a rare column. The weighted
Schur inequality combines each rare row or column bound with the full
complementary row bound, then adds the two exceptional operators.
Increasing the fixed precision pays their constant sum. This is why
\(J\) is chosen after the mesh, conductor, and pattern parameters.

The clean and exceptional parts sum to the literal operator \(E_o\).
Apply their squared estimates at precision \(E+1\). The inequality
\(\|u+v\|^2\le2\|u\|^2+2\|v\|^2\) gives

\[
\mathcal E(E_oF)\le4L^{-2(E+1)}\mathcal E(F)
                  \le L^{-2E}\mathcal E(F). \tag{T18}
\]

The exceptional estimate did not require cancellation. The same positive
row and column bounds, without the rare-event factor, apply if an
arbitrary coefficient of modulus at most one is attached to each
expanded transition. Thus, for a fixed exponent \(\kappa\),

\[
\mathcal E(E_o^\phi F)\le L^{2\kappa}\mathcal E(F),
\qquad |\phi|\le1. \tag{T19}
\]

The exponent is independent of \(J\); its formula is given below. This
distinction is needed in the forest proof: a sign multiplier at a
collision owner uses (T19), not (T18).

\Needspace{6\baselineskip}

\hypertarget{marked-correlation-background-transitions-and-the-constants-in-the-moment-bound}{%
\subsubsection{Background transitions and the constants in the moment
bound}\label{marked-correlation-background-transitions-and-the-constants-in-the-moment-bound}}

At visit \(i\), the background operator \(G_i\) leaves the SMALL state
and all active BIG marks fixed. In one BIG group let \(Z\) be the
indices of OLD atoms at residue zero. Choose \(A\subset Z\), delete
those OLD occurrences, and append \(b\) new atoms
\((p_1,0),\ldots,(p_b,0)\). The coefficient of this choice is exactly

\[
 \left(-\frac{\eta_i}{\rho_0}\right)^{|A|}
 \left(\frac{q_i}{\rho_0^2}\right)^{|Z|-|A|}
 \frac1{b!}\left(-\frac{\eta_i V_g}{\rho_0}\right)^b
 \prod_{j=1}^b\widetilde\mu_g(p_j).
\tag{T20}
\]

Sum over these choices, with the finite output cap, then take the
product over groups and the orbit average (T4). These appended
occurrences are \emph{background births}. A later use of an existing
occurrence is not a birth. Formula (T20) explains both the signs and the
single birth factorial. The OLD factorial in (T3) is the Hilbert measure
and is not inserted a second time in an edge coefficient. The powers of
\(\rho_0\) in (T5) and (T20) cancel at successive endpoints; they are an
auxiliary norm device, not a change to the original residue density.

Here are explicit constants for (F7)--(F11). They also distinguish the
row constant used in removing the OLD cutoff from the smaller
conditional-map constant (T16). For \(d_0=\ell+m\) put

\[
\begin{aligned}
C_{\rm prom}(\rho_0,d_0)&=\sum_{j=0}^{d_0}\frac{j!}{(1-\rho_0)^{j+1}},\\
C_{\rm small}&=\prod_{g\in\mathrm{SMALL}}
       \frac{\ell!}{V_g^\ell(1-\sqrt q)^{\ell+1}},\\
S_{\rm small}&=\sum_{x\in\prod_{g\in\mathrm{SMALL}}\mathscr S_g}
                    \prod_g\frac1{|\mathbb T_g|V_g^M},\\
a&=\mathcal K C_{\rm small}
 \prod_{g\in\mathrm{BIG}}
 \left[w_0^{d_0}4^{d_0}C_{\rm prom}(\rho_0,d_0)
                 \max(1,V_g^{-1})^{d_0}2^{d_0}\right].
\end{aligned}\tag{T21}
\]

Here \(\mathcal K\) is (T11), and \(w_0\) is specified in
\protect\hyperlink{marked-correlation-restoring-all-histories}{the
history restoration}. Thus \(S_{\rm small}\) is a finite total mass, not
a universal constant or a probability normalization. Both
\(C_{\rm small}\) and \(a\) depend on the actual group masses and \(L\);
their polynomial \textbf{exponents} are fixed. The bounds on \(V_g\) and
the number of groups make \(\mathcal K C_{\rm small}\le L^{C_E'}\) and
\(a\le L^{C_b}\) eventually. \(C_E'\) here is a row-budget exponent,
unrelated to the ordinary-slot exponent \(C_E\) in (B1).

An explicit exponent for (T19) can be read without choosing \(J\).
Choose \(C_G\) so that both group counts and
\(|\mathcal A|+|\mathcal B|\) are at most \(C_G\log L\) and put

\[
\begin{aligned}
R_v&=\max(1,v_{\max},v^{-1}),\\
B_{\rm row}&=\max\{1,4^{d_0}C_{\rm prom}(\rho_0,d_0)
                                  \max(1,v^{-1})^{d_0}2^{d_0}\},\\
\kappa&=mC_G\log2+4C_q+2C_\xi+1\\
&\quad+C_G\log\max\left(1,
             \frac{\ell!}{(1-\sqrt q)^{\ell+1}v^\ell}\right)
       +C_G\log B_{\rm row}+\frac{d_0C_G}2\log R_v.
\end{aligned}\tag{T22}
\]

To see the uniformity, at most \(d_0\) slots per group are unshared. If
\(z\) OLD occurrences have residue zero, choosing \(j\le d_0\) of them
costs at most \((z)_j\rho_0^{z-j}\le C_{\rm prom}\). The SMALL row is
bounded by the same falling-factorial calculation as (T16), with
\(\sqrt q\) in place of \(q\). The number of patterns and the character
caps give the first line of the exponent in (T22). The comparison of
forward and reverse harmonic masses contributes at most \(R_v^{d_0G}\)
to the squared bound. Weighted Schur, followed by (T4), gives (T19).
Applying the rare-row estimate (M12) to this same calculation gives
(M13), so \(J\) can indeed be chosen afterwards.

The background estimate is uniform even in unit multipliers on its
individual choices. For \(W_0,w_0,\rho_0\) as below, summing (T20) in
absolute value with weight \(w_0^{|\mathrm{OLD}|}\) uses

\[
\frac1{\rho_0w_0}+\frac{\sqrt q}{\rho_0^2}\le1,
\qquad \sum_{p\in P_g}\widetilde\mu_g(p)\le2.
\]

The deletion sum is binomial and the birth sum is exponential. In the
transposed sum the factorials in (T3) give the corresponding binomial
coefficients; hence the same weighted Schur argument gives
\(\mathcal E(G_i^\phi F)\le\Gamma^2\mathcal E(F)\), with
\(\Gamma=\exp(4W_0w_0G/\rho_0)\). Set

\[
l=L^{-S_0},\qquad h=L^\kappa,\qquad
b=L^{S_0+\kappa},\quad K_r=S_0+\kappa,\qquad
P_0=\left|\prod_{g\in\mathrm{BIG}}\prod_{p\in P_g}b_p\right|
                \mathcal E(v)\Gamma^{N+1},
\tag{T23}
\]

where \(v\) is one on states with empty OLD memory and zero otherwise.
The initial distinct-active mask is retained in the history sum and can
only reduce its boundary energy. Equations (T18)--(T23) supply all three
estimates in (F7); the forest argument below restores the omitted
distinctness between different birth times.

\Needspace{6\baselineskip}

\hypertarget{marked-correlation-restoring-all-histories}{%
\subsection{Restoring all
histories}\label{marked-correlation-restoring-all-histories}}

Equalities between prime labels in a history are indexed by forests. The
exact signed expansion, its coefficients, and the two bounds below are
proved in
\protect\hyperlink{forest-moments}{Section~\ref{forest-moments}}. Its
random-sign factorization explains why each identification loses the
local contraction at at most two visits. For
\(N=2\lceil L^{13/25}\rceil\), the number of available labels is bounded
by

\[
\mathcal A=(N+1)\#\mathrm{BIG}(J+1+\lceil L^2\rceil)\le L^4.
\]

A forest of rank \(r\) requires \(r\) identifications and there are at
most \(\mathcal A^{2r}\) such forests. The expansion has two bounds for
every forest contribution \(F_r\):

\[
|F_r|\le
\min\{P_0(L^{-S_0})^N b^{2r},\ P_0a^N\delta^r\}, \tag{M15}
\]

where

\[
\begin{gathered}
P_0\le L^{(C_v+2)N},\qquad b\le L^{K_r},\qquad a\le L^{C_b},\\
S_0=E+C_v+5,\qquad \delta=\frac2{v e^{L^{3/10}}}.
\end{gathered}
\]

The first bound uses the local signed contraction at each step; each
identification replaces at most two local factors by their row bounds.
The second uses absolute row bounds everywhere and charges one factor
\(\delta\) to each independent equality of BIG prime labels. This is a
finite reindexing of the same history expansion.

Split at \(r=\lceil NL^{-3/20}\rceil\). Below the split, the exponent
cost \((8+2K_r)r\) from the forest count and \(b^{2r}\) is at most \(N\)
eventually. The first bound in (M15) is therefore at most
\(L^{-(E+2)N}\) after counting forests of that rank. Above the split,

\[
(\mathcal A^2\delta)^r
 \le e^{-rL^{3/10}/2}
 \le e^{-NL^{3/20}/2},
\]

which dominates \(L^{(C_v+C_b+E+4)N}\). The second bound gives the same
\(L^{-(E+2)N}\). There are at most \(\mathcal A+1\le L^N\) ranks;
summing proves an upper bound \(L^{-EN}\) for the signed history sum
with the OLD cutoff \(\#\mathrm{OLD}\le\lceil L^2\rceil\).

There is no circular parameter choice here. The fixed volume exponent
\(C_v\) is chosen from the harmonic mass bounds and \(q\), before \(m\)
or \(J\). One may first choose \(\rho_0\in(q^{1/4},1)\) and then use

\[
W_0=\max(1,v_{\max}),\quad
w_0=\max\!\left(1,\frac{4W_0}{\rho_0(1-\sqrt q/\rho_0^2)}\right),
\quad C_v=\frac{4W_0w_0C_0}{\rho_0},
\]

where \(C_0\log L\) bounds the group count. Thus \(S_0\) is fixed before
the local parameters. The later exponents \(K_r,C_b\) affect only the
eventual threshold in the rank split.

To remove this cutoff, put \(T_0=\lceil L^2\rceil\) and
\(H_b=\lceil L^2/4\rceil\). The number of active BIG labels encountered
at all visits is at most \((N+1)G(J+1)\), where \(G\) is the BIG group
count. Eventually \(H_b+(N+1)G(J+1)\le T_0\). The first crossing of the
OLD cutoff therefore requires at least \(H_b\) background births. In the
positive history expansion, multiply each birth weight by \(2\). The
discarded mass is at most \(2^{-H_b}\) times this single positive root
sum; no factor counting histories is needed. Its explicit bound is

\[
S_{\rm small}e^{4(N+1)(1+\sqrt L)}
 \left[C_\nu C_{\rm small}
  \left(16\max(1,v^{-1})e^{\sqrt L}\right)^{(1+m)G}\right]^N,
\]

Here \(C_\nu=\mathcal K\), and

\[
C_\nu C_{\rm small}\le L^{C_E'},\qquad
S_{\rm small}\le
 \exp\!\bigl((C_0+1)L\log L(\sqrt L+\log\max(1,v^{-1}))\bigr).
\] The discarded mass is consequently at most
\(\exp(-\tfrac14(\log2)L^2+O_{\rm fixed}(L^{7/4}))\), hence at most
\(L^{-(E+1)N}\). Use the forest bound at precision \(E+1\) as well.
Adding the two errors gives \(2L^{-(E+1)N}\le L^{-EN}\) for the full
history sum.

\Needspace{6\baselineskip}

\hypertarget{marked-correlation-from-residue-averages-to-the-integer-interval}{%
\subsection{From residue averages to the integer
interval}\label{marked-correlation-from-residue-averages-to-the-integer-interval}}

For a fixed history \(\omega\), let \(h_i\) be its offsets. Its density
at an integer \(n\) is

\[
D_\omega(n)=\prod_p\prod_{i=0}^N
\begin{cases}
\mathbf1_{p\mid n+h_i},&p\text{ is active at visit }i,\\
q_i^{\mathbf1_{p\mid n+h_i}},&\text{otherwise},
\end{cases} \tag{M16}
\]

with \(q_0=q_N=\sqrt q\) and \(q_i=q\) internally. Its model density
\(v_\omega\) is the product of the exact averages over residues modulo
each prime. This retains coincidences between offsets modulo a prime.
Explicitly, let \(A_p\) be the visits at which \(p\) is active, and
\(I_{p,t}=\{i:h_i\equiv t\pmod p\}\). The local factor is

\[
v_{\omega,p}=\begin{cases}
0,&\text{the offsets on }A_p\text{ are not all congruent},\\
\displaystyle\frac1p\prod_{i\in I_{p,t}\setminus A_p}q_i,
 &A_p\ne\varnothing,\ h_i\equiv t\ (i\in A_p),\\
\displaystyle1-\frac1p\sum_{t\bmod p}
              \left(1-\prod_{i\in I_{p,t}}q_i\right),&A_p=\varnothing.
\end{cases}\tag{M16a}
\]

The minus sign between the initial integer residue and \(t\) disappears
under uniform averaging. There is one residue per \textbf{numerical
prime} in (M16a), even if it appears in many active lists.

We verify its connection to the memory expansion, which is needed before
applying an interval estimate. A prime lifespan is either empty or an
interval \([a,z]\subset[0,N]\). For a nonempty lifespan require
\(h_a\equiv h_z\pmod p\), require every active visit to lie in \([a,z]\)
and have offset congruent to \(h_a\), and allow \(a=z\) only when \(a\)
is active. Its weight is

\[
\frac1p\epsilon_a\epsilon_z
 \prod_{a<i<z,\ i\notin A_p}
       (1-\eta_i1_{h_i\equiv h_a\ (p)}),
\qquad
\epsilon_i=\begin{cases}1,&i\in A_p,\\-\eta_i,&i\notin A_p.\end{cases}
\tag{T24}
\]

The empty lifespan is allowed only if \(A_p=\varnothing\) and has weight
\(b_p=1-\sum_i\eta_i/p\). The sum of all these weights is exactly
(M16a). For an absent prime, expand the elementary identity

\[
\prod_{i\in I}(1-\eta_i)
 =1-\sum_{i\in I}\eta_i
  +\sum_{a<z,\ a,z\in I}\eta_a\eta_z
                      \prod_{a<i<z,\ i\in I}(1-\eta_i)
\tag{T25}
\]

in each residue class \(I=I_{p,t}\) and average over \(t\). The identity
follows by telescoping twice, or induction on \(|I|\). For a represented
prime, let \(a_0,z_0\) be its first and last active visits. The inactive
factors strictly before \(a_0\) have the expansion
\(1-\sum_{a<a_0}\eta_a1_{h_a\equiv h_{a_0}}  \prod_{a<i<a_0}(1-\eta_i1_{h_i\equiv h_{a_0}})\);
use the reversed expansion after \(z_0\) and leave the middle factors
unchanged. Multiplication gives exactly (T24), including the cases
\(a=a_0\), \(z=z_0\), and a single active visit.

To recover the transition description from these lifespans, follow each
prime through its active visits and keep it in OLD between separated
active runs. Leaving an active run stores it; returning promotes the
same occurrence. An inactive left endpoint of a lifespan is a background
birth, and an inactive right endpoint a background deletion. The
intermediate residue is \(h_i-h_a\pmod p\), so all zero tests are the
ones in (T24). Conversely, a history with distinct births determines
these lifespans and active runs uniquely, apart from the ordering of
simultaneous births and OLD representatives. Their permutation counts
cancel respectively the birth factorials in (T20) and the sector
factorials in (T3). At each prime the \(\rho_0\) powers cancel;
\(b_p\widetilde\mu_g(p)=1/(V_gp)\); the remaining \(V_g^{-1}\) factors
count active runs. These are exactly the factors obtained from the
integer matrix's state and edge masses. Thus multiplication over primes
and summation over the frame choices identify the restored signed memory
sum with \(\sum_\omega c_\omega v_\omega\). The frame coefficient
retains its sign, auxiliary ban and complex phase throughout this
equality.

Here is the finite approximation used for (M16). If active congruences
at some prime are incompatible, both the integer density and its residue
average vanish. Otherwise let \(Q_0\) be the product of the distinct
active primes. Their congruences specify one class modulo \(Q_0\), with
a fixed positive factor from the inactive visits at those primes. For
each remaining prime put

\[
f_p(r)=1-\prod_{i=0}^N q_i^{1_{r+h_i\equiv0\ (p)}}.
\]

Then \(0\le f_p\le1\), it is supported on at most \(N+1\) residue
classes, and its average is at most \((N+1)/p\le1/2\) eventually, since
every marked prime is at least \(e^{L^{1/10}}\). Put

\[
\lambda=\sum_p\frac1p\sum_{r\bmod p}f_p(r),\qquad
V_0=\prod_p\left(1-\frac1p\sum_r f_p(r)\right),\qquad
U_0=1+\sum_p\sum_r f_p(r).
\]

Bonferroni truncation at orders \(2K\) and \(2K+1\), followed by CRT on
each retained product, gives relative error on an interval of length
\(X\) at most

\[
\frac{e^{4\lambda}}{2^{2K+1}}
 +\frac{Q_0}{X}\,2(2K+2)e\,U_0^{2K+1}e^{2\lambda}. \tag{M17}
\]

Indeed, the omitted elementary symmetric sums are bounded using the sum
\(\lambda\) of the local averages; the normalization by \(V_0\) costs at
most \(e^{2\lambda}\). For each retained residue pattern, CRT counts its
interval representatives with error at most one. Summing the absolute
coefficients of the retained patterns gives the second term in (M17).

Take \(K=\lceil L^{11/20}\rceil\). The actual ranges give

\[
\lambda=O(N\log L),\qquad
\log Q_0=O_J(L^{23/25}\log L),\qquad
K\log U_0=O(L^{19/20}),
\]

whereas \(X\ge e^{\gamma L}\) for a fixed \(\gamma>0\). Both terms in
(M17) are smaller than \(L^{-EN}\) for every fixed \(E\) eventually. For
complex history coefficients, the transfer is made before discarding any
signs:

\[
\left|\sum_\omega c_\omega\sum_nD_\omega(n)
 -X\sum_\omega c_\omega v_\omega\right|
\le XL^{-EN}\sum_\omega|c_\omega|v_\omega. \tag{M18}
\]

The absolute mass on the right is at most \(L^{CN}\). Its exponent is
fixed after \(J\); increasing the requested CRT accuracy absorbs it. The
signed main sum on the left is the one already bounded by (M15).

We spell out the finite physical matrix to which these history estimates
apply. Fix a seed interval \(I=[a,b)\cap\mathbb Z\) of positive
integers, of length \(X=b-a\ge e^{\gamma L}\). Put
\(R=\lceil L^{13/25}\rceil\),
\(r_{\rm path}=\lceil2R\Delta_{\max}\rceil\), with \(\Delta_{\max}\) as
in (M20a), and take the carrier positions to be

\[
I^+=[a-r_{\rm path},b+r_{\rm path}]\cap\mathbb Z.
\]

Every prefix of every \(2R\)-step history from \(I\) belongs to \(I^+\):
each displacement has absolute value at most \(\Delta_{\max}\), so the
triangle inequality bounds a prefix by \(2R\Delta_{\max}\). Thus
truncating the matrix to this common carrier introduces no
history-dependent restriction on its starting interval. All residue
counts below use the same \(I\) and the same length \(X\).

Put \(M=J+1\). A state is an integer position \(n\in I^+\), together
with, for each marked group \(g\), an ordered list \(\boldsymbol p_g\)
of \(M\) distinct primes dividing \(n\). The residues of \(n\) determine
all residue coordinates of the state. All states have the same reference
mass

\[
m_0=\prod_g V_g^{-M};\qquad
u_n=\sqrt{m_0}\sum_{\substack{s\text{ a state}\\\operatorname{pos}(s)=n}}e_s.
\tag{M18a}
\]

Thus no mark fibre is normalized to have mass one.

A frame consists of a raw pattern or a nonempty pair \((a,b)\) from
(M2), an auxiliary prime for each freed BIG probe coordinate, and
independent source and target permutations \(\sigma_g,\tau_g\in S_M\) in
every group. For a comparison pattern, let
\(C_g\subseteq\{0,\ldots,m-1\}\) be its freed coordinates in group
\(g\); for a raw pattern and for SMALL groups, \(C_g=\varnothing\). The
shared coordinates are \(S_g=\{0,\ldots,J-1\}\setminus C_g\). Write the
auxiliary primes as \(\pi_{g,i}\) and give them the law

\[
\mu_{\rm aux}=\prod_{g\in\mathrm{BIG}}\prod_{i\in C_g}
                 \frac1{V_g\pi_{g,i}}.
\]

Auxiliary values may repeat. Each must differ from every source and
target mark in its own group. For states \(s,t\), a frame is allowed
when their permuted marks agree at each shared coordinate and

\[
\begin{aligned}
D_f&=\prod_g\prod_{i\in S_g}p^s_{g,\sigma_g(i)}
       \prod_{g\in\mathrm{BIG}}\prod_{i\in C_g}\pi_{g,i},\\
\operatorname{pos}(t)&=\operatorname{pos}(s)+2kD_f.
\end{aligned} \tag{M18b}
\]

All factors of \(D_f\) retain their multiplicities. For the raw pattern
let \(K_f=1\). For a comparison pattern put

\[
K_f=-(-1)^{|a|+|b|}K_A(D_a,X_a)K_B(D_b,X_b),
\]

where \(D_a,D_b\) are the corresponding products of auxiliary primes,
\(X_a\) is the product of the \textbf{target} marks in the first probe
block, and \(X_b\) is the product of the \textbf{source} marks in the
second block. Both endpoint products use the permuted lists. The two
kernels are exactly (M4), with the actual cell masses for these probe
products.

For an allowed frame define

\[
w_f(s,t)=\mu_{\rm aux}\prod_g
 \left[(M!)^{-2}V_g^{-(M-|S_g|)}
 q^{(\omega_g(\operatorname{pos}(s))+
       \omega_g(\operatorname{pos}(t))-2M)/2}\right].
\]

The powers of \(q\) are the two endpoint dampings. Their exponents are
nonnegative because each endpoint has \(M\) distinct marked divisors.
The physical matrix is

\[
A_{s,t}=\sum_{f\text{ allowed}}w_f(s,t)K_fD_f^{i\zeta}. \tag{M18c}
\]

Reversing the two endpoint permutations preserves \(D_f\) and the real
weight. Reversing the orientation conjugates the whole complex
coefficient, including \(K_f\). The backwards matrix is therefore
exactly \(A^*\).

Expanding \((AA^*)^R\) gives the alternating histories above. Insert
\(\mathrm e(2kD_f\theta)\) on a forward edge and its conjugate on a
backward edge. The product is \(\mathrm e(\theta h_N)\), so integration
for \(0\le\theta\le1\) keeps precisely the histories returning to the
initial integer position. The initial and final mark lists remain free;
the vectors \(u_n\) in (M18a) perform those sums. The restored history
bound and (M18) consequently give

\[
T:=\sum_{n\in I}\langle u_n,(AA^*)^Ru_n\rangle
\le XL^{-2RE},\qquad R=\lceil L^{13/25}\rceil. \tag{M19}
\]

\Needspace{6\baselineskip}

\hypertarget{marked-correlation-the-endpoint-bound}{%
\subsection{The endpoint
bound}\label{marked-correlation-the-endpoint-bound}}

The first endpoint has the form \(f=\sum_{n\in I}c_nu_n\); its
coefficient depends on the integer position but not on the individual
marks. The second endpoint \(g\) may depend on all marks. The finite
spectral estimate used here is

\[
|\langle f,Ag\rangle|
\le\sqrt3\,\|g\|\,\|f\|^{1-1/R}
 (H_{\rm loc}C_{\rm sup}^2)^{1/(2R)}T^{1/(2R)}. \tag{M20}
\]

Here \(|c_n|\le C_{\rm sup}\) and a nonzero matrix entry changes the
integer position by at most the positive integer \(H_{\rm loc}\). To
prove the estimate, partition the positions into disjoint intervals
\(I_j\) of \(H_{\rm loc}\) integers. Let \(f_j\) be the restriction of
\(f\) to \(I_j\), and let \(g_j\) be the restriction of \(g\) to
\(I_{j-1}\cup I_j\cup I_{j+1}\). Locality gives
\(\langle f,Ag\rangle=\sum_j\langle f_j,Ag_j\rangle\) and
\(\sum_j\|g_j\|^2\le3\|g\|^2\). Put \(P=AA^*\). Spectral interpolation
for the nonnegative matrix \(P\), followed by Cauchy--Schwarz, gives

\[
|\langle f_j,Ag_j\rangle|
 \le\|g_j\|\,\|f_j\|^{1-1/R}
            \langle f_j,P^Rf_j\rangle^{1/(2R)}.
\]

Indeed, the spectral measure of \(P\) at \(f_j\) satisfies
\(\langle f_j,Pf_j\rangle\le \|f_j\|^{2-2/R}\langle f_j,P^Rf_j\rangle^{1/R}\).
Writing \(f_j=\sum_{n\in I\cap I_j}c_nu_n\) and applying Cauchy--Schwarz
to the vectors \(P^{R/2}u_n\) yields

\[
\langle f_j,P^Rf_j\rangle
 \le H_{\rm loc}C_{\rm sup}^2
       \sum_{n\in I\cap I_j}\langle u_n,P^Ru_n\rangle.
\]

Hölder's inequality in \(j\), with exponents \(2\), \(2R/(R-1)\) and
\(2R\), proves (M20), since \(\sum_j\|f_j\|^2=\|f\|^2\). For \(R=1\),
omit the middle factor and use Cauchy--Schwarz. In the application
\(R>1\).

In the application we may take \(C_{\rm norm}=9\) and
\(C_{\rm sup}=e^{5\sqrt L}\), with \(\|f\|,\|g\|\le\sqrt X L^9\)
eventually. To specify the window, put

\[
\begin{aligned}
C_D&=4(J+1)(4+1/\log2),\\
\Delta_{\max}&=2L^{C_k}\exp(C_D\log L\,L^{2/5}),\\
H_{\rm loc}&=2\lceil2R\Delta_{\max}\rceil+1.
\end{aligned}\tag{M20a}
\]

Every frame satisfies \(D_f\le\exp(C_D\log L\,L^{2/5})\), including
repeated auxiliaries, and every \(2R\)-step path from \(n\) stays within
\(\lceil2R\Delta_{\max}\rceil\) of \(n\). The reference measure in (M20)
is ordinary counting measure on the finite physical states;
\(\sqrt{m_0}\) is in the vectors, not the norm. Since
\(\log H_{\rm loc}=O_J(\log L\,L^{2/5})\) and \(R\ge L^{13/25}\), the
local factor in (M20) is bounded. Apply (M19) with
\(E'=E+2C_{\rm norm}+1\). Then (M20) is at most \(XL^{-E}\) eventually.
The \protect\hyperlink{physical-lift}{physical-lift argument} gives the
exact endpoint vectors, the Fourier pairing identity (L8), and the
shared-prime restoration error (L4). It obtains the \(J+1\) marked
physical states by multiplying both original integers by a common
product of \(J\) primes per group; it does not require the original
integers in (B17) to have that many distinct marked divisors.

This proves the saving for the original term minus its arithmetic
comparison term. It remains to estimate that comparison term itself.

\Needspace{6\baselineskip}

\hypertarget{marked-correlation-the-comparison-term-and-the-fine-discrepancy}{%
\subsection{The comparison term and the fine
discrepancy}\label{marked-correlation-the-comparison-term-and-the-fine-discrepancy}}

We give the endpoint used in the comparison explicitly. If \(G\) is the
union of the marked prime groups, let \(n_*\) be the part of \(n\) with
all prime factors in \(G\) removed. For bounded coefficients
\(\alpha,c\), a prime slot \(Q_*\) disjoint from \(G\), a character
\(\chi\), and real frequencies \(\sigma_0,\sigma_1\), set

\[
\operatorname{Core}(n_*)=
 \sum_{\substack{m p\mid n_*\\p\in Q_*}}
 \alpha(m)\chi(m)m^{i\sigma_0}p^{i\sigma_1}c(n_*/mp).
 \tag{M21}
\]

For each marked group \(P_g\), let \(\mathcal L_g(n)\) consist of its
ordered lists of \(t_g\) distinct prime divisors of \(n\), and let
\(\omega_g(n)\) count those distinct divisors. For a function \(b\) of
the BIG lists with \(|b|\le1\), the literal marked weight is

\[
W_{\boldsymbol t}^{,b}(n)=
 \prod_g\frac{q^{\omega_g(n)-t_g}}{V_g^{t_g}}
 \sum_{\substack{\boldsymbol p_g\in\mathcal L_g(n)\\\text{all }g}}
 b((\boldsymbol p_g)_{g\in\mathrm{BIG}}). \tag{M22}
\]

This includes the distinguished SMALL group; in the source its number of
marks is \(k+1\). Multiplicities in \(n\) are not restricted by the
distinctness of the marks. With a smooth cutoff \(\Lambda(n)\), the
endpoint Fourier polynomial is

\[
\widehat E(\theta)=\sum_{1\le n\le N_{\rm cap}}
 \Lambda(n)\varepsilon(n)\operatorname{Core}(n_*)
 W_{\boldsymbol t}^{,b}(n)\mathrm e(n\theta). \tag{M23}
\]

The cell/character kernel (M4) has a concrete rank decomposition: each
\((j,\chi)\) contributes the product of
\(1_{D\in I_j}\overline{\chi(D)}\) and \(1_{X\in I_j}\chi(X)\), with
coefficient \(\nu_j^{-1}\). There are at most \(n_AQ^2\) such terms in
the first kernel and \(n_BQ^2\) in the second. Thus a comparison pattern
has at most \(n_An_BQ^4\) pairs, each coefficient at most \(\xi^{-2}\).
The support counts (M9a) bound this rank by a fixed power of \(L\). For
retained cells there is a sharper bound: their disjoint probabilities
sum to at most one, so their number is at most \(\xi^{-1}\) in either
block. Including all patterns gives

\[
\begin{aligned}
\#\{\text{ranks}\}&\le L^{C_{\rm pat}}Q^4\xi^{-2},\\
\sum_{\text{ranks}}\nu_{j_A}^{-1}\nu_{j_B}^{-1}
 &\le L^{C_{\rm pat}}Q^4\xi^{-4},\qquad
C_{\rm pat}=m(3+1/\log2)\log2.
\end{aligned}\tag{M23a}
\]

Here \(3+1/\log2\) is an upper group-count constant for the actual BIG
family. The two inverse cell masses stay outside the endpoint tests. The
remaining endpoint cell indicators and characters are the bounded list
coefficient \(b\) in (M22).

The auxiliary sum retains its dependence on the shift and cutoff. More
explicitly, let \(\nu_A(u),\nu_B(v)\) be the probability laws of the two
auxiliary products. Independence of the two disjoint probe blocks gives
the exact finite identity

\[
\begin{aligned}
&\sum_{\rm aux}\mu_{\rm aux}
 1_{D_A\in I_{j_A},D_B\in I_{j_B}}
 \overline{\chi_A(D_A)\chi_B(D_B)}\,H(D_AD_B)\\
&\quad=\sum_{u\in I_{j_A}}\sum_{v\in I_{j_B}}
 \nu_A(u)\nu_B(v)\overline{\chi_A(u)\chi_B(v)}\,H(uv).
\end{aligned}\tag{M23b}
\]

The inverse cell masses and the pattern sign multiply both sides. After
the shared-divisor restoration described below, \(H(c)\) retains the
shift \(2kc\) and the cutoff \(\Psi(n/(Zc))\). Thus (M23b) does not
replace this dependence by a mass-one average. Fourier inversion of the
cutoff as in (L8) separates its Mellin factor; additive Fourier
inversion separates the two endpoint integers. This yields the pair of
polynomials (M23), with an auxiliary bilinear phase
\(\mathrm e(-2kuv\theta)\) on the right of (M23b). Partition \(u,v\)
into the coarse dyads of (M9a) when applying (M7). The scales of a
retained cell pair are \(H'=e^{j_Ah_A+j_Bh_B}\) and \(Y=ZH'\); its
actual product lies between \(H'\) and \(e^{h_A+h_B}H'\). A fixed smooth
cutoff equal to one on this ratio interval gives the endpoint support
used in (M25). All rank, coefficient and Fourier costs are fixed before
requesting the discrepancy precision below.

Here is the resulting comparison identity with the endpoint factors
displayed. For a pattern \((a,b)\) let \(t_g=M-|S_g|\) be its number of
free marks; thus \(t_g=1\) in SMALL and \(t_g=1+|C_g|\) in BIG. For rank
\(r=(j_A,\chi_A,j_B,\chi_B)\) define the two list functions

\[
b_{1,r}=\overline{1_{X_B\in I_{j_B}}\chi_B(X_B)},\qquad
b_{2,r}=1_{X_A\in I_{j_A}}\chi_A(X_A).
\tag{M23c}
\]

Here \(X_B\) is the source probe product and \(X_A\) the target probe
product, with the shared entries deleted and the free entries kept in
their original coordinate order. Both functions have modulus at most one
and contain no auxiliary variables. With the two cores of (L7), define

\[
U_r(n)=\varepsilon(n)n^{it_1}\operatorname{Core}_1(n_*)
                  W_{\boldsymbol t}^{,b_{1,r}}(n),\qquad
V_r(n)=\varepsilon(n)n^{it_2}\operatorname{Core}_2(n_*)
                  W_{\boldsymbol t}^{,b_{2,r}}(n).
\tag{M23d}
\]

The first core is conjugated when \(U_r\) is conjugated in the pairing.
For (B17) the phases are \(t_1=-(\tau+\xi_3)\), \(t_2=-\tau\),
\(\zeta=-\xi_3\), as in (L7)--(L8). Write \(c=uv\) and
\(s_{a,b}=-(-1)^{|a|+|b|}\). The independent-rank main contribution of
this pattern and rank to the signed comparison part of \(A\) (hence to
\(-Q_{\rm cmp}\) in (L8)) is exactly

\[
\begin{aligned}
\frac{s_{a,b}}{\nu_{j_A}\nu_{j_B}}
\sum_{\substack{u\in I_{j_A}\\v\in I_{j_B}}}
 &\nu_A(u)\nu_B(v)\overline{\chi_A(u)\chi_B(v)}\,c^{i\zeta}\\
&\cdot\sum_{\substack{n>0\\n+2kc>0}}
 \frac{\Psi(n/(Zc))}{n}\overline{U_r(n)}V_r(n+2kc).
\end{aligned}\tag{M23e}
\]

To check the removal of the shared divisor \(d\), substitute the
original positions \(dn,dn'\) and \(D_f=dc\). The displacement becomes
\(n'=n+2kc\) and the cutoff becomes \(\Psi(n/(Zc))\). The ordinary parts
of the two integers are unchanged, since all prime factors of \(d\) are
marked. The common parity is squared. The phase in \(d\) is
\(d^{i(-t_1+t_2+\zeta)}=1\). The shared harmonic factors then sum to one
on the unrestricted shared family; the exclusions and repeated-mark
cases produce the weighted restoration error (M28), not an exact
mass-one assertion on the restricted family.

There is one further explicit difference in (M23e): the original
auxiliary ban forbids a free endpoint mark from equalling an auxiliary
prime. Partitioning the complete pair of free mark lists into the
allowed and forbidden sets gives, exactly, the independent product of
the two weights minus its collision contribution. A forbidden pair
forces \(p\mid n\) for an auxiliary prime \(p\ge e^{L^{3/10}}\); if it
divides the other endpoint, it also divides \(n\), since \(p\mid c\).
Reindex \(n=pr\). The factor \(1/n\) and the divisor-moment bound for
the two cores give a factor \(1/p\) with a fixed logarithmic cost. If
\(L^{E_{\rm mark}}\) bounds either full marked weight, the error for one
pattern and one fixed auxiliary assignment is at most

\[
C_JK_\Psi Q^4\xi^{-2}
 L^{2E_{\rm mark}+16+C_\psi}e^{-L^{3/10}}.
\tag{M23f}
\]

Here \(C_J=2J(3+1/\log2)\,64\cdot66\cdot3^{15}\); \(E_{\rm mark}\)
depends only on \(J\). Summing auxiliary probabilities and at most
\(L^{C_{\rm pat}}\) patterns preserves an error smaller than every fixed
inverse power of \(L\). Thus (M23e), summed over ranks and patterns,
equals that signed comparison up to the two specified errors (M28) and
(M23f).

For completeness, the Fourier pairing of (M23e) also has an exact form.
Put \(\psi_{\log}(t)=\Psi(e^t)\). Choose a fixed real smooth cutoff
\(\eta\), equal to one at both \(n/Y\) and \((n+2kc)/Y\) whenever the
summand in (M23e) is nonzero, and supported in a fixed compact
subinterval of \((0,\infty)\). Such a choice is uniform because
\(1\le c/H'\le e^{h_A+h_B}\) and \(|k|/Z\to0\) uniformly. Set

\[
A_{r,s,Y}(n)=\frac{\eta(n/Y)}{n}n^{-2\pi is}U_r(n),\qquad
B_{r,Y}(n)=\eta(n/Y)V_r(n),
\tag{M23g}
\]

and use \(\widehat A(\theta)=\sum_n A(n)\mathrm e(n\theta)\), likewise
for \(B\). Then (M23e) is

\[
\int_{\mathbb R}\widehat\psi_{\log}(s)
 \int_0^1\mathcal K_{r,s}(\theta)
       \overline{\widehat A_{r,s,Y}(\theta)}
       \widehat B_{r,Y}(\theta)\,d\theta\,ds,
\tag{M23h}
\]

where the full coupled kernel is

\[
\mathcal K_{r,s}(\theta)=
\frac{s_{a,b}Z^{-2\pi is}}{\nu_{j_A}\nu_{j_B}}
\sum_{\substack{u\in I_{j_A}\\v\in I_{j_B}}}
\nu_A(u)\nu_B(v)\overline{\chi_A(u)\chi_B(v)}
(uv)^{i\zeta-2\pi is}\mathrm e(-2kuv\theta).
\tag{M23i}
\]

The circle integral imposes \(n'=n+2kuv\), and the Mellin integral
restores \(\Psi(n/(Zuv))\), proving (M23h) term by term. The cell masses
in the denominator cancel the total restricted product probability, so
\(|\mathcal K_{r,s}(\theta)|\le1\) for every \(s,\theta\). On minor arcs
(M7)--(M9) give the stronger saving, retaining both inverse cell masses
in the coefficient budget (M23a). The endpoints (M23g) are now
independent of \(u,v\), which is the precise interface needed for the
arc-energy estimate.

The coefficient \(\alpha\) has the complete-cell discrepancy (C3): for
any prescribed \(B\), the fine mesh can be chosen so that

\[
\left|\sum_{a\le m\le b}
       \frac{\alpha(m)\chi(m)m^{it}}m\right|\le L^{-B},
\qquad q,|t|\le L^B,\quad [a,b]\subset[1,x^2]. \tag{M24}
\]

Here \(\alpha\) already includes its restriction to the factor interval
\(K\subset[x^\beta,x^2]\) of multiplicative width at most two, with
fixed \(\beta>0\). Thus the displayed sum means its intersection with
\(K\). If nonempty, that intersection satisfies the lower endpoint and
width conditions of (C3). The modulus \(q\) in this display is an
integer conductor variable, distinct from the fixed damping \(q=1/2\).

Let the two nonempty probe-product scales be \(U_1,U_2\), put
\(H'=U_1U_2\) and \(Y=ZH'\), and suppose \(Z\ge xL^{-C_Z}\). The marked
arc estimate obtained from (M24) is

\[
\int_{|\beta|\le L^{C_{\rm arc}}/H'}
 |\widehat E(h/d+\beta)|^2\,d\beta
 \le K_{\rm arc}YL^{-T}. \tag{M25}
\]

Here \(Q_*\) is the complete prime interval used by the comparison,
contained in \([x^\tau,x^\eta]\) with \(0<\tau<\eta<1/4\); the support
of \(\alpha\) is \(e^{\sqrt L}\)-rough. The cutoff has the prescribed
polynomial bounds on its logarithmic derivatives, and
\(2\le H'/L^{C_{\rm arc}}\le Y/(8A_\psi)\). These conditions hold for
the actual data from (B16). Under them, (M25) is uniform for the
polynomial conductor, frequency, and coefficient budgets in the
comparison expansion. The support has \(e^{3L/4}\le Y\le e^{3L/2}\);
take \(N_{\rm cap}=\lceil A_\psi Y\rceil\le x^2\), where \(A_\psi\)
bounds the cutoff support.

Here are the estimates establishing (M25). Peeling one distinguished
SMALL mark leaves the core and parity at the original integer and
evaluates only the remaining marked weight at \(n/p\). Separate the
smooth logarithmic factors by Fourier inversion, and write
\(H=H'/L^{C_{\rm arc}}\). The finite local-energy inequality (E1) and
the exact separation (E9) in
\protect\hyperlink{local-energy}{Section~\ref{local-energy}} give the
following four terms. With \(B_{\rm amp}\ll L^{E_{\rm amp}}\) and
\(B_s\ll L^{E_s}\) the resulting amplitude bounds, those terms, up to
fixed constants, are

\[
\begin{aligned}
&YV_s^{-2}d^2(1+\log_2N_{\rm cap})^8
                   L^{-(T+2c_m+8)},\\
&Y^{-1}L^{-(j-2)}
 \bigl(LY/H+16N_{\rm cap}(1+\log N_{\rm cap})\bigr)
 B_{\rm amp}^2N_{\rm cap}(1+\log N_{\rm cap})^{63},\\
&(H/Y)^2 B_{\rm amp}^2N_{\rm cap}(1+\log N_{\rm cap})^{63},\\
&3L B_s^2N_{\rm cap}(1+\log N_{\rm cap})^{16}e^{-L^{a_s}}.
\end{aligned} \tag{M26}
\]

The first is the finite Mellin-window contribution (E13), obtained from
(M24) and the small- and long-prime estimates, with \(d\le L^{c_m}\).
The second is the weighted Mellin tail; choose
\(j=\lceil2E_{\rm amp}+64+T\rceil+2\). The third is the error (E3) in
replacing linear windows by logarithmic windows. Since \(Y=ZH'\),
\(H/Y=1/(ZL^{C_{\rm arc}})\), so its exponential decay pays for every
displayed logarithmic factor. The fourth accounts for squares of a
peeled SMALL prime, whose size is at least \(e^{L^{a_s}}\); here
\(a_s=1/10\) is available. Divisor-moment estimates give the displayed
powers \(8,63,16\). Each line of (M26) is \(O(YL^{-T})\) with the stated
choices, proving (M25). The window replacement follows from the Taylor
estimate (E2) and finite overlap, with no prime distribution estimate on
an interval of length \(H\).

The first Fourier polynomial in the actual comparison has an additional
factor \(1/n\). With its cutoff written as \(\eta(n/Y)\),

\[
\frac{\eta(n/Y)}n=Y^{-1}\frac{\eta(n/Y)}{n/Y}.
\]

The cutoff is supported away from zero. Applying (M25) to \(\eta(u)/u\)
therefore gives energy \(K_{\rm arc}Y^{-1}L^{-T}\) for this first
polynomial; the second has energy \(K_{\rm arc}YL^{-T}\).
Cauchy--Schwarz cancels the two powers of \(Y\). Summing the rational
arcs of the comparison gives \(2K_{\rm arc}L^{2c_m-T}\). Off those arcs,
the auxiliary-product bilinear exponential-sum estimate gives
\(L^{-A_1}\) against the global polynomial endpoint energy. Choose
\(T,A_1\) to pay for the finite rank and pattern budgets. This makes the
whole comparison term smaller than any requested logarithmic power.

Subtracting this separately bounded term from the signed-operator
estimate (M20) bounds the original correlation.

\Needspace{6\baselineskip}

\hypertarget{marked-correlation-caps-shared-divisors-and-the-parameter-order}{%
\subsection{Caps, shared divisors, and the parameter
order}\label{marked-correlation-caps-shared-divisors-and-the-parameter-order}}

The finite comparison first imposes a cap on the marked part and allows
marks shared between the two endpoints. Both modifications are removed
quantitatively. The Rankin estimate from the weight chapter gives

\[
\sum_{\substack{u\le Y\\u_P>\lfloor e^{L/88}\rfloor}}W_1(u)
 \le YL^{(3+1/\log2)\log(6e^{12})}e^{-L^{3/5}/88}. \tag{M27}
\]

Stripping shared divisors has error

\[
O\!\left(L^{8(J+1)\log2+2E_{\rm mark}+17+C_\psi}
                e^{-L^{1/10}}\right). \tag{M28}
\]

The quotient cap is transported through the exact finite identity before
using (M27); its final specialization has shared divisor one and zero
extra marks. Thus no assertion that the marked part is squarefree, or
that a cap is unchanged by arbitrary division, is needed. Both (M27) and
(M28) are negligible at every fixed logarithmic accuracy.

The choices can now be made in order. Fix the desired saving and all
external frequency, shift, and cutoff bounds. Choose the local
contraction precision and \(m\) using (M1)--(M3), then the arithmetic
mesh, conductor, and retained-mass parameters using (M4)--(M10). Choose
\(J\) using (M13). The forest row bounds, total path mass, comparison
rank, Fourier-tail orders, and CRT precision are chosen after these
fixed parameters. Finally choose the discrepancy order \(B\) and the
fine-cell exponent \(S=2(B+1)+4\).

These choices precede the positive lower exponent \(\beta\) of the fine
interval and its positive roughness exponent \(\gamma\). Only the
eventual threshold in \(L\) depends on \(\beta,\gamma\). Combining the
endpoint bound (M20), the comparison estimate (M25), and the two
restoration errors (M27)--(M28) gives, with precisely the coefficients
of (B16)--(B17),

\[
|\mathcal C_{k,Z}|\le ZL^{-D_1}
\]

uniformly for \(0<|k|\le L^{C_k}\), \(xL^{-C_Z}\le Z\le xL^{C_Z}\), the
stated polynomial frequency bounds, and either multiplicity-parity
choice. This is (B18).

\clearpage
\hypertarget{physical-lift}{%
\section{From one mark to the physical matrix}\label{physical-lift}}

This chapter supplies the finite change of variables between (B17) and
the matrix (M18c). Extra marks are introduced by a common dilation of
the two integers. They are then removed by an exact finite reindexing
and a quantitative estimate for collisions.

\Needspace{6\baselineskip}

\hypertarget{physical-lift-the-raw-shared-prime-sum}{%
\subsection{The raw shared-prime
sum}\label{physical-lift-the-raw-shared-prime-sum}}

Put \(q=1/2\), \(M=J+1\), and let \(G\) denote the set of all marked
groups. For each \(g\in G\) draw \(J\) labelled primes
\(s_{g,i}\in P_g\) independently with probabilities \(1/(V_gs_{g,i})\).
Write

\[
D_s=\prod_{g,i}s_{g,i},\qquad
\mu(s)=D_s^{-1}\prod_gV_g^{-J},\qquad \sum_s\mu(s)=1.
\]

Repeated draws are included. For positive \(w\), let \(\mathcal A_s(w)\)
be the set of complete ordered lists of \(M\) distinct primes in each
group, all dividing \(D_sw\), whose first \(J\) coordinates are the
given \(s_{g,i}\). In particular this set is empty if any shared list
repeats a prime. Put

\[
H(n)=\prod_g q^{\max(\omega_g(n)-M,0)/2},
\qquad
R_s(w,w')=
\frac{H(D_sw)^2H(D_sw')^2}{D_s\prod_gV_g^{J+2}}
 |\mathcal A_s(w)|\,|\mathcal A_s(w')|. \tag{L1}
\]

On a nonempty completion the maximum in \(H\) can be omitted. Formula
(L1) is exactly the raw physical slice: the state mass contributes
\(\prod V_g^{-M}\), the raw row contributes \(\prod V_g^{-1}\), and its
two half dampings combine with the two endpoint half dampings. Averaging
the endpoint permutations gives the same fixed-slot completion sum; each
relabelling has the inverse permutation as its unique inverse. The
factor \(D_s^{-1}\) comes from the harmonic factor \(1/n\) after putting
\(n=D_sw\).

Call a draw clean if its primes are distinct in each group and none
divides \(ww'\). Then \(\omega_g(D_sw)=J+\omega_g(w)\) and each
completion is specified by one prime divisor of \(w\) in each group.
Substitution in (L1) gives

\[
R_s(w,w')=\mu(s)W(w)W(w')\quad\text{on clean draws}. \tag{L2}
\]

This explains why an original integer with only one marked divisor in a
group is allowed: its physical integer is \(D_sw\), which has the
additional \(J\) shared marked divisors on a clean draw.

Let \(I=J|G|\), \(P_{\min}\le\min_{p\in\cup P_g}p\), and
\(v\le\min_gV_g\). The probability of a collision or a prime dividing
\(ww'\) is at most

\[
\frac{I(\log_2 w+\log_2 w')+I^2}{vP_{\min}}. \tag{L3}
\]

Indeed every atom has mass at most \(1/(vP_{\min})\); a positive integer
has at most its base-two logarithm distinct prime divisors. Sum these
bounds over the draws and over pairs of draws. On any draw, completions
can only lose free choices compared with the one-mark lists of \(w,w'\).
At most \(J\) shared primes per group already divide either integer, so
removing the associated damping costs at most \(q^{-2M|G|}\). Therefore

\[
\left|\sum_sR_s(w,w')-W(w)W(w')\right|
\le(q^{-2M|G|}+1)W(w)W(w')
 \frac{I(\log_2 w+\log_2 w')+I^2}{vP_{\min}}. \tag{L4}
\]

This is obtained by summing \(R_s-\mu(s)W(w)W(w')\), which vanishes on
clean draws by (L2). Thus no conditioning or renormalization of the
prime law occurs. Since \(J\) is fixed, \(|G|=O(\log L)\) and
\(P_{\min}\ge e^{L^{1/10}}\), the factor in (L4) is a fixed power of
\(L\) times \(e^{-L^{1/10}}\) when \(w,w'\le e^{2L}\).

\Needspace{6\baselineskip}

\hypertarget{physical-lift-the-arithmetic-factors-and-the-exact-reindexing}{%
\subsection{The arithmetic factors and the exact
reindexing}\label{physical-lift-the-arithmetic-factors-and-the-exact-reindexing}}

Write \(a_1=\tau+\xi_3\), \(a_2=-\tau\) for the phases in (B17), and set
\(E=\overline F\). Then the product of the two cores is
\(\overline{E(w)}G(w')\). Let \(\Psi(t)=t\psi(t)\) and define

\[
K(w,w')=\frac{\Psi(w/Z)}w\,w^{ia_1}(w')^{ia_2}
 \overline{E(w)}G(w')\varepsilon(w)\varepsilon(w').
\]

Consequently
\(\mathcal C_{k,Z}/Z= \sum_{w>0,\,w+2k>0}K(w,w+2k)W(w)W(w+2k)\). The sum
is finite by the support of \(\psi\).

For a raw frame set \(n=D_sw\), \(n'=D_sw'\). Its equation
\(n'-n=2kD_s\) is equivalent to \(w'-w=2k\); the change of variables is
bijective on the positive integers divisible by \(D_s\). Moreover

\[
(D_sw)_*=w_*,\qquad
\varepsilon(D_sw)\varepsilon(D_sw')
 =\varepsilon(w)\varepsilon(w'). \tag{L5}
\]

The first identity uses only that \(D_s\) is supported on marked primes;
the second uses complete multiplicativity and \(\varepsilon(D_s)^2=1\).
They hold even on dirty draws and preserve the original multiplicities.
Taking the frame phase \(D_s^{i\zeta}\) with \(\zeta=-(a_1+a_2)=-\xi_3\)
also cancels the common dilation in the two imaginary powers. The cutoff
is evaluated at \(n/(ZD_s)=w/Z\). Expanding the raw physical pairing and
reindexing therefore gives exactly

\[
Q_{\rm raw}=
\sum_{w>0,\,w+2k>0}K(w,w+2k)\sum_s R_s(w,w+2k). \tag{L6}
\]

Choose the finite physical window to contain these dilated positions and
their edge neighbours. This is possible uniformly: \(D_s\) is at most
\(\exp(O_J(L^{2/5}\log L))\), while \(w\asymp Z=xL^{O(1)}\). No
contributing term is added or lost when the common finite carrier
replaces the individual quotient carriers.

Applying (L4) in (L6), the pointwise logarithmic bounds for \(W\), and
the divisor-moment bounds for the cores shows that
\(Q_{\rm raw}-\mathcal C_{k,Z}/Z\) is smaller than any fixed inverse
power of \(L\). For example, Cauchy--Schwarz bounds the sum of the two
core moduli over \(w\asymp Z\) by \(ZL^{O(1)}\), and the factor \(1/w\)
cancels \(Z\). The remaining factor is \(L^{O_J(1)}e^{-L^{1/10}}\).
Physical marked-part caps are removed separately by (M27); one must not
identify a cap at \(D_sw\) with the same cap at \(w\).

\Needspace{6\baselineskip}

\hypertarget{physical-lift-the-endpoint-vectors-and-the-signed-difference}{%
\subsection{The endpoint vectors and the signed
difference}\label{physical-lift-the-endpoint-vectors-and-the-signed-difference}}

Partition the first integer position into dyadic intervals \([X,2X)\).
On a physical state \(s\) at position \(n\), put

\[
\begin{aligned}
f_{X,\nu}(s)&=\sqrt{m_0}\,1_{X\le n<2X}\frac Xn
 (n/Z)^{-2\pi i\nu}n^{-ia_1}E(n)\varepsilon(n)H(n),\\
g_X(s)&=\sqrt{m_0}\,n^{ia_2}G(n)\varepsilon(n)H(n).
\end{aligned} \tag{L7}
\]

Restrict \(g_X\) to the finite neighbourhood reachable from the first
interval; this does not change the pairing. These formulas specify the
extra endpoint damping, separately from the damping already in the
matrix. The first vector depends only on the integer position.

Let \(f_\Psi(t)=\Psi(e^t)\) and use
\(\widehat f_\Psi(\nu)=\int_{\mathbb R}f_\Psi(t)e^{-2\pi i\nu t}\,dt\).
Fourier inversion of the actual frame cutoff, followed by finite
summation, gives

\[
Q_{\rm raw}-Q_{\rm cmp}
 =\sum_X\frac1X\int_{\mathbb R}\widehat f_\Psi(\nu)
   \langle f_{X,\nu},A(\zeta-2\pi\nu)g_X\rangle\,d\nu.
\tag{L8}
\]

Here \(A(v)\) is precisely (M18c) with frame phase \(D_f^{iv}\), and
\(Q_{\rm cmp}\) is the sum of its nonempty comparison patterns with the
signs in (M2). The conjugation of the first vector supplies
\((n/Z)^{2\pi i\nu}\); its product with \(D_f^{-2\pi i\nu}\) is
\(e^{2\pi i\nu\log(n/(ZD_f))}\). This proves (L8) without replacing
\(D_f\) by a representative cell value. Compact support and smoothness
make \(\widehat f_\Psi\) integrable. Fourier tails are bounded with the
fixed derivative budgets before the polynomial frequency restriction in
the correlation estimate is applied.

The vector norms in (M20) follow by first fixing the complete distinct
mark product \(P_0\). It divides \(n\), and the core has size at most a
fixed logarithmic factor times \(d_2(n/P_0)^2\). Thus its squared sum
over \(n\le3X\) is at most \((3X/P_0)L^{O(1)}\), using the fourth
divisor moment. Finally
\(m_0\sum_{\text{distinct lists}}1/P_0\le \prod_g V_g^{-M}(\sum_{p\in P_g}1/p)^M=1\).
This gives \(\|f_{X,\nu}\|,\|g_X\|\le\sqrt X L^{C_{\rm norm}}\) with
exponent independent of \(J\). The rough-factor counting bound for the
core gives the supremum bound in (M20). Applying that estimate in (L8),
paying for the logarithmic number of dyads and the Fourier norm, and
then using (L4)--(L6), proves the claimed saving for the original
correlation minus its comparison term.

\clearpage
\hypertarget{forest-moments}{%
\section{The forest expansion of the history sum}\label{forest-moments}}

We give the finite identity and the two estimates used in (M15). The
distinction between a new prime label and a later use of the same label
is essential: the forest expansion applies to the list of births, not to
the list of all visits.

\Needspace{6\baselineskip}

\hypertarget{forest-moments-a-linear-functional-on-birth-lists}{%
\subsection{A linear functional on birth
lists}\label{forest-moments-a-linear-functional-on-birth-lists}}

Fix the moment length \(N=2R\), the OLD cutoff \(T\), and all SMALL
labels. Expand the product of the \(N\) signed transitions, including
their comparison patterns and independent endpoint permutations. At the
initial state there are \(G(J+\ell)\) potential BIG birth slots. At each
of the \(N+1\) background steps there are \(GT\) slots, and each of the
\(N\) active transitions has \(G(J+\ell)\) slots. Here \(G\) is the
number of BIG groups. Thus the list has fixed length

\[
A=(N+1)G(J+\ell+T).
\]

An unused slot has value \(\varnothing\). Retaining a label already
present in the state creates no new birth. All local distinctness
conditions, transition legality conditions, OLD gates and the initial
distinctness condition remain in the coefficients of this expansion.
Only the additional condition that different births have different
numerical prime labels is set aside temporarily.

Let \(\mathcal H\) be this finite collection of expanded histories, let
\(b(H)=(b_0(H),\ldots,b_{A-1}(H))\) be its birth list, and let \(c_H\)
be the product of its transition and boundary coefficients, with the
common baseline and tilted prime masses described below. In particular
\(c_H\) includes the sign of every comparison pattern, all normalized
prime and permutation masses, the initial and terminal boundary factors,
and the damping factors. Define the linear functional

\[
\Lambda(f)=\sum_{H\in\mathcal H}c_H f(b(H)). \tag{F3}
\]

The projected history sum is
\(\Lambda(1_{\text{nonempty births distinct}})\). This construction does
not replace the coefficients by their absolute values. On histories with
distinct births its density agrees with the prime residue averages in
(M16). Factoring out the contribution of primes not selected as births
gives the common baseline \(B_{\mathrm{base}}=\prod_p b_p\) and the
tilted laws \(\widetilde\mu_g(p)=1/(V_gp b_p)\) from
\protect\hyperlink{marked-correlation-the-arithmetic-state-space-and-the-local-contraction}{the
state-space construction}.

For a birth list \(b\), replace an unused entry at address \(j\) by a
symbol \(\star_j\) used at no other address. Denote the resulting list
by \(x_0,\ldots,x_{A-1}\). A forest is a choice, for each \(j\), of
either no parent or one parent \(p(j)<j\). Its rank \(r(p)\) is the
number of chosen parents. Define

\[
T_p(b)=(-1)^{r(p)}\prod_{j:\,p(j)\text{ chosen}}
                  1_{x_j=x_{p(j)}},\qquad F_p=\Lambda(T_p).
\tag{F4}
\]

The exact identity is

\[
1_{x_0,\ldots,x_{A-1}\text{ distinct}}
 =\prod_{j=0}^{A-1}\left(1-\sum_{i<j}1_{x_j=x_i}\right)
 =\sum_p T_p(b). \tag{F5}
\]

For a distinct list every factor in the product is one. Otherwise, at
the first repeated entry the previous entries are distinct, so exactly
one summand is one and the product vanishes. Expanding the product gives
the second equality. The parent is always earlier, so cycles are
impossible. There are at most \(A^{2r}\) forests of rank \(r\): choose
at most \(A^r\) lists of children and at most \(A^r\) lists of parents.
Applying \(\Lambda\) to (F5) proves that the original signed history sum
is exactly \(\sum_p F_p\).

\Needspace{6\baselineskip}

\hypertarget{forest-moments-the-bound-preserving-signed-cancellation}{%
\subsection{The bound preserving signed
cancellation}\label{forest-moments-the-bound-preserving-signed-cancellation}}

The equality constraints in (F4) cannot be inserted into the local
contraction by conditioning. Instead use independent random sign fields.
For each selected edge \(e=(j,p(j))\), let \(\epsilon_e(a)\) be
independent uniform signs, indexed by the finite alphabet of primes and
private symbols. Then

\[
1_{x_j=x_{p(j)}}=
\mathbb E\bigl[\epsilon_e(x_j)\epsilon_e(x_{p(j)})\bigr].
\tag{F6}
\]

Independence between the edge fields gives the product of the equality
indicators. For a fixed field assignment, group its signs by address. To
make the multiplier at an unused slot equal to one, divide the
multiplier at address \(j\) by its value at \(\star_j\), and keep the
product of those values as an overall scalar. That scalar and every slot
multiplier have modulus one. Consequently \(F_p\) is exactly
\((-1)^{r(p)}\) times an average of the same history functional with
unit slot multipliers and a unit scalar outside it.

Assign the active birth slots in transition \(i\) to that transition;
assign background and initial slots to adjacent visits. Only the owners
of selected children and their parents can acquire a nontrivial slot
multiplier. There are at most \(2r(p)\) such owners. The background
operator bound is uniform in all unit multipliers. At an unmarked active
transition the original signed contraction still applies; at a marked
transition use the bound for an arbitrary unit multiplier.

Use exactly the orbit norm \(\mathcal E\) defined in (T1)--(T4) of
\protect\hyperlink{marked-correlation-the-arithmetic-state-space-and-the-local-contraction}{Section~\ref{marked-correlation-the-arithmetic-state-space-and-the-local-contraction}}.
Its OLD residues have counting measure and its ordered memory sectors
have the single factorial in (T3). The local estimates, proved in
(T6)--(T23), are

\[
\mathcal E(E_i f)\le l^2\mathcal E(f),\qquad
\mathcal E(E_i^{\phi}f)\le h^2\mathcal E(f)
\quad(|\phi|\le1),\qquad
\mathcal E(G_i^{\phi}f)\le\Gamma^2\mathcal E(f).
\tag{F7}
\]

Here \(E_i\) is the active signed transition, \(G_i\) is the background
transition, and \(\phi\) multiplies the actual transition coefficient.
In the notation (T21)--(T23), \(l=L^{-S_0}\), \(h=L^\kappa\),
\(b=L^{S_0+\kappa}\) and \(K_r=S_0+\kappa\). With the constants
\(W_0,w_0,\rho_0\) specified after (M15), the background bound is
\(\Gamma=\exp(4W_0w_0G/\rho_0)\).

Iterating (F7) and applying Cauchy--Schwarz to the initial and terminal
vectors bounds each fixed-field scalar by

\[
|B_{\mathrm{base}}|\,\mathcal E(v)\,
 \Gamma^{N+1}\prod_{i=0}^{N-1}
 \begin{cases}h,&i\text{ has a marked owner},\\l,&\text{otherwise}.
 \end{cases}
\]

The initial vector includes the retained initial distinctness mask and
the conjugated initial slot phase; its energy is at most that of the
original boundary vector \(v\). Thus, putting
\(P_0=|B_{\mathrm{base}}|\mathcal E(v)\Gamma^{N+1}\) and averaging the
unit scalars, we obtain

\[
|F_p|\le P_0 l^N b^{2r(p)}. \tag{F8}
\]

The initial boundary energy occurs once. In particular, (F8) does not
pay an additional factor for the number of histories or sign fields.

\Needspace{6\baselineskip}

\hypertarget{forest-moments-the-bound-from-prescribed-prime-values}{%
\subsection{The bound from prescribed prime
values}\label{forest-moments-the-bound-from-prescribed-prime-values}}

For the second bound keep the equality constraints and use absolute
coefficients. The tilted law satisfies

\[
\sum_p\widetilde\mu_g(p)\le2,\qquad
\sum_{p:\,p=t}\widetilde\mu_g(p)
 \le\delta\sum_p\widetilde\mu_g(p),\qquad
\delta=\frac2{v e^{L^{3/10}}}. \tag{F9}
\]

Expose births in address order. A selected child must equal a label
whose value has already been exposed. In an active transition with \(s\)
selected new coordinates the pinned absolute row bound is

\[
\sum_{\text{legal transitions with the prescribed values}}
 |c(x,y)|\,w_0^{|\mathrm{OLD}(y)|}
 \le \delta^s a\,w_0^{|\mathrm{OLD}(x)|}. \tag{F10}
\]

To obtain this bound, expand the comparison pattern, the auxiliary prime
choices and the endpoint permutations before summing the new
coordinates. Their masses are normalized. Each prescribed new value
contributes the atom bound in (F9); cap the character kernel pointwise
as in (M4). The remaining SMALL row sums and promotion damping give the
same absolute row constant \(a\) independently of the prescribed values.
For fixed \(m\) it satisfies \(a\le L^{C_b}\). This argument retains
multiplicities of transitions with the same output.

The corresponding background bound is \(\delta^s\Gamma\) with the same
weight \(w_0^{|\mathrm{OLD}|}\). Its geometric sum converges because
\(1/(\rho_0w_0)+\sqrt q/\rho_0^2\le1\). Apply these two bounds backwards
to successive suffixes of the history. Each selected child is charged
exactly once; a child forced to equal an incompatible group or a private
unused symbol gives zero. A selected child in the initial block also
gives zero, because initial distinctness was retained. Every surviving
forest therefore receives all \(r(p)\) atom factors, including forests
with several children of one parent.

The terminal boundary is bounded by \(w_0^{|\mathrm{OLD}|}\). The suffix
induction, followed by the initial boundary sum, yields

\[
|F_p|\le P_0 a^N\delta^{r(p)}. \tag{F11}
\]

No independence of the conditioned history is asserted in this argument:
(F10) is uniform in the already exposed state and prescribed values.
Combining (F8) and (F11) proves (M15), for the identical coefficients
\(F_p\) defined in (F3)--(F4). The rank split and removal of the OLD
cutoff then proceed as in
\protect\hyperlink{marked-correlation-restoring-all-histories}{the
history restoration}.

\clearpage
\hypertarget{local-energy}{%
\section{From complete-cell discrepancy to additive arc
energy}\label{local-energy}}

We supply the local-energy inequality and the separation of variables
used in (M26). The term with factor \((H/Y)^2\) is the error in
replacing a linear window by a logarithmic window. The separate error in
peeling a SMALL prime is supported on its square multiples.

\Needspace{6\baselineskip}

\hypertarget{local-energy-a-finite-local-energy-inequality}{%
\subsection{A finite local-energy
inequality}\label{local-energy-a-finite-local-energy-inequality}}

Let \((a_n)\) be a finite complex sequence supported on
\(Y/A\le n\le AY\), where \(A\ge1\), and assume \(1\le H\le Y/(8A)\).
Choose a fixed smooth function \(K\) supported on \([-\delta,\delta]\),
\(0<\delta\le1\), with Lipschitz constant \(M_K\) and

\[
|\widehat K(t)|\ge c_K>0\quad(|t|\le1).
\]

Such a function is obtained by shrinking a nonnegative smooth bump of
integral one. Its Fourier transform stays close to one on \([-1,1]\).
Put

\[
\begin{aligned}
P(\beta)&=\sum_n a_n\mathrm e(n\beta),\qquad
P_a(t)=\sum_n a_n n^{it},\\
Q(v)&=\sum_n a_nK((n-v)/H),\qquad
R(v)=\sum_n a_nK(v\log(n/v)/H),\quad v>0.
\end{aligned}
\]

For every integer \(j\ge2\) there is a constant \(C_j\), depending only
on \(A,K,j\), such that

\[
c_K^2\int_{-1/H}^{1/H}|P(\beta)|^2\,d\beta
\le \frac{C_j}{Y}\int_{\mathbb R}
 \frac{|P_a(t)|^2}{(1+(H/Y)|t|)^j}\,dt
 +2C_{\rm lin}(H/Y)^2\sum_n|a_n|^2,
\tag{E1}
\]

where \(C_{\rm lin}=4\delta(4\delta+1)(16A\delta^2M_K)^2\). The constant
is uniform in the sequence, \(H\) and \(Y\).

Indeed, \(\widehat Q(-\beta)=H\widehat K(-H\beta)P(\beta)\). Plancherel
bounds the left side by \(H^{-2}\int|Q(v)|^2\,dv\). The support of \(Q\)
lies in \(V=[Y/(2A),2AY]\). If either kernel in \(Q(v)-R(v)\) is
nonzero, then \(|n-v|\le2\delta H\). Taylor's estimate for
\(\log(1+u)\), using \(v\ge Y/(2A)\) and \(H\le Y/(8A)\), gives

\[
\left|K((n-v)/H)-K(v\log(n/v)/H)\right|
 \le16A\delta^2M_K H/Y. \tag{E2}
\]

At each \(v\) at most \((4\delta+1)H\) integers contribute; for each
\(n\) the contributing \(v\) lie in an interval of length \(4\delta H\).
Cauchy--Schwarz and integration give

\[
H^{-2}\int_V|Q(v)-R(v)|^2\,dv
 \le C_{\rm lin}(H/Y)^2\sum_n|a_n|^2. \tag{E3}
\]

To estimate \(R\), choose a fixed smooth compactly supported \(\chi\)
equal to one on \(\log(V/Y)\), set \(h=H/Y\), and write

\[
\mathcal Q(u)=\int_{\mathbb R}
 \chi(u)e^{-u}\widehat K(he^{-u}t/(2\pi))
 e^{-itu}Y^{-it}P_a(t)\,dt.
\]

Fourier inversion with \(t=2\pi v\xi/H\) gives exactly

\[
\chi(\log(v/Y))R(v)=\frac H{2\pi Y}
                   \mathcal Q(\log(v/Y)). \tag{E4}
\]

The amplitude is smooth and compactly supported in \(u\). Its first two
\(u\)-derivatives, multiplied by any fixed power of \(1+h|t|\), are
bounded uniformly in \(h,t\), by Schwartz decay of \(\widehat K\). Two
integrations by parts bound its Fourier transform by
\(C(1+h|t|)^{-j}(1+|s|)^{-2}\). Including \(e^{-itu}\) replaces the last
factor by \((1+|s+t/(2\pi)|)^{-2}\). Its row and column integrals are
bounded constants. Weighted Schur and Plancherel yield

\[
\int|\mathcal Q(u)|^2\,du
 \le C'_j\int(1+h|t|)^{-j}|P_a(t)|^2\,dt.
\]

The Jacobian of \(v=Ye^u\) is \(Ye^u=O(Y)\) on the support of \(\chi\).
Together with (E4) and division by \(H^2\), this leaves precisely
\(Y^{-1}\). Now \(|Q|^2\le2|R|^2+2|Q-R|^2\) proves (E1).

\Needspace{6\baselineskip}

\hypertarget{local-energy-exact-separation-of-the-arithmetic-coefficient}{%
\subsection{Exact separation of the arithmetic
coefficient}\label{local-energy-exact-separation-of-the-arithmetic-coefficient}}

Use (M21)--(M23) at rational center \(h_0/d\). Denote the distinguished
SMALL prime set by \(S\) and its harmonic mass by \(V_s\). After peeling
its distinguished mark, the main coefficient is

\[
a_n=\psi(n/Y)n^{iv_1}\mathrm e(h_0n/d)\varepsilon(n)
 \operatorname{Core}(n_*)\frac1{V_s}
 \sum_{s\in S,\ s\mid n}W_{\rm lo}(n/s).
\tag{E5}
\]

The weight \(W_{\rm lo}\) has one fewer mark in \(S\) and keeps all
other groups and the bounded list coefficient \(b\). Expanding the core
writes \(n=msp r\), where \(m\) is rough, \(s\in S\) is the peeled
prime, and \(p\) lies in the complete ordinary prime slot. The integers
\(m,p\) have no marked prime factors, and all of \(m,s,p\) are coprime
to \(d\). Thus the remaining marked weight is exactly \(W_{\rm lo}(r)\),
while the parity factor is \(\varepsilon(s)\varepsilon(r)\). The main
term does not require \(s\) to occur only to the first power.

For \(d_0=\gcd(r,d)\) put \(q=d/d_0\). On units modulo \(q\) expand

\[
\mathrm e(h_0u/q)=\sum_{\xi\bmod q}c_{q,\xi}\xi(u),\qquad
c_{q,\xi}=\frac1{\varphi(q)}\sum_{u\in(\mathbb Z/q\mathbb Z)^\times}
 \mathrm e(h_0u/q)\overline{\xi(u)}.
\tag{E6}
\]

All characters, including principal and imprimitive ones, occur. We have
\(|c_{q,\xi}|\le1\); their total number over \(d_0\mid d\) is
\(\sum_{d_0\mid d}\varphi(d/d_0)=d\). Define

\[
d_{d_0,\xi}(r)=1_{\gcd(r,d)=d_0}\,
 c(r_*)W_{\rm lo}(r)\varepsilon(r)\xi(r/d_0). \tag{E7}
\]

This coefficient is independent of \(m,s,p\). In particular
\(c(n_*/mp)\) has become \(c(r_*)\); it is not absorbed into \(\alpha\).
The marked-weight bound gives \(|d_{d_0,\xi}(r)|\le C K^{g+1}\), where
\(g=O(\log L)\) is the group count and \(K\) depends only on the fixed
mark counts, damping and lower harmonic-mass bound. This is a fixed
polynomial in \(L\).

Partition each of \(m,s,p,r\) into canonical dyads. On one box put

\[
\begin{aligned}
S_\xi(t)&=\sum_{s\in S\text{ in the dyad}}
             \varepsilon(s)\xi(s)s^{-1+it},\\
P_\xi(t)&=\sum_{p\text{ in the slot dyad}}\xi(p)p^{-1+it},\\
A_\xi(t)&=\sum_{m\text{ in the factor interval and dyad}}
             \alpha(m)\chi_0(m)\xi(m)m^{-1+it},\\
D_{d_0,\xi}(t)&=\sum_{r\text{ in the dyad}}
             d_{d_0,\xi}(r)r^{-1+it},\\
F_{d_0,\xi}(t)&=S_\xi(t)P_\xi(t+\sigma_1)
                    A_\xi(t+\sigma_0)D_{d_0,\xi}(t).
\end{aligned} \tag{E8}
\]

With \(f(u)=e^u\psi(e^u)\), Fourier inversion of the product cutoff is

\[
\frac{P_a(t)}Y=
\frac1{V_s}\sum_{d_0\mid d}\sum_{\xi\bmod d/d_0}
 c_{d/d_0,\xi}\sum_{\rm boxes}
 \int_{\mathbb R}\widehat f(\nu)Y^{-2\pi i\nu}
 F_{d_0,\xi}(t+v_1+2\pi\nu)\,d\nu.
\tag{E9}
\]

All integer sums are finite. The cutoff restricts the product to
\([Y/A,AY]\subset[1,N_{\rm cap}]\), so the independent four dyads
neither omit nor add a nonzero term. The reciprocal factors in (E8)
multiply to \(1/n\); the factor \(e^u=n/Y\) in \(f\) explains \(P_a/Y\).
There are at most \((1+\log_2N_{\rm cap})^4\) boxes.

For \(T_0+|\sigma_0|\le L^B\), (M24) bounds \(A_\xi(t+\sigma_0)\) on
\(|t|\le T_0\). Its joint character has modulus at most
\(q_{\rm base}d\), within the requested budget. If the other coefficient
bounds are \(C_s,C_p,C_r\), their reciprocal sums on dyads are at most
\(2C_s,2C_p,2C_r\). Hence

\[
\int_{-T_0}^{T_0}|F_{d_0,\xi}(t)|^2\,dt
 \le2T_0\bigl[(2C_s)(2C_p)(2C_r)L^{-B}\bigr]^2.
\tag{E10}
\]

This applies discrepancy only after exact separation, with no
\(m\)-dependent coefficient left in the other factors.

\Needspace{6\baselineskip}

\hypertarget{local-energy-the-rest-of-the-mellin-window}{%
\subsection{The rest of the Mellin
window}\label{local-energy-the-rest-of-the-mellin-window}}

The window extends beyond polynomial frequencies, so (E10) alone does
not suffice. Put \(T=2LY/H\). Discard boxes whose product interval
misses the cutoff support: their contribution before Fourier separation
is zero. On each remaining box with small-prime scale \(P\), long-prime
scale \(N\), and other scales \(M,R\), the actual comparison cells imply

\[
T\le MNR,\qquad Y\asymp PMNR,\qquad
\lceil Y\rceil^{3/5}\le MR\le\lceil Y\rceil.
\tag{E11}
\]

Here \(\log P\in[L^{a_s},L^{b_s}]\) up to fixed factors,
\(a_s<b_s<3/10\), and \(x^\tau/2\le N\le x^\eta\), \(\eta<1/4\). The two
nonempty BIG probe products ensure \(H'\gg LP L^{C_{\rm arc}}\), giving
the first inequality. The lower bound for \(MR\) follows from
\(Y\ge e^{3L/4}\) and \(PN\le e^{o(L)}x^\eta\); the upper bound follows
from the product support. Fixed comparison constants are absorbed by the
strict exponent margins and a larger threshold.

Choose \(\delta=L^{-A_s}\). Where \(|S_\xi(t)|\le\delta\), collapse the
other three factors to a Dirichlet polynomial of length at most
\(8MNR\). Its mean square on \([-T,T]\) is bounded by

\[
[2T+128MNR(1+\log(8MNR))]
 \frac{(C_mC_pC_r)^2}{MNR}(1+\log(8MNR))^{16}.
\]

Integrating the off-diagonal terms uses
\(|\log(n/m)|^{-1}\le\max(m,n)/|n-m|\); the divisor moment bounds the
collapsed coefficients. Multiplication by \(\delta^2\) and a sufficient
choice of \(A_s\) give any required logarithmic saving.

Cover the remaining times by occupied unit intervals. The factorial
moment estimate for the small-prime polynomial implies that their number
in each class of indices modulo three is at most
\(\lceil Y\rceil^{1/20}\), eventually. Raise the polynomial to a power
of order \(\log Y/\log P\). Unique factorization bounds ordered
representations by the factorial of that power. The separated sampling
estimate then gives an exceptional count \(\exp(o(\log Y))\): the
factorial and threshold losses have logarithm
\(O((\log Y/\log P)\log L)=o(\log Y)\), and \(\log P=o(\log Y)\). This
proves the stated exponent for every fixed \(A_s\).

On the occupied intervals the residual satisfies the explicit bound

\[
\sum_{j\in\mathcal J}\int_j^{j+1}
 |A_\xi(t+\sigma_0)D_{d_0,\xi}(t)|^2\,dt
 \le10752(C_mC_r)^2(1+\log(4\lceil Y\rceil))^4.
\tag{E12}
\]

Indeed, its collapsed coefficients on \([U,4U]\), \(U=MR\), have squared
norm at most \((C_mC_r/U)^2\,4U(1+\log(4U))^3\). For \(s\) one-separated
sample times in \([-\lceil Y\rceil,\lceil Y\rceil]\), the complete Gram
row is at most

\[
448U(1+\log(2\lceil Y\rceil))+448s\sqrt{\lceil Y\rceil}.
\]

The first-derivative bound for \(\sum_{U\le n\le4U}n^{it}\) is
\(O(U/|t|)\) for \(1\le|t|\le U\); the second-derivative bound is
\(O(\sqrt{\lceil Y\rceil})\) for \(U<|t|\le2\lceil Y\rceil\). Their sum
at separated times gives this row estimate. Since
\(s\le\lceil Y\rceil^{1/20}\) and \(U\ge\lceil Y\rceil^{3/5}\), the
second term is absorbed by \(448U\). Schur's inequality gives
\(3584(C_mC_r)^2(1+\log(4\lceil Y\rceil))^4\). Choose independently in
each interval a residual maximizer and a point witnessing the large
small-prime sum. Indices in one class modulo three make both families
one-separated. Summing the three classes proves (E12).

On occupied intervals above a sufficiently large fixed power of \(L\),
the long-prime estimate (U1), proved below, gives
\(|P_\xi(t+\sigma_1)|\le K_D L^{-D}\) for every fixed \(D\), uniformly
for the polynomial conductor budget and \(|t+\sigma_1|\le x^2\).
Multiply it by (E12), using \(|S_\xi|\le2\). The other occupied
intervals lie at polynomial frequencies and are covered by (E10). Thus
the entire window has arbitrary logarithmic saving for each four-factor
box.

Its constants, and those of the factorial-moment and Gram bounds,
precede the cutoff and box endpoints. To use (E9), restrict \(|\nu|\) to
a fixed polynomial range, enlarge the phase budgets, and apply
Minkowski. Bound the omitted Fourier integral by repeated integration by
parts. The logarithmic derivative budget for \(\psi\) supplies the
required norms and tails. Cauchy--Schwarz over the \(d\) character terms
and four dyads then gives, for any prescribed \(D\) after requesting
sufficient precision in (M24),

\[
\int_{-2LY/H}^{2LY/H}|P_a(t)/Y|^2\,dt
 \le V_s^{-2}d^2(1+\log_2N_{\rm cap})^8 L^{-D}.
\tag{E13}
\]

\Needspace{6\baselineskip}

\hypertarget{local-energy-the-high-frequency-prime-input}{%
\subsection{The high-frequency prime
input}\label{local-energy-the-high-frequency-prime-input}}

\textbf{Long-prime estimate.} Fix \(0<\tau<\eta<1/4\) and \(C,A>0\).
There are \(B_0,K>0\) and \(x_0\ge3\) such that, for \(x\ge x_0\),
uniformly for \(x^\tau/2\le N\le x^\eta\), \(1\le q\le(\log x)^C\),
every Dirichlet character \(\chi\bmod q\), every interval
\(I\subset[N,2N]\), and \((\log x)^{B_0}\le|u|\le x^2\),

\[
\left|\sum_{\substack{p\in I\\p\ \mathrm{prime}}}
       \chi(p)p^{-1+iu}\right|\le K(\log x)^{-A}.
\tag{U1}
\]

Principal and imprimitive characters are included. In particular, the
estimate holds at \(u=t+\sigma_1\) when \(L^{B_0}+|\sigma_1|\le|t|\) and
\(|t+\sigma_1|\le x^2\). This lower bound on \(|t|\) explains the
polynomial-frequency region handled separately by (E10).

The analytic input is the Vinogradov--Korobov zero-free region

\[
L(s,\chi)\ne0\quad\text{for}\quad
\Re s\ge1-
 \frac1{10.5\log q+61.5(\log|\Im s|)^{2/3}
                              (\log\log|\Im s|)^{1/3}},
\quad q\ge3,\quad |\Im s|\ge10.
\tag{U2}
\]

This is \href{https://arxiv.org/pdf/2210.06457\#page=1}{Khale, Theorem
1.1}, with the displayed conductor and height range.

Put \(L=\log x\) and choose

\[
\epsilon=L^{-(A+5)},\quad H=L^{A+7},\quad c=L^{-1},\quad
Q=3L^C,\quad V=3e^{3L},
\]

\[
D=10.5\log Q+61.5(\log V)^{2/3}(\log\log V)^{1/3},
\qquad \delta=(100D)^{-1},\qquad B_0=A+10.
\tag{U3}
\]

For moduli one and two induce to modulus \(3q\). The extra Euler factor
is nonzero in the positive half-plane, so (U2) gives the needed region
for the original character as well. Smooth the interval \(I\) with
relative boundary width \(\epsilon\). The Mellin transform of this
cutoff has decay constant \(B_j\ll_j\epsilon^{-j}\le C_jL^{(j+1)(A+5)}\)
at order \(j\). Mellin inversion expresses the smoothed Mangoldt sum
with coefficient \(\Lambda(n)\chi(n)n^{-1+iu}\) using
\(-L'/L(1-iu+z,\chi)\) on \(\Re z=c\). Shift the central segment
\(|\Im z|\le H\) to \(\Re z=-\delta\). Here \(|u|\ge L^{A+10}\) keeps
the rectangle away from the principal pole and low heights. The
denominator \(D\) dominates the zero-free widths for all its conductors
and heights. The logarithmic-derivative bound in this smaller rectangle,
the cutoff decay, and \(N^c\le e^\eta\) bound the left side, the
horizontal sides, and the remaining right-line tails by

\[
\mathcal M_j\ll B_j\left(
 HLN^{-\delta}
 +(c+\delta)e^\eta LH^{-j}
 +e^\eta L^2H^{1-j}\right).
\tag{U4}
\]

These bounds are uniform in the interval endpoints. Since
\(D\ll_C L^{3/4}\) and \(\log N\ge\tau L-\log2\), \(N^{-\delta}\) is
smaller than every fixed inverse power of \(L\). Choose a fixed integer
\(j>2A+30\). The other terms in (U4) are also \(O(L^{-A-2})\):
substituting (U3) gives exponents at most \(-2j+2A+14\). Removing the
smoothing, deleting higher prime powers, and partial summation to remove
\(\log p\) give

\[
\left|\sum_{p\in I}\chi(p)p^{-1+iu}\right|
\le\frac1{\log N}\left[
 \frac{\mathcal M_j}{2\pi}
 +\frac{2\log(3N)}N(2\epsilon N+2)
 +\frac{4\log(4N)}{\sqrt N}\right].
\tag{U5}
\]

The last two terms are respectively the boundary collars and the prime
powers. The exponential lower bound on \(N\) makes (U5) \(O(L^{-A})\),
proving (U1). Finally \(|t+\sigma_1|\ge|t|-|\sigma_1|\) proves the
shifted assertion.

\Needspace{6\baselineskip}

\hypertarget{local-energy-the-four-terms-in-m26}{%
\subsection{The four terms in
(M26)}\label{local-energy-the-four-terms-in-m26}}

For the peeled coefficient \(|a_n|\le B_{\rm amp}d_2(n)^3\). With
\(N=\lceil AY\rceil\) this gives

\[
\sum_n|a_n|^2\le B_{\rm amp}^2N(1+\log N)^{63}, \tag{E14}
\]

because \(d_2(n)^6\le d_{64}(n)\) and
\(\sum_{n\le N}d_{64}(n)\le N(1+\log N)^{63}\). The finite
Dirichlet-polynomial mean-square estimate, applied on dyadic time
intervals and multiplied by \((1+(H/Y)|t|)^{-j}\), bounds the part
\(|t|>LY/H\) by

\[
4L^{-(j-2)}[LY/H+16N(1+\log N)]
 B_{\rm amp}^2N(1+\log N)^{63}. \tag{E15}
\]

In (E1), (E13) contributes \(Y\) times its right side, and (E15)
contributes \(Y^{-1}\) times its right side. Formula (E3), with (E14),
gives the third line of (M26).

The difference between the original coefficient and (E5) vanishes unless
\(s^2\mid n\) for a peeled prime \(s\in S\). The marked-weight bound and
the fourth divisor moment on square multiples bound its squared
coefficient mass by

\[
3L B_s^2N(1+\log N)^{16}e^{-L^{a_s}}. \tag{E16}
\]

The factor \(3L\) pays for the dyadic SMALL-prime ranges, and
\(e^{-L^{a_s}}\) comes from summing \(s^{-2}\) for \(s\ge e^{L^{a_s}}\).
Parseval bounds its arc energy by this full coefficient mass. Combining
(E13)--(E16) with (E1) and \(|u+v|^2\le2|u|^2+2|v|^2\) gives the four
terms in (M26), up to constants fixed before \(L\).

\clearpage
\hypertarget{congruences}{%
\section{Distribution in residue classes}\label{congruences}}

We prove the negative-parity estimate used in
\protect\hyperlink{sieve}{the sieve argument}. All weights, prime bands,
and normalizing factors are those of \protect\hyperlink{proof}{the main
proof}. Put \(x=e^L\), \(Y=\lfloor2x\rfloor\), and
\(X=X_{-1}=\sum_{u\le Y}w_{-1}(u)\). For odd squarefree \(d\), set

\[
r(d)=\sum_{\substack{u\le Y\\u\equiv-2\pmod d}}w_{-1}(u)
       -\frac X{\varphi(d)}.
\]

For fixed \(0<\vartheta<1/2\) and every \(D>0\), the result is

\[
\sum_{\substack{d\le x^\vartheta\\d\text{ odd, squarefree}}}|r(d)|
 \ll_{\Psi,\vartheta,D}xL^{-D}+XL^{-18}. \tag{G1}
\]

The proof below concerns \(w_{-1}\). No distribution assertion for the
other parity sign is needed here.

\Needspace{6\baselineskip}

\hypertarget{congruences-orthogonality-and-the-nonunit-correction}{%
\subsection{Orthogonality and the nonunit
correction}\label{congruences-orthogonality-and-the-nonunit-correction}}

Extend Dirichlet characters by zero on nonunits, and write

\[
S_d(\chi)=\sum_{u\le Y}w_{-1}(u)\chi(u),\qquad
N(d)=\sum_{\substack{u\le Y\\(u,d)>1}}w_{-1}(u).
\]

Since \(-2\) is a unit modulo an odd \(d\), character orthogonality
gives

\[
\varphi(d)r(d)
 =\sum_{\chi\ne\chi_0}\chi((-2)^{-1})S_d(\chi)-N(d).
 \tag{G2}
\]

The correction \(N(d)\) occurs because the principal character sums to
\(X-N(d)\), rather than \(X\). In particular,

\[
|r(d)|\le\frac1{\varphi(d)}
 \left(\sum_{\chi\ne\chi_0}|S_d(\chi)|+N(d)\right). \tag{G3}
\]

Let

\[
H(N)=\sum_{\substack{d\le N\\d\text{ squarefree}}}\frac1{\varphi(d)}.
\]

For a fixed \(u\), assign a noncoprime squarefree \(d\) to any prime
\(p\mid(u,d)\) and write \(d=pe\). Dropping the conditions \((p,e)=1\)
and \(e\le N/p\) gives

\[
\sum_{\substack{d\le N\\d\text{ squarefree}\\(u,d)>1}}
 \frac1{\varphi(d)}
 \le H(N)\sum_{p\mid u}\frac1{p-1}. \tag{G4}
\]

The squarefree totient sum satisfies
\(H(\lfloor x^\vartheta\rfloor)\le eL\) eventually. Every prime divisor
of a supported \(u\le2x\) is at least \(L^{20}\). Also
\(2^{\omega(u)}\le u\), where \(\omega\) counts distinct prime divisors.
Therefore, eventually,

\[
\sum_{p\mid u}\frac1{p-1}
 \le\frac{2\omega(u)}{L^{20}}
 \le\frac4{\log2}L^{-19}.
\]

Multiply (G4) by \(w_{-1}(u)\ge0\) and interchange the finite sums. This
proves

\[
\sum_{\substack{d\le x^\vartheta\\d\text{ odd, squarefree}}}
 \frac{N(d)}{\varphi(d)}
 \le\frac{4e}{\log2}XL^{-18}. \tag{G5}
\]

\Needspace{6\baselineskip}

\hypertarget{congruences-exposing-one-prime-variable}{%
\subsection{Exposing one prime
variable}\label{congruences-exposing-one-prime-variable}}

Reserve one labelled slot in \(Q_0\), whose prime values satisfy

\[
x^\alpha\le q\le x^{6\alpha/5},\qquad \alpha=1/44.
\]

The ordered-slot expansion of the ordinary coefficient gives an exact
convolution

\[
w_{-1}(u)=\Psi(u/x)
 \sum_{\substack{q\in Q_0\\q\mid u}}b_L(u/q). \tag{G6}
\]

Here \(b_L\) is the sum over the remaining ordinary slots and all marked
group factors, with the original factorial normalization and negative
parity projection. Deleting the slot changes the number of slot
variables, not that normalization. If a prime occurs several times, its
different slot assignments remain in this finite expansion. After
factorial normalization the ordinary-slot coefficient is at most one,
and the marked factors are bounded by a fixed power of \(L\). Hence, for
a fixed \(c\ge0\),

\[
0\le b_L(R)\le L^c. \tag{G7}
\]

If \(\chi\) modulo \(d\) is induced by a primitive character \(\chi^*\)
of conductor \(f\mid d\), put \(h=d/f\). Inserting (G6) gives the exact
sum

\[
S_d(\chi)=\sum_{\substack{q\in Q_0,\ R\ge1\\(q,h)=(R,h)=1}}
 b_L(R)\Psi(qR/x)\chi^*(q)\chi^*(R). \tag{G8}
\]

The two coprimality masks are retained. Since \(d\) is squarefree,
\((f,h)=1\) and \(\varphi(d)=\varphi(f)\varphi(h)\). Summing the
complementary level \(h\) costs at most
\(H(\lfloor x^\vartheta\rfloor)=O(L)\).

Partition the two variables into binary dyadic intervals \(q\asymp P\),
\(R\asymp U\). The support \(x\le qR\le2x\) leaves \(O(L)\) pairs of
scales, and they satisfy

\[
P\le x^{6\alpha/5},\qquad U\le4x^{1-\alpha},\qquad
x/4<PU\le2x. \tag{G9}
\]

It remains to bound the primitive-character sums on each such rectangle.

\Needspace{6\baselineskip}

\hypertarget{congruences-small-and-large-conductors}{%
\subsection{Small and large
conductors}\label{congruences-small-and-large-conductors}}

The two analytic inputs are the prime Siegel--Walfisz estimate and the
multiplicative large sieve. In the forms used here, Siegel--Walfisz
gives arbitrary logarithmic saving for prime partial sums twisted by a
nonprincipal character of modulus at most a fixed power of \(L\). The
large sieve bounds the primitive-character quadratic mean of a sequence
on an interval of length \(N\) by its squared \(\ell^2\) norm times
\(Q^2+\pi N\), with the usual \(f/\varphi(f)\) conductor weight for
\(f\le Q\).

Choose a conductor cutoff \(2^K\le L^{C_*}\). For \(f\le2^K\), apply
Siegel--Walfisz to the \(q\) variable in (G8). Abel summation
incorporates \(\Psi(qR/x)\); its supremum and total variation are fixed
constants. The excluded primes \(q\mid h\) number at most
\(\vartheta/\alpha\), since \(q\ge x^\alpha\) and \(h\le x^\vartheta\).
Thus the small-conductor contribution on one retained rectangle is at
most

\[
L^{C_*+c}\left(
 2xB_\Psi L^{-A}V_\Psi
 +4x^{1-\alpha}\frac\vartheta\alpha S_\Psi\right), \tag{G10}
\]

where \(B_\Psi,S_\Psi,V_\Psi\) are fixed nonnegative constants from the
prime partial-sum estimate and the cutoff. The exponent \(A\) may be
chosen arbitrarily large once \(C_*\) is fixed.

For \(F<f\le2F\), use Fourier inversion in the logarithmic variable to
separate \(\Psi(qR/x)\). If \(\psi(t)=\Psi(e^t)\), write
\(M_\Psi=\int_{\mathbb R}|\widehat\psi(s)|\,ds<\infty\). At each
frequency, Cauchy--Schwarz in the primitive characters, followed by the
large sieve for the two factors, yields

\[
\mathcal B_F\le M_\Psi L^c
 \sqrt{\frac{(4F^2+\pi P)P}{F}}
 \sqrt{\frac{(4F^2+\pi U)U}{F}}. \tag{G11}
\]

Here \(\mathcal B_F\) is the sum of the absolute bilinear character sums
weighted by \(1/\varphi(f)\), for the fixed complementary level \(h\)
and rectangle. The masks in (G8) only decrease the coefficient norms.
Expanding the right side with \(\sqrt{a+b}\le\sqrt a+\sqrt b\), using
(G9), and summing over \(2^K<F\le2^T\) gives

\[
\sum_F\mathcal B_F\le M_\Psi L^c\left(
8\cdot2^T x^{1/2}
+12\sqrt\pi\,T x^{1-\alpha/2}
+\frac{4\pi x}{2^K}\right). \tag{G12}
\]

For a requested exponent \(D\), take

\[
k=c+D+2,\qquad C_*=k+1,\qquad A=C_*+c+D+2,
\]

and choose integer \(K,T\) with

\[
L^k\le2^K\le L^{k+1},\qquad
x^\vartheta\le2^T\le2x^\vartheta.
\]

The two remaining costs, the complementary levels and the rectangle
count, are each \(O(L)\). They turn the first term of (G10) and the last
term of (G12) into \(O(xL^{-D})\). All other terms contain a fixed
negative power of \(x\) after division by \(x\): the first term in (G12)
uses exactly \(\vartheta<1/2\). These powers absorb every fixed
logarithmic factor. We conclude that

\[
\sum_{\substack{d\le x^\vartheta\\d\text{ odd, squarefree}}}
 \frac1{\varphi(d)}\sum_{\chi\ne\chi_0}|S_d(\chi)|
 \ll xL^{-D}. \tag{G13}
\]

Equations (G3), (G5), and (G13) prove (G1). At \(D=230\), the lower mass
bound \(X\gg xL^{-221}\) makes its right side \(o(X/L)\).

\clearpage
\hypertarget{block-sieve}{%
\section{The quantitative finite block sieve}\label{block-sieve}}

We prove the finite estimate used for the rectangle upper bound and the
one-dimensional sieves. Its constants are independent of the block
count.

\Needspace{6\baselineskip}

\hypertarget{block-sieve-statement}{%
\subsection{Statement}\label{block-sieve-statement}}

Let \(z>1\), let \(h\ge2\) be even, and let \(P_j\), \(0\le j<J\), be
finite sets of primes in \((z^{2^{-(j+1)}},z^{2^{-j}}]\). These blocks
are disjoint; put \(P=\bigcup_jP_j\). Let \(\mathcal X\) be a finite set
of objects with weights \(w(x)\ge0\) and bad-prime sets \(B(x)\).
Suppose \(X\ge0\), \(0\le g(p)\le1-\eta\) for \(p\in P\), where
\(\eta>0\), \(C\ge0\), and \(\sum_{p\in P_j}g(p)\le C\). For every
\(s\subseteq P\), assume

\[
\sum_{x\in\mathcal X}w(x)\prod_{p\in s}1_{p\in B(x)}
 =X\prod_{p\in s}g(p)+r\!\left(\prod_{p\in s}p\right). \tag{BS1}
\]

Write

\[
\begin{gathered}
V=\prod_{p\in P}(1-g(p)),\qquad D=z^{4h+2},\qquad
\mathcal E=\sum_{\substack{d\mid\prod_{p\in P}p\\d\le D}}|r(d)|,\\
K(C,\eta)=\frac{\exp(eC+C/\eta)}{1-e^{-1}},\qquad
\rho=K(C,\eta)e^{-h},\qquad
S=\sum_{x\in\mathcal X}w(x)1_{B(x)\cap P=\varnothing}.
\end{gathered}
\]

Then

\[
XV(1-\rho e^\rho)-\mathcal E
 \le S\le XV e^\rho+\mathcal E. \tag{BS2}
\]

Every divisor in \(\mathcal E\) is squarefree. Both \(J\) and \(h\) may
vary.

\Needspace{6\baselineskip}

\hypertarget{block-sieve-proof}{%
\subsection{Proof}\label{block-sieve-proof}}

Set \(k_j=(j+1)h\). For variables \(t_p\in[0,1]\), define

\[
\begin{aligned}
a_j(t)&=\prod_{p\in P_j}(1-t_p),\\
u_j(t)&=\sum_{n=0}^{k_j}(-1)^n
             \sum_{\substack{s\subseteq P_j\\|s|=n}}\prod_{p\in s}t_p,\\
e_j(t)&=\sum_{\substack{s\subseteq P_j\\|s|=k_j+1}}\prod_{p\in s}t_p.
\end{aligned}
\]

Even Bonferroni truncation gives \(a_j\le u_j\) and
\(0\le u_j-a_j\le e_j\). To verify it, take independent Bernoulli
variables of means \(t_p\) and average
\(\sum_{n=0}^k(-1)^n\binom{q}{n}=(-1)^k\binom{q-1}{k}\) for \(q>0\),
handling \(q=0\) separately. A telescoping expansion now gives

\[
\prod_j u_j-\sum_j e_j\prod_{i\ne j}u_i
 \le\prod_j a_j\le\prod_j u_j. \tag{BS3}
\]

Apply (BS3) with \(t_p=1_{p\in B(x)}\), multiply by \(w(x)\) and sum.
Both bounding polynomials have monomial coefficients of absolute value
at most one. A subset contributes to the upper polynomial only when its
degree in every block is at most \(k_j\). It contributes to the
subtracted sum only when exactly one block has degree \(k_j+1\) and all
other degrees are at most their cutoffs. These cases are disjoint, and
the exceptional block is unique.

Every surviving monomial has prime product at most \(D\): its exponent
of \(z\) is at most \(h\sum_{j\ge0}(j+1)2^{-j}+1=4h+1<4h+2\). Distinct
prime subsets have distinct products. Inserting (BS1) therefore costs at
most \(\mathcal E\) for each bounding polynomial, with no extra divisor
coefficients or factor \(J\).

Evaluate next at \(t_p=g(p)\). Abbreviate \(a_j=a_j(g)>0\). Since
\(\log(1-t)\ge-t/\eta\) for \(0\le t\le1-\eta\), \(a_j\ge e^{-C/\eta}\).
Moreover

\[
e_j(g)\le\frac{C^{k_j+1}}{(k_j+1)!}\le\exp(eC-k_j-1),
\]

by retaining one term in the series for \(\exp(eC)\). Thus
\(\delta_j=e_j(g)/a_j\) satisfies \(\sum_j\delta_j\le\rho\). From
\(a_j\le u_j(g)\le a_j(1+\delta_j)\) we obtain

\[
V\le\prod_j u_j(g)\le Ve^\rho,
\qquad
\sum_j e_j(g)\prod_{i\ne j}u_i(g)\le V\rho e^\rho.
\]

Together with (BS3) and the remainder bound, these inequalities prove
(BS2).

\Needspace{6\baselineskip}

\hypertarget{block-sieve-constants-in-the-applications}{%
\subsection{Constants in the
applications}\label{block-sieve-constants-in-the-applications}}

For \protect\hyperlink{rectangles}{rectangles}, the local estimates give
\(C=12\), \(\eta=2/9\), and \(h=2\). Hence
\(\rho=e^{12e+52}/(1-e^{-1})\) and \(D=z^{10}\), exactly as in (R7). The
block count may grow.

For the preliminary sieve on \(u+2\), the density is zero at \(2\) and
\(1/(p-1)\) at odd primes. The reciprocal-prime block bounds give
\(C=8\), and \(\eta=1/2\). Take \(h=2\lfloor1/(40b)\rfloor\). If
\(0<b\le1/1000\), then \(h\ge1/(40b)\), \(h\ge50\) and \(\rho\le1\).
Consequently

\[
\rho e^\rho\le C_s e^{-1/(40b)},\qquad C_s=3K(8,1/2).
\]

This is the constant in (S7). The level exponent satisfies
\(b(4h+2)\le1/5+2b\). For the conditioned sieve in (S8b), use the same
\(C,\eta\) and let an even depth tend to infinity while keeping
\(W^{4h+2}\) below the remaining level after conditioning on the prime.
Estimate (BS2) is uniform in this depth and the block count.

For ordinary rough integers the density is \(1/p\), with block bound
\(C=4\) and gap \(\eta=1/2\). Put \(K_0=K(4,1/2)\) and
\(C_{\rm sieve}=K_0e^{K_0}\). Since \(\rho=K_0e^{-h}\le K_0\), both
relative errors in (BS2) are at most \(C_{\rm sieve}e^{-h}\). Counting
multiples of each squarefree divisor has error at most one, so the total
error is at most \(2D\), as used in (S8d). Both \(C_s\) and
\(C_{\rm sieve}\) are fixed before \(b\) is chosen.

\clearpage
\hypertarget{rectangles}{%
\section{A positive estimate on factor rectangles}\label{rectangles}}

We prove the finite estimate underlying (S2) in
\protect\hyperlink{sieve}{the sieve argument}. The notation for the
ordinary bands \(Q_j\), marked groups \(P_g\), and harmonic factors
\(H_Q,H_P\) is fixed in \protect\hyperlink{weights}{the weight chapter}.
Write \(x=e^L\) and

\[
\theta_0=\frac32\frac1{44}2^{-20},\qquad z=x^{\theta_0},\qquad
Z=\lfloor z^{10}\rfloor.
\]

Let \(K\subset[A,B)\) and \(N\subset[C,E)\) be finite sets of positive
integers, with natural endpoints \(A\le B\), \(C\le E\). Set

\[
\mathcal A=(B-A)(E-C),\qquad T=(B-A)+(E-C).
\]

For \(2\le W\le z\), use the indicator \(R_W\) from (C1), which vanishes
at \(0\) and \(1\). The finite rough-factor sum is

\[
\mathcal R_\rho(K,N;W)
=\sum_{u\le Y}w_\rho(u)
 \sum_{\substack{m\in K,\ n\in N\\mn=u+2}}
 R_W(m)R_W(n).
\]

The argument only uses \(|\rho|\le1\), \(0\le\Psi\le1\), and
\(\Psi(0)=0\); the outer cutoff \(Y\) may be any natural number.
Uniformly in these data, for sufficiently large \(L\),

\[
\begin{aligned}
\mathcal R_\rho(K,N;W)\le{}&
 K_0\mathcal A H_QH_PV(W)^2V(z)\\
&+200^{400}H_Q
 \left(2TZx^{1/20}+4Z^2x^{1/10}\right)(1+\log Z)^3,
\end{aligned} \tag{R1}
\]

where \(V(y)=\prod_{p\le y}(1-1/p)\) and \(K_0\) is an absolute
constant. The proof keeps the additional factor \(V(z)\) in the main
term.

\Needspace{6\baselineskip}

\hypertarget{rectangles-assignments-and-random-permissions}{%
\subsection{Assignments and random
permissions}\label{rectangles-assignments-and-random-permissions}}

Since \(A_\rho\le A_0\), it suffices to majorize the unsigned weight.
Keep the twenty ordinary slots in each band \(j\ge20\), and choose one
marked prime \(p_g\in P_g\) from each group. For an assignment \(a\),
put

\[
D_a=\prod_{j=20}^d\prod_{i=1}^{20}q_{j,i}\prod_gp_g,
\qquad
\eta=\frac1{(20!)^{d+1-20}}\prod_g\frac1{M_g}.
\]

The geometric bound on the ordinary scales and the marked-band bounds
give \(D_a\le x^{1/20}\) eventually. Every prime divisor of \(D_a\) is
at least \(L^{20}\), so \(D_a\) is odd.

Let \(\mathcal P=\bigcup_gP_g\) and
\(U_a=\{p\in\mathcal P:p\nmid D_a\}\). For each subset \(S\subset U_a\),
give \(S\) weight \(2^{-|U_a|}\) and call its primes permitted. Ban
every odd prime at most \(z\) that neither divides \(D_a\) nor belongs
to \(S\). For a fixed \(u\), independent choices of permission give the
exact identity

\[
2^{-|U_a|}\sum_{S\subset U_a}
 \prod_{p\in U_a\setminus S}1_{p\nmid u}
=\prod_{p\in U_a}
 \begin{cases}1/2,&p\mid u,\\1,&p\nmid u.\end{cases} \tag{R2}
\]

When the marked prime is summed over the distinct divisors from a group,
(R2) recovers its weight \(\omega_g(u)2^{-(\omega_g(u)-1)}/M_g\). Prime
multiplicities remain in \(u\); the permissions concern distinct primes.

The omitted ordinary bands \(j<20\) consist of primes larger than \(z\).
Thus, while the full ordinary factorization is still retained, every
prime at most \(z\) in the ordinary part must divide the tail
assignment. This proves the deterministic exclusions needed for the mask
majorant. We may then drop the remaining nonnegative factorial
coefficients and enlarge the tail-assignment sum. Consequently the
weighted rectangle sum is bounded by the sum, over \(a\) and masks
\(S\), of \(\eta2^{-|U_a|}\) times the number of pairs in the rectangle
satisfying

\[
mn\equiv2\pmod{D_a},\qquad R_W(m)R_W(n)=1,
\qquad p\nmid mn-2\ \text{for every banned }p. \tag{R3}
\]

The assignment moments used to sum this bound are

\[
\begin{aligned}
\sum_a\frac\eta{D_a}
 &=\prod_{j=20}^d\frac{V_j^{20}}{20!}
 \le200^{400}H_Q,\\
\sum_a\eta&\le200^{400}x^{1/20}H_Q,\\
\sum_a\eta D_a&\le200^{400}x^{1/10}H_Q.
\end{aligned} \tag{R4}
\]

In the first identity the marked reciprocal sums cancel \(M_g\) exactly.
The inequality uses \(V_j^{20}/20!\ge200^{-20}\) for the twenty omitted
bands. The other two inequalities follow from \(D_a\le x^{1/20}\).

\Needspace{6\baselineskip}

\hypertarget{rectangles-local-densities-and-the-finite-sieve}{%
\subsection{Local densities and the finite
sieve}\label{rectangles-local-densities-and-the-finite-sieve}}

Fix an assignment and mask. Write \(D=D_a\). The base congruence
\(mn\equiv2\pmod D\) has exactly \(\varphi(D)\) residue pairs modulo
\(D\): choose the unit \(m\), then \(n=2m^{-1}\). At an odd prime
\(p\nmid D\), the bad residue pairs are \(mn=0\) when \(p\le W\), and
\(mn=2\) when \(p\) is banned. These events are disjoint and have
\(2p-1\) and \(p-1\) solutions, respectively. Thus their combined
density is

\[
g(p)=\frac{(2p-1)1_{p\le W}+(p-1)1_{p\text{ banned}}}{p^2},
\qquad 0\le g(p)\le7/9,\quad g(p)\le3/p. \tag{R5}
\]

For squarefree \(d\) coprime to \(D\), the Chinese remainder theorem
makes the densities multiplicative. Counting each residue class in the
two integer intervals gives

\[
\#\{\text{base pairs bad at every }p\mid d\}
=\frac{\varphi(D)\mathcal A}{D^2}\prod_{p\mid d}g(p)+r_a(d),
\]

with the explicit boundary estimate

\[
|r_a(d)|\le\tau(d)^2(2T+4Dd). \tag{R6}
\]

Indeed, the count in a residue class differs from its interval length
divided by the modulus by at most one in each coordinate; expansion of
the product leaves the two side terms and a corner term. Summing over
the allowed CRT classes and using the local bounds in (R5) gives (R6).

Apply the \protect\hyperlink{block-sieve}{finite block sieve} with
parameter \(h=2\). The reciprocal-prime block bounds give block mass at
most \(12\), and (R5) gives the uniform gap \(1-g(p)\ge2/9\). The finite
sieve is obtained by upper Bonferroni truncation in each block; it
bounds the sifted count by

\[
e^{\mathfrak r_2}\frac{\varphi(D)\mathcal A}{D^2}
 \prod_{\substack{3\le p\le z\\p\nmid D}}(1-g(p))
+\sum_{d\le Z}|r_a(d)|,
\qquad
\mathfrak r_2=\frac{e^{12e+52}}{1-e^{-1}}. \tag{R7}
\]

Here the sieve level is \(z^{10}\), so the integer cutoff is \(Z\). The
explicit constant follows from (BS2) with \(C=12\), \(\eta=2/9\). Its
block count may grow with \(L\); the displayed bound is uniform in that
count by the finite proof given there.

The elementary divisor estimate
\(\sum_{d\le Z}\tau(d)^2\le Z(1+\log Z)^3\) now gives

\[
\sum_{d\le Z}|r_a(d)|
 \le(2TZ+4DZ^2)(1+\log Z)^3. \tag{R8}
\]

The mask weights sum to one. Averaging (R8) over masks and then applying
the last two moments in (R4) gives exactly the error in (R1).

\Needspace{6\baselineskip}

\hypertarget{rectangles-averaging-the-euler-product}{%
\subsection{Averaging the Euler
product}\label{rectangles-averaging-the-euler-product}}

Put \(v_p=1-1/p\). Direct comparison of the four possibilities in (R5)
gives

\[
1-g(p)\le
v_p^{\,2\,1_{p\le W}+1_{p\text{ banned}}}
\left(1+\frac1{p(p-1)}\right). \tag{R9}
\]

The product of the correction factors over odd primes is at most
\(e^{1/2}\), since \(\log(1+t)\le t\) and
\(\sum_{n\ge3}1/(n(n-1))=1/2\). Before deleting assignment primes or
permitting marked primes, the remaining product is

\[
\prod_{3\le p\le z}v_p^{\,1+2\,1_{p\le W}}
=8V(W)^2V(z). \tag{R10}
\]

The factor eight restores the three omitted factors at the prime two.
Permitting \(p\in S\) removes one power of \(v_p\). Averaging these
removals over the masks gives

\[
2^{-|U_a|}\sum_{S\subset U_a}\prod_{p\in S}v_p^{-1}
=\prod_{p\in U_a}\left(1+\frac1{2(p-1)}\right)
\le\prod_gZ_g(1/2)\le H_P. \tag{R11}
\]

Deleting primes dividing \(D_a\) removes at most three powers per prime.
Their roughness and \(D_a\le x^{1/20}\) imply, eventually,

\[
\prod_{p\mid D_a}v_p^{-3}\le e. \tag{R12}
\]

For example, \(-\log(1-1/p)\le2/p\), and the number of distinct prime
divisors is at most \(\log D_a/\log2\); hence the logarithm of the left
side tends to zero uniformly.

Finally \(\varphi(D_a)/D_a^2\le1/D_a\). Combining (R4) and (R7)--(R12)
proves the main term in (R1), with the explicit possible value

\[
K_0=200^{400}\,8e^{3/2+\mathfrak r_2}.
\]

\Needspace{6\baselineskip}

\hypertarget{rectangles-normalization-by-the-actual-mass}{%
\subsection{Normalization by the actual
mass}\label{rectangles-normalization-by-the-actual-mass}}

Assume \(\mathcal A\le C_{\rm area}x\) and
\(T\le C_{\rm side}x^{27/50}\), with fixed constants. Since
\(10\theta_0\le1/20\), we have \(Z\le x^{1/20}\) and
\((1+\log Z)^3\le8L^3\) for \(L\ge1\). The error in (R1) is therefore at
most a fixed constant times

\[
H_QL^3(x^{16/25}+x^{1/5}), \tag{R13}
\]

where \(27/50+1/20+1/20=16/25\) and \(2/20+1/10=1/5\). The mass bound
\(X_\rho\gg xH_QH_P/L\), Mertens' estimate \(V(z)\asymp L^{-1}\), and
\(H_P\ge1\) show that

\[
X_\rho V(W)^2\gg xH_Q/L^3\qquad(2\le W\le z).
\]

The main term of (R1) is \(O(X_\rho V(W)^2)\), and the ratio of (R13) to
this scale is at most a constant times
\(L^6(e^{-9L/25}+e^{-4L/5})\to0\). Thus there is \(C_{\rm box}>0\) such
that

\[
\mathcal R_\rho(K,N;W)\le C_{\rm box}X_\rho V(W)^2
\qquad(\rho=\pm1).
\]

Both the constant and the eventual threshold precede the two signs, the
admissible threshold \(W\), the finite carriers, and their endpoints. In
particular this applies to \(W=e^{\sqrt L}\).

\clearpage
\hypertarget{sieve}{%
\section{Sieve estimates and the remaining composites}\label{sieve}}

We use the definitions of \(w_\rho\), \(X_\rho\), and \(\mathfrak S\) in
\protect\hyperlink{proof}{the main proof}. The common outer cutoff is
\(Y=\lfloor2e^L\rfloor\), because \(\Psi\) is supported in \([1,2]\).

\Needspace{6\baselineskip}

\hypertarget{sieve-local-densities}{%
\subsection{Local densities}\label{sieve-local-densities}}

Write

\[
V(y)=\prod_{p\le y}(1-1/p),\qquad
V_2(y)=\prod_{3\le p\le y}(1-1/(p-1)).
\]

Every prime factor of a supported integer \(u\) is at least \(L^{20}\).
For large \(L\) this makes \(u\), and therefore \(u+2\), odd. At an odd
prime the excluded unit residue is \(-2\), and the density of the
divisibility condition is \(g(p)=1/(p-1)\). The identity

\[
1-\frac1{p-1}=(1-1/p)\left(1-\frac1{(p-1)^2}\right)
\]

gives \(V_2(y)/V(y)\to\mathfrak S\). Consequently
\(V_2(x^b)\sim\mathfrak S e^{-\gamma_E}/(bL)\) for fixed \(b>0\). The
product defining \(\mathfrak S\) converges because
\(\sum_{p\ge3}(p-1)^{-2}<\infty\). It is at least one: each factor lies
in \((0,1)\), so comparison with all integer indices gives

\[
\prod_{p\ge3}\left(1-\frac1{(p-1)^2}\right)
\ge\prod_{n=2}^{\infty}(1-n^{-2})=\frac12.
\]

The congruence estimate used here is specifically for the negative
parity weight. Put

\[
r(d)=\sum_{d\mid u+2}w_{-1}(u)-\frac{X_{-1}}{\varphi(d)}.
\]

For every fixed \(0<\vartheta<1/2\), some \(C\ge0\) satisfies,
eventually,

\[
\sum_{\substack{d\le x^\vartheta\\d\text{ odd and squarefree}}}|r(d)|
\le C\bigl(xL^{-230}+X_{-1}L^{-18}\bigr). \tag{S1}
\]

The \protect\hyperlink{congruences}{residue-class chapter} proves the
stronger version with any fixed accuracy in place of \(230\). Character
orthogonality separates nonprincipal characters from the nonunit
correction. A reserved ordinary prime slot, Siegel--Walfisz, and the
multiplicative large sieve control the former; the \(L^{20}\)-rough
support controls the latter. By \(X_{-1}\ge c_0xL^{-221}\), the right
side of (S1) is \(o(X_{-1}/L)\): the two ratios are bounded by constants
times \(L^{-8}\) and \(L^{-17}\).

\Needspace{6\baselineskip}

\hypertarget{sieve-positive-rectangles}{%
\subsection{Positive rectangles}\label{sieve-positive-rectangles}}

Let \(A,B,C,E\) be natural numbers with \(A\le B\) and \(C\le E\), and
let \(K\subset[A,B)\) and \(N\subset[C,E)\) be finite sets of positive
integers. For a roughness threshold \(W\), define the nonnegative sum

\[
\mathcal R_\rho(K,N;W)
=\sum_{u\le Y}w_\rho(u)
  \sum_{\substack{m\in K,\ n\in N\\mn=u+2}}
    R_W(m)R_W(n).
\]

Fix positive constants \(C_{\rm area},C_{\rm side}\). If

\[
(B-A)(E-C)\le C_{\rm area}x,
\qquad (B-A)+(E-C)\le C_{\rm side}x^{27/50},
\]

then, for \(W=e^{\sqrt L}\), there is a constant \(C_{\rm box}>0\) such
that

\[
\mathcal R_\rho(K,N;W)\le C_{\rm box}X_\rho V(W)^2. \tag{S2}
\]

The constant and the sufficiently large threshold for \(L\) precede both
parity signs and all the finite carriers and endpoints.

This is the specialization \(W=e^{\sqrt L}\) of
\protect\hyperlink{rectangles-normalization-by-the-actual-mass}{the
normalized rectangle estimate}. Its extra sieve factor
\(V(x^{\theta_0})\) controls prime divisors of \(mn-2\), while the two
factors \(V(W)\) impose roughness on \(m,n\). That normalization proves
the stated uniformity before any balanced width is chosen.

\Needspace{6\baselineskip}

\hypertarget{sieve-counting-balanced-rectangles}{%
\subsection{Counting balanced
rectangles}\label{sieve-counting-balanced-rectangles}}

Put

\[
\mathcal B_\kappa=[x^{1/2-2\kappa},x^{1/2+2\kappa}],
\qquad0\le\kappa\le1/50.
\]

Partition each positive factor by its unique binary index
\(k=\lfloor\log_2m\rfloor\). Only rectangles containing a pair with
\(mn=u+2\) and \(\Psi(u/x)\ne0\) need to be retained. If \((k,j)\) is
such a pair of indices, its geometry implies

\[
x/4<2^k2^j\le4x,
\qquad
x^{1/2-2\kappa}<2^{k+1},\quad 2^k\le x^{1/2+2\kappa}.
\]

The number of first indices is at most \(4\kappa L/\log2+2\). For each
first index there are at most five second indices, because their powers
of two lie in an interval with ratio sixteen. Thus the number of active
rectangles is at most

\[
C_{\rm dyad}(\kappa L+1),\qquad
C_{\rm dyad}=5(4/\log2+2). \tag{S5}
\]

The active rectangles partition the supported factor pairs. Each has
area at most \(4x\) and side sum at most
\(2x^{1/2+2\kappa}\le2x^{27/50}\), so (S2) applies with fixed constants.

Summing (S2) and applying (S5) gives

\[
\sum_{\substack{mn=u+2\\m,n\in\mathcal B_\kappa\\
P^-(m),P^-(n)>W}}w_\rho(u)
\le C_1(\kappa L+1)X_\rho V(W)^2. \tag{S6}
\]

The constant \(C_1\) is chosen before \(\kappa\).

\Needspace{6\baselineskip}

\hypertarget{sieve-removing-the-unbalanced-composites}{%
\subsection{Removing the unbalanced
composites}\label{sieve-removing-the-unbalanced-composites}}

The approximation estimate in
\protect\hyperlink{correlation}{Section~\ref{correlation}} permits
replacement of the roughness test on the cofactor of a least prime
divisor. The number of bins is fixed before the approximation
parameters; the parameter theorem remains uniform over the subsequent
choice of endpoints.

For small fixed \(b>0\), define the rough-number density by

\[
D_b(w)=\frac1w+
\sum_{k\ge2}\frac1{k!}
\int_{\substack{t_1,\ldots,t_{k-1}\ge b\\
                 t_1+\cdots+t_{k-1}\le w-b}}
\frac{dt_1\cdots dt_{k-1}}
 {t_1\cdots t_{k-1}(w-t_1-\cdots-t_{k-1})}.
\]

Only finitely many terms can be nonzero, since the domain is empty when
\(kb>w\). The \protect\hyperlink{block-sieve}{finite block sieve} gives
the preliminary-sieve constant \(C_s=3\exp(8e+16)/(1-e^{-1})\),
independent of \(b\). Set

\[
A_b=\frac{e^{-\gamma_E}}b
       \left(1-C_s e^{-1/(40b)}\right).
\]

The preliminary sieve and (S1) give, eventually,

\[
\sum_{P^-(u+2)>x^b}w_{-1}(u)
\ge\left(\mathfrak S A_b-\frac{\mathfrak S}{10}\right)
       \frac{X_{-1}}L. \tag{S7}
\]

Choose a finite monotone partition of \((b,1/2-\kappa]\). For a bin
\((\gamma,\gamma']\), the cofactor test is \(1_{P^-(m)>x^\gamma}\) and
the other factor is a prime \(n\in(x^\gamma,x^{\gamma'}]\). The sum over
all such ordered pairs dominates the composites whose least prime factor
is in that bin: choose \(n=P^-(u+2)\) and \(m=(u+2)/n\). No uniqueness
of the other representations is needed because their weights are
nonnegative.

For a fixed bin \((\gamma,\eta]\), put

\[
C_{\gamma,\eta}=\sup_{\gamma\le t\le\eta}D_\gamma(1-t),
\qquad
F_{\gamma,\eta}=
\sum_{x^\gamma<p\le x^\eta}
\sum_{\substack{p\mid u+2\\P^-(u+2)>W}}w_{-1}(u),
\quad W=e^{\sqrt L}.
\]

The cofactor \(m=(u+2)/p\) has logarithmic exponent in
\([1-\eta,1-\gamma+o(1)]\). The full-cell roughness approximation
satisfies, for every fixed \(\delta>0\), eventually on these pairs,

\[
0\le B_{R_{x^\gamma}}(m)
\le\frac{C_{\gamma,\eta}+\delta}{L V(W)}. \tag{S8a}
\]

To obtain (S8a), count products of a fixed number of primes in each
complete logarithmic cell. The prime number theorem and partial
summation give the reciprocal-simplex integrals defining \(D_\gamma\).
There are only finitely many possible numbers of prime factors on the
fixed compact exponent range. The denominator in the cell ratio is its
length times \(V(W)(1+o(1))\), by the finite rough sieve. Joint
continuity of \(D_\gamma(w)\) on this compact range absorbs the \(o(1)\)
overshoot at \(1-\gamma\) and gives the displayed supremum bound.

If the cofactor approximation is nonzero, then \(m\) is \(W\)-rough.
Since the prime \(p>x^\gamma\) also exceeds \(W\), the product
\(u+2=mp\) is \(W\)-rough, even when \(p\mid m\). Reindexing the ordered
pair sum by \(p\) and \(u\) therefore bounds its approximated value by
\((C_{\gamma,\eta}+\delta)F_{\gamma,\eta}/(LV(W))\).

The conditioned sieve gives

\[
\frac{F_{\gamma,\eta}}{X_{-1}V(W)}
\longrightarrow\mathfrak S\log(\eta/\gamma). \tag{S8b}
\]

Here one chooses \(\eta<\vartheta<1/2\) in (S1). For a fixed
\(p\in(x^\gamma,x^\eta]\), sieve the condition \(p\mid u+2\) by primes
at most \(W\). Since \(p>W\), its local density separates as
\(g(pd)=g(p)g(d)\) for each sieve divisor \(d\). The main term is
\(X_{-1}g(p)V_2(W)\). Summing over \(p\) gives
\(\sum g(p)\to\log(\eta/\gamma)\) by the prime reciprocal-sum estimate.
The sieve level can be chosen below \(x^\vartheta/p\); its logarithm
divided by \(\log W\) tends to infinity uniformly in this bin. Thus its
relative sieve error tends to zero. In the remainder sum, each integer
\(pd\le x^\vartheta\) has at most \(\vartheta/\gamma\) possible prime
factors from the bin, a fixed bound. Estimate (S1) consequently controls
the summed remainders. Finally \(V_2(W)/V(W)\to\mathfrak S\), proving
(S8b).

Combining (S8a) and (S8b), and taking \(\delta\) sufficiently small,
gives for any \(\varepsilon>0\) the per-bin bound

\[
U^{\rm proxy}_{\gamma,\eta}
\le\left(\mathfrak S C_{\gamma,\eta}\log(\eta/\gamma)
          +\varepsilon\right)\frac{X_{-1}}L. \tag{S8c}
\]

For completeness, the choice of partition in this sum is independent of
all subsequent approximation parameters. Set \(c=1/2-\kappa\). The
function \(F(s,t)=D_s(1-t)\) is continuous on \([b,c]^2\). Given
\(\varepsilon>0\), uniform continuity permits an affine grid
\(t_i=b+(c-b)i/N\), with \(N>0\) fixed, for which

\[
\sum_{i<N}C_{t_i,t_{i+1}}\log(t_{i+1}/t_i)
\le\int_b^c D_t(1-t)\,\frac{dt}{t}+\varepsilon.
\]

Indeed, on each cell the supremum differs from \(F(t,t)\) by at most
\(\varepsilon/\log(c/b)\); integration against \(dt/t\) makes these
errors sum to at most \(\varepsilon\).

Partitioning each prime-factor simplex by its least coordinate gives

\[
\int_b^{1/2}D_t(1-t)\,\frac{dt}{t}=D_b(1)-1.
\]

The one-factor term is \(1\). In a \(k\)-factor term, choose its least
coordinate \(t\) and integrate the remaining \(k-1\) coordinates with
lower bound \(t\); the \(k\) choices change \(1/k!\) to \(1/(k-1)!\).
The sets where two coordinates tie have measure zero. All terms are
nonnegative and only finitely many can be nonzero, which justifies the
partition and integration. Restricting the integral to \([b,c]\) only
reduces it. Choose the grid above with \(\varepsilon=1/10\). Sum (S8c)
over this fixed grid and use the approximation error \(o(X_{-1}/L)\)
from \protect\hyperlink{correlation}{Section~\ref{correlation}}.
Allocating the finite per-bin errors and the replacement error gives

\[
U\le\left(\mathfrak S(D_b(1)-1+1/10)
                   +\frac{\mathfrak S}{10}\right)\frac{X_{-1}}L.
\tag{S8}
\]

The needed small-\(b\) estimate is the one-sided bound

\[
D_b(1)\le\frac{e^{-\gamma_E}}b
 \left(1+C_{\rm sieve}e^{-1/(40b)}\right),\qquad0<b\le1/100.
\tag{S8d}
\]

To prove this, let \(R_b(z)\) count the integers in \((0,z]\) with no
prime factor at most \(z^b\), and take the even sieve depth
\(h=2\lfloor1/(40b)\rfloor\). The finite rough sieve gives

\[
R_b(z)\le zV(z^b)(1+C_{\rm sieve}e^{-h})
                 +2z^{b(4h+2)}.
\]

For \(b\le1/100\), \(h\ge1/(40b)\) and \(b(4h+2)\le1/5+2b<1\). The last
term disappears after multiplication by \(\log z/z\). Mertens' formula
gives \((\log z)V(z^b)\to e^{-\gamma_E}/b\). On the other hand, counting
products of distinct primes in compact subsets of the simplices by the
prime number theorem gives

\[
D_b(1)\le\liminf_{z\to\infty}\frac{\log z}{z}R_b(z).
\]

To see the factorial, count ordered \(k\)-tuples and divide by \(k!\);
distinct tuples have exactly \(k!\) orderings. Repeated-prime tuples may
be discarded in this lower bound. Refine the compact grids and let their
boundaries tend to the simplex boundary; positivity gives the complete
integral. Again \(k\le1/b\) is bounded for fixed \(b\). Combining the
lower comparison with the upper sieve bound proves (S8d).

Now put \(b=1/(40T)\), with \(T\ge25\) chosen so large that

\[
40e^{-\gamma_E}(C_s+C_{\rm sieve})Te^{-T}\le1/10.
\]

Such a \(T\) exists because \(Te^{-T}\to0\). Then \(0<b\le1/1000\), and
(S8d) gives

\[
A_b-(D_b(1)-1+1/10)
\ge\frac9{10}
 -\frac{e^{-\gamma_E}}b(C_s+C_{\rm sieve})e^{-1/(40b)}
\ge\frac45.
\]

Subtracting (S8) from (S7) leaves at least
\((4/5-1/10-1/10)\mathfrak S X_{-1}/L\), which is at least
\((\mathfrak S/2)X_{-1}/L\).

Every composite surviving that subtraction has least prime factor
greater than \(x^{1/2-\kappa}\). The finite partition theorem accounts
for the strict lower and closed upper endpoints of every bin. Hence

\[
\mathcal M_{-1}+\mathcal C_{-1}(L,\kappa)
\ge\frac{\mathfrak S}{2}\frac{X_{-1}}L. \tag{S9}
\]

\Needspace{6\baselineskip}

\hypertarget{sieve-balanced-composites}{%
\subsection{Balanced composites}\label{sieve-balanced-composites}}

For sufficiently large \(x\), a composite \(u+2\le2x+2\) whose least
prime factor exceeds \(x^{1/2-\kappa}\) has exactly two prime factors,
counted with multiplicity: three factors would give
\(u+2>x^{3(1/2-\kappa)}>2x+2\). Both factors lie in
\(\mathcal B_\kappa\). The argument includes a square of a prime. Let
\(I(m)=1_{m\text{ prime}}\), and let \(B(m)\) be its full-cell
approximation. The complete-cell prime estimate gives
\(0\le B(m)\le3/(LV(W))\) on the balanced range, with rough support.
Therefore (S6) bounds the pair sum with two approximations by

\[
9C_1(\kappa+1/L)\frac{X_\rho}L.
\]

For fixed \(\kappa>0\), eventually \(1/L\le\kappa\). Thus the constant
\(C_{\rm bal}=18C_1\) works for the approximated sum, independently of
\(\kappa\) and of the cell precision.

There are exactly two replacement steps:

\[
I(m)I(n)-B(m)B(n)
=(I(m)-B(m))I(n)+B(m)(I(n)-B(n)).
\]

The companion in the second step is bounded but need not be smooth.
Their total is bounded by \(\mathcal E(L)=o(X_\rho/L)\) from
(C9)--(C12).

It follows that for every fixed \(0<\kappa<1/50\) and \(\varepsilon>0\),
eventually for both signs,

\[
\mathcal C_\rho(L,\kappa)
\le(C_{\rm bal}\kappa+\varepsilon)\frac{X_\rho}L. \tag{S10}
\]

The bound holds first for the ordered prime-pair sum and hence for the
composite mass it covers. For (S10), the auxiliary partition may be
fixed to one admissible bin with lower endpoint \(1/200\). Equations
(S9) and (S10) are the two estimates used in the main proof.


\appendix
\clearpage
\hypertarget{lean-guide}{%
\section{Reading the Lean proof}\label{lean-guide}}

The public theorem is {\small\nolinkurl{Green64.green_64}} in
\href{../Green64/Main.lean}{\nolinkurl{Main.lean}}. Its statement is the
direct proposition

\begin{verbatim}
{p : ℕ | p.Prime ∧ Odd (Ω (p - 2))}.Infinite
\end{verbatim}

The proof first constructs a prime above every natural bound.
{\small\nolinkurl{Green64.primes_unbounded}} records positivity of
\(p-2\) as well as its odd prime-factor count.
{\small\nolinkurl{scripts/StatementTests.lean}} checks both exact
propositions, including the multiplicity convention, and
{\small\nolinkurl{scripts/AxiomAudit.lean}} prints their dependency
axioms. Mathlib's {\small\nolinkurl{Ω}} is the length of
{\small\nolinkurl{Nat.primeFactorsList}}, so it counts prime factors
with multiplicity.

The project separates a short final argument from the supporting
analytic library. Module imports form an acyclic graph. Supporting
declarations remain source code in {\small\nolinkurl{Green64/Library}};
they are not external packages or precompiled proof objects. The final
argument uses the following modules.

\begin{longtable}[]{@{}
  >{\raggedright\arraybackslash}p{(\columnwidth - 4\tabcolsep) * \real{0.3333}}
  >{\raggedright\arraybackslash}p{(\columnwidth - 4\tabcolsep) * \real{0.3333}}
  >{\raggedright\arraybackslash}p{(\columnwidth - 4\tabcolsep) * \real{0.3333}}@{}}
\toprule\noalign{}
\begin{minipage}[b]{\linewidth}\raggedright
Mathematical step
\end{minipage} & \begin{minipage}[b]{\linewidth}\raggedright
Module
\end{minipage} & \begin{minipage}[b]{\linewidth}\raggedright
Main declaration
\end{minipage} \\
\midrule\noalign{}
\endhead
\bottomrule\noalign{}
\endlastfoot
Infinitely many primes & \href{../Green64/Main.lean}{\nolinkurl{Main}} &
{\small\nolinkurl{Green64.green_64}} \\
Witness above any bound &
\href{../Green64/Infinitude.lean}{\nolinkurl{Infinitude}} &
{\small\nolinkurl{prime_witnesses_unbounded}} \\
Positive mass gives a prime of the required parity &
\href{../Green64/PrimeExtraction.lean}{\nolinkurl{PrimeExtraction}} &
{\small\nolinkurl{negative_parity_prime_candidates}} \\
A fixed smooth cutoff &
\href{../Green64/Localization.lean}{\nolinkurl{Localization}} &
{\small\nolinkurl{localizingBump}} \\
Subtract balanced composites &
\href{../Green64/Sieve/PrimeMass.lean}{\nolinkurl{Sieve/PrimeMass}} &
{\small\nolinkurl{negative_prime_mass_lower_bound}} \\
Cover the composite remainder by prime pairs &
\href{../Green64/Sieve/CompositeMass.lean}{\nolinkurl{Sieve/CompositeMass}}
& {\small\nolinkurl{balanced_composite_mass_bound}} \\
Transfer the proxy bound to sharp prime indicators &
\href{../Green64/Sieve/BalancedMass.lean}{\nolinkurl{Sieve/BalancedMass}}
& {\small\nolinkurl{balanced_prime_pair_and_composite_bound}} \\
Bound the two prime proxies &
\href{../Green64/Sieve/BalancedApproximation.lean}{\nolinkurl{Sieve/BalancedApproximation}}
& {\small\nolinkurl{balanced_proxy_bound_uniform}} \\
Sum positive bounds over balanced rectangles &
\href{../Green64/Sieve/BalancedBoxes.lean}{\nolinkurl{Sieve/BalancedBoxes}}
& {\small\nolinkurl{balanced_rough_pair_bound}} \\
Uniform weighted rectangle bound &
\href{../Green64/Sieve/WeightedBoxes.lean}{\nolinkurl{Sieve/WeightedBoxes}}
& {\small\nolinkurl{weighted_box_bound}} \\
Linear count of compatible dyadic rectangles &
\href{../Green64/Sieve/DyadicCounting.lean}{\nolinkurl{Sieve/DyadicCounting}}
& {\small\nolinkurl{compatible_dyad_labels_card}} \\
Exact finite partition &
\href{../Green64/Sieve/DyadicPartition.lean}{\nolinkurl{Sieve/DyadicPartition}}
& the active-rectangle partition identities \\
Support and rectangle geometry &
\href{../Green64/Sieve/DyadicGeometry.lean}{\nolinkurl{Sieve/DyadicGeometry}}
& the bounds for attained dyadic scales \\
Normalize the finite main term &
\href{../Green64/Sieve/BoxNormalization.lean}{\nolinkurl{Sieve/BoxNormalization}}
& {\small\nolinkurl{box_main_Xrho_eventually}} \\
Bound the finite remainder &
\href{../Green64/Sieve/BoxRemainder.lean}{\nolinkurl{Sieve/BoxRemainder}}
& {\small\nolinkurl{box_finite_error_exp_bound}} \\
Make the remainder negligible &
\href{../Green64/Sieve/BoxError.lean}{\nolinkurl{Sieve/BoxError}} &
{\small\nolinkurl{box_error_uniform_little_o}} \\
Euler-product estimates &
\href{../Green64/Sieve/EulerProduct.lean}{\nolinkurl{Sieve/EulerProduct}}
& {\small\nolinkurl{box_euler_bounds}} \\
Account for both replacements &
\href{../Green64/Sieve/ApproximationError.lean}{\nolinkurl{Sieve/ApproximationError}}
& {\small\nolinkurl{balanced_sharp_le_proxy_error}} \\
Close the analytic parameter choices &
\href{../Green64/Sieve/Parameters.lean}{\nolinkurl{Sieve/Parameters}} &
{\small\nolinkurl{proxy_parameters_exist}} \\
\end{longtable}

Unless a row says otherwise, the public analytic declarations are in
namespace {\small\nolinkurl{Green64Analytic.Parity}}. Lower-level
declarations retain their specific namespaces. Search a declaration with
{\small\nolinkurl{rg}} if using an editor without Lean's
go-to-definition support.

\Needspace{6\baselineskip}

\hypertarget{lean-guide-correspondence-with-the-mathematical-text}{%
\subsection{Correspondence with the mathematical
text}\label{lean-guide-correspondence-with-the-mathematical-text}}

The following map records the exact source results used in the exposition.
The probability notation in the weight chapter rewrites normalized sums;
it does not introduce a probabilistic assumption.

\begin{longtable}[]{@{}
  >{\raggedright\arraybackslash}p{(\columnwidth - 4\tabcolsep) * \real{0.3333}}
  >{\raggedright\arraybackslash}p{(\columnwidth - 4\tabcolsep) * \real{0.3333}}
  >{\raggedright\arraybackslash}p{(\columnwidth - 4\tabcolsep) * \real{0.3333}}@{}}
\toprule\noalign{}
\begin{minipage}[b]{\linewidth}\raggedright
Text
\end{minipage} & \begin{minipage}[b]{\linewidth}\raggedright
Lean declarations
\end{minipage} & \begin{minipage}[b]{\linewidth}\raggedright
Source
\end{minipage} \\
\midrule\noalign{}
\endhead
\bottomrule\noalign{}
\endlastfoot
Ordinary bands and moving depth &
{\small\nolinkurl{sourceOriginalQ0Scale}},
{\small\nolinkurl{sourceOriginalQActualBands}},
{\small\nolinkurl{sourceOriginalLiteralDepth}} &
\href{../Green64/Library/Correlation/Part065.lean}{\nolinkurl{Correlation/Part065}},
\href{../Green64/Library/Weights/Part097.lean}{\nolinkurl{Weights/Part097}} \\
Every admissible marked band &
{\small\nolinkurl{sourceOriginalPGroupIndex}},
{\small\nolinkurl{source_original_P_all_groups_mem_iff}},
{\small\nolinkurl{sourceOriginalPAllActualBands}} &
\href{../Green64/Library/Weights/SourcePOriginalGroupGeometry.lean}{group
geometry},
\href{../Green64/Library/Weights/Part019.lean}{\nolinkurl{Weights/Part019}},
\href{../Green64/Library/Correlation/Part032.lean}{\nolinkurl{Correlation/Part032}} \\
Factorial coefficient and marked weight &
{\small\nolinkurl{sourceOriginalOrdinaryCoefficient}},
{\small\nolinkurl{sourceOriginalW1}},
{\small\nolinkurl{sourceOriginalA0}} &
\href{../Green64/Library/Analysis/Estimates/TrueCellPrimeFibre.lean}{\nolinkurl{TrueCellPrimeFibre}} \\
Parity projection & {\small\nolinkurl{parityWeight}} &
\href{../Green64/Library/Sieve/Foundations/Part044.lean}{\nolinkurl{Foundations/Part044}} \\
Exact multiplicity sign &
{\small\nolinkurl{b628_parity_liouville_equals_P_phase434}},
{\small\nolinkurl{b628_parity_negative_candidate_exact_target434}} &
\href{../Green64/Library/Weights/ParityLiteralParityTarget.lean}{\nolinkurl{ParityLiteralParityTarget}} \\
Ordinary harmonic lower bounds &
{\small\nolinkurl{b635_actual_moving_Q_twelfth_floor434}},
{\small\nolinkurl{b635_parity_admissible_band_count434}},
{\small\nolinkurl{b635_parity_band_factor_numeric434}},
{\small\nolinkurl{b635_parity_actual_HQ_lower_220434}} &
\href{../Green64/Library/Weights/Part102.lean}{\nolinkurl{Weights/Part102}} \\
Ordinary harmonic upper bound &
{\small\nolinkurl{source_original_Q_actual_band_harmonic_le_four}} &
\href{../Green64/Library/Weights/Part105.lean}{\nolinkurl{Weights/Part105}} \\
Marked harmonic bounds &
{\small\nolinkurl{source_original_P_all_actual_reciprocal_le_four434}},
{\small\nolinkurl{source_original_P_all_reciprocal_floor_eventually}} &
\href{../Green64/Library/Correlation/Part032.lean}{\nolinkurl{Correlation/Part032}},
\href{../Green64/Library/Weights/Part107.lean}{\nolinkurl{Weights/Part107}} \\
Euler normalization and Rankin tail &
{\small\nolinkurl{source_original_P_group_harmonic_tsum_eq_Hg}},
{\small\nolinkurl{source_original_P_actual_polylog_rankin_tail}} &
\href{../Green64/Library/Weights/Part103.lean}{\nolinkurl{Weights/Part103}} \\
Product law and both parity probabilities &
{\small\nolinkurl{source_original_global_actual_P_normalized_harmonic_law}},
{\small\nolinkurl{parity_original_P_global_actual_uniform_positive_parity434}}
&
\href{../Green64/Library/Weights/Part104.lean}{\nolinkurl{Weights/Part104}} \\
Exact localized weight identity and corridor &
{\small\nolinkurl{parity_original_XA_QP_literal_weights_exact_parity434}},
{\small\nolinkurl{source_original_QP_actual_good_rest_corridor}} &
\href{../Green64/Library/Weights/Part106.lean}{\nolinkurl{Weights/Part106}} \\
Strong mass lower bound &
{\small\nolinkurl{parity_original_XA_actual_HQ_HP_floor_eventually_parity434}},
{\small\nolinkurl{parity_original_XA_literal_depth_lower_eventually_parity434}}
&
\href{../Green64/Library/Sieve/Foundations/Part005.lean}{\nolinkurl{Foundations/Part005}} \\
Shared \(X_\rho\ge c_0e^LL^{-221}\) &
{\small\nolinkurl{localizedMass_lower_bound}} &
\href{../Green64/Library/Sieve/Foundations/Part050.lean}{\nolinkurl{Foundations/Part050}} \\
Full logarithmic cells and roughness proxy &
{\small\nolinkurl{sourceD20LogCell}},
{\small\nolinkurl{sourceD20FullProxy}},
{\small\nolinkurl{sourceD20FullDelta}} &
\href{../Green64/Library/Weights/Part012.lean}{\nolinkurl{Weights/Part012}} \\
Prime proxy and character estimate &
{\small\nolinkurl{b627PrimeFullProxy434}},
{\small\nolinkurl{b627_actual_fine_rough_all_character_saving434}} &
\href{../Green64/Library/Analysis/Estimates/PrimeFineProxy.lean}{\nolinkurl{PrimeFineProxy}} \\
Norm of the second proxy &
{\small\nolinkurl{parity_original_actual_second_proxy_beta_half_norm_le_one434}}
&
\href{../Green64/Library/Sieve/Foundations/Part045.lean}{\nolinkurl{Foundations/Part045}} \\
Actual signed source and common analytic precision &
{\small\nolinkurl{parity_eventually_original_signed_quarter_b629_fine_one_slot_Psi_Mellin_source_anylog_common_precision434}}
&
\href{../Green64/Library/Sieve/Foundations/Part050.lean}{\nolinkurl{Foundations/Part050}} \\
Whole rectangles and both replacement orders &
{\small\nolinkurl{parity_eventually_original_both_parity_full_rectangles_proxy_steps434}}
&
\href{../Green64/Library/Sieve/Foundations/Part052.lean}{\nolinkurl{Foundations/Part052}} \\
All finite error terms and normalization &
{\small\nolinkurl{requiredProxyError}},
{\small\nolinkurl{allOrdersProxyError}},
{\small\nolinkurl{parity_original_littleO_exp222_to_Xrho434}} &
\href{../Green64/Library/Sieve/Foundations/Part053.lean}{\nolinkurl{Foundations/Part053}} \\
Singular product &
{\small\nolinkurl{singularSeries_converges_and_ge_one}} &
\href{../Green64/Library/Sieve/Foundations/Part006.lean}{\nolinkurl{Foundations/Part006}} \\
Slot allocation, diagonal estimate, Fourier tails & Detailed source map
in
\protect\hyperlink{bilinear-source-correspondence}{Section~\ref{bilinear-source-correspondence}}
&
\href{../Green64/Library/Sieve/Foundations/Part050.lean}{\nolinkurl{Foundations/Part050}} \\
Alternating resampling, physical moments, and comparison subtraction &
Derivation in
\protect\hyperlink{marked-correlation}{Section~\ref{marked-correlation}};
source map in
\protect\hyperlink{marked-correlation-correspondence-with-lean}{Section~\ref{marked-correlation-correspondence-with-lean}}
&
\href{../Green64/Library/Weights/Part092.lean}{\nolinkurl{Weights/Part092}},
\href{../Green64/Library/Sieve/Foundations/Part008.lean}{\nolinkurl{Foundations/Part008}} \\
Negative-parity congruence estimate &
{\small\nolinkurl{parity_original_negative_congruence_D230_bound434}} &
\href{../Green64/Library/Sieve/Foundations/Part053.lean}{\nolinkurl{Foundations/Part053}} \\
Character decomposition, deleted-slot convolution, and conductor
estimates & Detailed source map in
\protect\hyperlink{congruences-source-correspondence}{Section~\ref{congruences-source-correspondence}}
&
\href{../Green64/Library/Weights/Part107.lean}{\nolinkurl{Weights/Part107}},
\href{../Green64/Library/Weights/Part108.lean}{\nolinkurl{Part108}} \\
Assignment masks, CRT, and the finite rectangle sieve & Detailed source
map in
\protect\hyperlink{rectangles-source-correspondence}{Section~\ref{rectangles-source-correspondence}}
&
\href{../Green64/Library/Sieve/Foundations/Part056.lean}{\nolinkurl{Foundations/Part056}},
\href{../Green64/Library/Sieve/Foundations/Part057.lean}{\nolinkurl{Part057}} \\
Rough-number density & {\small\nolinkurl{sourceSimplexTerm}},
{\small\nolinkurl{sourceDensity}} &
\href{../Green64/Library/Analysis/Estimates/Part003.lean}{\nolinkurl{Analysis/Estimates/Part003}} \\
Remove unbalanced composites &
{\small\nolinkurl{parity_original_exists_small_b_unbalanced_margin434}},
{\small\nolinkurl{negative_prime_or_composite_mass_lower_bound}} &
\href{../Green64/Library/Sieve/Foundations/Part055.lean}{\nolinkurl{Foundations/Part055}} \\
\end{longtable}

\Needspace{6\baselineskip}

\hypertarget{lean-guide-the-analytic-library}{%
\subsection{The analytic
library}\label{lean-guide-the-analytic-library}}

{\small\nolinkurl{Library/Weights}} develops labelled prime assignments,
multiplicities, harmonic normalization, and the finite combinatorial
operators used by the correlation argument.
{\small\nolinkurl{Library/Correlation}} contains their signed
correlations and reindexing identities.
{\small\nolinkurl{Library/Sieve/Foundations}} connects those identities
to the sieve and the finite factor sums.
{\small\nolinkurl{Library/Analysis}}, {\small\nolinkurl{Library/OAI}},
{\small\nolinkurl{Library/ChenTheorem}}, and the other retained source
directories provide the analytic and arithmetic estimates used in those
proofs. See \href{../THIRD_PARTY.md}{THIRD\_PARTY.md} for their origins.

The \protect\hyperlink{bilinear}{bilinear chapter} follows the exact
Cauchy square, the diagonal divisor bound, and the integration over the
Fourier ball and its tail. The
\protect\hyperlink{marked-correlation}{marked-correlation chapter}
supplies the cancellation used inside that ball, including the
separately bounded comparison term. The final analytic parameter
construction discharges the hypotheses of these lower-level estimates.

The {\small\nolinkurl{PartNNN}} files contain supporting results; the
named files in {\small\nolinkurl{Green64/Sieve}} expose the final
argument.

\Needspace{6\baselineskip}

\hypertarget{lean-guide-checking-what-is-proved}{%
\subsection{Checking what is
proved}\label{lean-guide-checking-what-is-proved}}

A successful build alone is not the axiom audit. Run
{\small\nolinkurl{scripts/verify.py}} and inspect its
{\small\nolinkurl{VALIDATION.json}} receipt for the public theorem
names, axioms, source hashes, dependency revisions, and command exit
codes. The audit checks the infinite set proposition itself, together
with the stronger unbounded-witness theorem. It does not substitute a
conditional statement or a wrapper with an assumed analytic conclusion.

\clearpage
\hypertarget{source-correspondence}{%
\section{Further source correspondence}\label{source-correspondence}}

\Needspace{10\baselineskip}

\hypertarget{correlation-source-correspondence}{%
\subsection{Replacing prime and roughness
tests}\label{correlation-source-correspondence}}

\begin{longtable}[]{@{}
  >{\raggedright\arraybackslash}p{(\columnwidth - 2\tabcolsep) * \real{0.5000}}
  >{\raggedright\arraybackslash}p{(\columnwidth - 2\tabcolsep) * \real{0.5000}}@{}}
\toprule\noalign{}
\begin{minipage}[b]{\linewidth}\raggedright
Text
\end{minipage} & \begin{minipage}[b]{\linewidth}\raggedright
Definitions and proved identities
\end{minipage} \\
\midrule\noalign{}
\endhead
\bottomrule\noalign{}
\endlastfoot
(C1)--(C3), complete-cell discrepancies &
\href{../Green64/Library/Weights/Part012.lean}{\nolinkurl{Weights/Part012}}:
{\small\nolinkurl{sourceD20Rough}},
{\small\nolinkurl{sourceD20LogCell}},
{\small\nolinkurl{sourceD20FullProxy}},
{\small\nolinkurl{sourceD20FullDelta}},
{\small\nolinkurl{sourceD20_full_delta_zero}};
\href{../Green64/Library/Analysis/Estimates/Part005.lean}{\nolinkurl{Estimates/Part005}}:
{\small\nolinkurl{sourceD20_principal_interval_bound}},
{\small\nolinkurl{sourceD20_mesh_scalar}},
{\small\nolinkurl{sourceD21_full_nonprincipal_saving}};
\href{../Green64/Library/Analysis/Estimates/PrimeFineProxy.lean}{\nolinkurl{PrimeFineProxy}}:
{\small\nolinkurl{b627PrimeFullProxy434}},
{\small\nolinkurl{b627_actual_fine_rough_all_character_saving434}},
{\small\nolinkurl{b627_prime_full_delta_eq_actual_rough434}} \\
(C5)--(C6), common bilinear precision &
\href{../Green64/Library/Sieve/Foundations/Part050.lean}{\nolinkurl{Foundations/Part050}}:
{\small\nolinkurl{parity_original_TypeII_target_K0_precision_compatible434}},
{\small\nolinkurl{parity_eventually_original_signed_quarter_b629_fine_one_slot_Psi_Mellin_source_anylog_common_precision434}};
the exponents are defined in
\href{../Green64/Library/Sieve/Foundations/Part013.lean}{\nolinkurl{Foundations/Part013}}:
{\small\nolinkurl{parityOriginalW1PolylogExponent434}},
\href{../Green64/Library/Sieve/Foundations/Part011.lean}{\nolinkurl{Part011}}:
{\small\nolinkurl{parityOriginalSelectedECoefficientPolylogExponent434}},
and
\href{../Green64/Library/Sieve/Foundations/Part010.lean}{\nolinkurl{Part010}}:
{\small\nolinkurl{parityOriginalB590FixedFourierPolynomialExponent434}} \\
(C7)--(C8), two replacements and dyadic partition &
\href{../Green64/Library/Sieve/Foundations/Part045.lean}{\nolinkurl{Foundations/Part045}}:
{\small\nolinkurl{parity_original_actual_fullcell_proxy_beta_norm_le_two434}},
{\small\nolinkurl{parity_original_actual_second_proxy_beta_half_norm_le_one434}};
\href{../Green64/Library/Sieve/Foundations/Part050.lean}{\nolinkurl{Foundations/Part050}}:
{\small\nolinkurl{parity_original_parity_positive_Icc_eq_quarter_nested_sources434}} \\
(C9)--(C12), total error and normalization &
\href{../Green64/Library/Sieve/Foundations/Part053.lean}{\nolinkurl{Foundations/Part053}}:
{\small\nolinkurl{requiredProxyError}},
{\small\nolinkurl{allOrdersProxyError}},
{\small\nolinkurl{parity_eventually_original_sieve_all_proxy_errors434}},
{\small\nolinkurl{parity_original_littleO_exp222_to_Xrho434}};
\href{../Green64/Sieve/Parameters.lean}{\nolinkurl{Parameters}}:
{\small\nolinkurl{proxy_parameters_exist}} \\
\end{longtable}

\Needspace{10\baselineskip}

\hypertarget{bilinear-source-correspondence}{%
\subsection{The bilinear
reduction}\label{bilinear-source-correspondence}}

\begin{longtable}[]{@{}
  >{\raggedright\arraybackslash}p{(\columnwidth - 2\tabcolsep) * \real{0.5000}}
  >{\raggedright\arraybackslash}p{(\columnwidth - 2\tabcolsep) * \real{0.5000}}@{}}
\toprule\noalign{}
\begin{minipage}[b]{\linewidth}\raggedright
Argument
\end{minipage} & \begin{minipage}[b]{\linewidth}\raggedright
Lean source
\end{minipage} \\
\midrule\noalign{}
\endhead
\bottomrule\noalign{}
\endlastfoot
Common-precision rectangle statement &
\href{../Green64/Library/Sieve/Foundations/Part050.lean}{\nolinkurl{Foundations/Part050}}:
{\small\nolinkurl{parity_eventually_original_signed_quarter_b629_fine_one_slot_Psi_Mellin_source_anylog_common_precision434}} \\
Smooth allocation functions &
\href{../Green64/Library/Weights/Part003.lean}{\nolinkurl{Weights/Part003}}:
{\small\nolinkurl{sourceBufferedGateReal}} \\
Slot ordering and frequency scales &
\href{../Green64/Library/Correlation/Part065.lean}{\nolinkurl{Correlation/Part065}}:
{\small\nolinkurl{sourceOrderedLabels}},
{\small\nolinkurl{sourceActualSortedFrequencyScales}} \\
Prefix bounds, partition, and the exponent \(88\) &
\href{../Green64/Library/Weights/Part117.lean}{\nolinkurl{Weights/Part117}} \\
Fourier phases and endpoint sums &
\href{../Green64/Library/Weights/Part020.lean}{\nolinkurl{Weights/Part020}},
\href{../Green64/Library/Sieve/Foundations/Part009.lean}{\nolinkurl{Foundations/Part009}},
\href{../Green64/Library/Sieve/Foundations/Part011.lean}{\nolinkurl{Part011}} \\
Intrinsic \(H\) and \(a_v\) &
\href{../Green64/Library/Sieve/Foundations/Part017.lean}{\nolinkurl{Foundations/Part017}}:
{\small\nolinkurl{parityOriginalIntrinsicH434}},
{\small\nolinkurl{parityOriginalIntrinsicSignedAV434}} \\
Windowed \(a_e\) and source comparison &
\href{../Green64/Library/Sieve/Foundations/Part018.lean}{\nolinkurl{Foundations/Part018}}:
{\small\nolinkurl{parityOriginalHistoricalWindowAE434}} \\
Finite \(\mathcal B'\) and divisor-moment diagonal &
\href{../Green64/Library/Sieve/Foundations/Part032.lean}{\nolinkurl{Foundations/Part032}}:
{\small\nolinkurl{parityOriginalHistoricalUngatedCauchyAtDeltaNoUCap434}} \\
Smooth square and exact off-diagonal &
\href{../Green64/Library/Sieve/Foundations/Part022.lean}{\nolinkurl{Foundations/Part022}},
\href{../Green64/Library/Sieve/Foundations/Part025.lean}{\nolinkurl{Part025}},
\href{../Green64/Library/Sieve/Foundations/Part028.lean}{\nolinkurl{Part028}} \\
Shift coordinates and residue mask &
\href{../Green64/Library/Sieve/Foundations/Part019.lean}{\nolinkurl{Foundations/Part019}},
\href{../Green64/Library/Sieve/Foundations/Part029.lean}{\nolinkurl{Part029}},
\href{../Green64/Library/Weights/Part005.lean}{\nolinkurl{Weights/Part005}} \\
Exact two-endpoint cofactor identity &
\href{../Green64/Library/Weights/Part102.lean}{\nolinkurl{Weights/Part102}}:
{\small\nolinkurl{b633_parity_full_physical_pair434}} \\
Physical correlation saving (B18) &
\href{../Green64/Library/Sieve/Foundations/Part008.lean}{\nolinkurl{Foundations/Part008}}:
{\small\nolinkurl{parity_eventually_actual_original_Parity417_physical_cofactor_saving_common_precision434}} \\
Actual Fourier moments (B19a) and the extra precision in (B19b) &
\href{../Green64/Library/Sieve/Foundations/Part026.lean}{\nolinkurl{Foundations/Part026}}:
{\small\nolinkurl{parity_original_historical_full_odd_3d_fourier_moment_uniform_polylog434}};
\href{../Green64/Library/Sieve/Foundations/Part032.lean}{\nolinkurl{Part032}}:
{\small\nolinkurl{parity_original_historical_box_source_BH_uniform_polylog434}};
\href{../Green64/Library/Sieve/Foundations/Part037.lean}{\nolinkurl{Part037}}:
{\small\nolinkurl{parity_original_actual_capfree_Bprime_M0_uniform_polylog434}},
{\small\nolinkurl{parity_original_actual_capfree_Bprime_primary_inverse_log_from_M0434}} \\
Fourier ball, tails, and common precision &
\href{../Green64/Library/Sieve/Foundations/Part046.lean}{\nolinkurl{Foundations/Part046}},
\href{../Green64/Library/Sieve/Foundations/Part049.lean}{\nolinkurl{Part049}},
\href{../Green64/Library/Sieve/Foundations/Part050.lean}{\nolinkurl{Part050}} \\
\end{longtable}

\Needspace{10\baselineskip}

\hypertarget{marked-correlation-correspondence-with-lean}{%
\subsection{Cancellation in the marked
correlation}\label{marked-correlation-correspondence-with-lean}}

The transfer formulas have the following more detailed correspondence.
Here {\small\nolinkurl{k}} in a local operator declaration means
\(s=2k\) in (B14), {\small\nolinkurl{rho}} in the memory declarations
means \(\rho_0\) (not parity), and {\small\nolinkurl{mu}} in a
physical-state declaration is the tilted law (T2).

\begin{longtable}[]{@{}
  >{\raggedright\arraybackslash}p{(\columnwidth - 2\tabcolsep) * \real{0.5000}}
  >{\raggedright\arraybackslash}p{(\columnwidth - 2\tabcolsep) * \real{0.5000}}@{}}
\toprule\noalign{}
\begin{minipage}[b]{\linewidth}\raggedright
Text and mathematical operation
\end{minipage} & \begin{minipage}[b]{\linewidth}\raggedright
Definitions and proved identities
\end{minipage} \\
\midrule\noalign{}
\endhead
\bottomrule\noalign{}
\endlastfoot
(M0), arbitrary joint endpoint tests and permutation averaging &
\href{../Green64/Library/Weights/Part090.lean}{\nolinkurl{Weights/Part090}}:
{\small\nolinkurl{sourceOriginalUnfilteredIdealEndpointMean434}};
\href{../Green64/Library/Weights/Part091.lean}{\nolinkurl{Part091}}:
{\small\nolinkurl{source_original_unfiltered_ideal_ENDPOINT_PAIR_MEAN434}} \\
(M0a), the phase and the asymmetric splice in (T7) &
\href{../Green64/Library/Weights/Part024.lean}{\nolinkurl{Weights/Part024}}:
{\small\nolinkurl{originalEdgeScalar}};
\href{../Green64/Library/Weights/Part091.lean}{\nolinkurl{Part091}}:
{\small\nolinkurl{sourceOriginalUnfilteredIdealHeadPhase434}},
{\small\nolinkurl{source_original_unfiltered_ideal_raw_splice_phase434}} \\
(M4a), product matching and the resampling reference &
\href{../Green64/Library/Weights/Part091.lean}{\nolinkurl{Weights/Part091}}:
{\small\nolinkurl{source_original_unfiltered_ideal_raw_splice_phase434}};
\href{../Green64/Library/Weights/Part092.lean}{\nolinkurl{Part092}}:
{\small\nolinkurl{source_original_unfiltered_ideal_RAW_FORM_RESIDUAL434}}
(the displayed matching kernel explains the finite product-law
substitution in these identities) \\
(T1)--(T3), actual state spaces and measures &
\href{../Green64/Library/Analysis/Estimates/Part003.lean}{\nolinkurl{Estimates/Part003}}:
{\small\nolinkurl{SmallState}}, {\small\nolinkurl{stateMass}};
\href{../Green64/Library/Weights/Part002.lean}{\nolinkurl{Weights/Part002}}:
{\small\nolinkurl{FockState}}, {\small\nolinkurl{fockMass}},
{\small\nolinkurl{BigState}}, {\small\nolinkurl{bigStateMass}};
\href{../Green64/Library/Correlation/Part006.lean}{\nolinkurl{Correlation/Part006}}:
{\small\nolinkurl{physicalSourceBaseline}},
{\small\nolinkurl{physicalSourceMuTilde}} \\
(T4), quotient map, mass and norm &
\href{../Green64/Library/Correlation/Part005.lean}{\nolinkurl{Correlation/Part005}}:
{\small\nolinkurl{physicalFockProjection}},
{\small\nolinkurl{physicalOrbitMass}};
\href{../Green64/Library/Correlation/Part006.lean}{\nolinkurl{Part006}}:
{\small\nolinkurl{fullPhysicalOrbitAverage}};
\href{../Green64/Library/Weights/Part093.lean}{\nolinkurl{Weights/Part093}}:
{\small\nolinkurl{sourceOriginalUnfilteredIdealOrbitEnergy434}} \\
(T5), legal transitions, fresh births and the two OLD gates &
\href{../Green64/Library/Weights/Part002.lean}{\nolinkurl{Weights/Part002}}:
{\small\nolinkurl{PhysicalEdgeLegal}},
{\small\nolinkurl{physicalBigBranchCoefficient}};
\href{../Green64/Library/Weights/Part024.lean}{\nolinkurl{Part024}}:
{\small\nolinkurl{sourceOrientedPhysicalCoefficient}};
\href{../Green64/Library/Weights/Part093.lean}{\nolinkurl{Part093}}:
{\small\nolinkurl{sourceOriginalUnfilteredEdgeLiteralKeep434}},
{\small\nolinkurl{sourceOriginalUnfilteredEdgeSeenErasedPhase434}} \\
(T6)--(T7), frozen form and raw/comparison/ban identity &
\href{../Green64/Library/Correlation/B26FrozenCleanPhysicalFourier.lean}{\nolinkurl{B26FrozenCleanPhysicalFourier}}:
{\small\nolinkurl{sourceOriginalB26FrozenFreshLaw434}},
{\small\nolinkurl{sourceOriginalB26FrozenBigWeight434}};
\href{../Green64/Library/Weights/Part092.lean}{\nolinkurl{Weights/Part092}}:
{\small\nolinkurl{source_original_unfiltered_ideal_FULL_FROZEN_PAIR434}},
{\small\nolinkurl{eventually_source_original_unfiltered_ideal_FULL_UNIT434}} \\
(T8)--(T12), common density, fresh collisions and exact three-term
identity &
\href{../Green64/Library/Weights/Part090.lean}{\nolinkurl{Weights/Part090}}:
{\small\nolinkurl{source_original_unfiltered_ideal_COMMON_ENERGY434}},
{\small\nolinkurl{source_original_unfiltered_ideal_PHYSICAL_COMMON_DENSITY434}};
\href{../Green64/Library/Weights/Part092.lean}{\nolinkurl{Part092}}:
{\small\nolinkurl{source_original_physical_FROZEN_BAD_FAMILY434}};
\href{../Green64/Library/Weights/Part094.lean}{\nolinkurl{Part094}}:
{\small\nolinkurl{source_original_physical_EXACT_THREE_TERM434}},
{\small\nolinkurl{eventually_source_original_physical_FULL_TILTED_NODUP_BIG434}} \\
(T13)--(T15), SMALL fibres and Fourier transform of the actual
translation &
\href{../Green64/Library/Weights/B26CleanSmallCoherentPairing.lean}{\nolinkurl{B26CleanSmallCoherentPairing}}:
{\small\nolinkurl{sourceOriginalB26SmallLift434}},
{\small\nolinkurl{source_original_B26_MOVING_SMALL_CLEAN_FRAME_PAIR434}},
{\small\nolinkurl{source_original_B26_ORIGINAL_CLEAN_SMALL_FOURIER434}} \\
(T16)--(T17), both endpoint energies and orbit transfer &
\href{../Green64/Library/Analysis/Estimates/Part013.lean}{\nolinkurl{Estimates/Part013}}:
{\small\nolinkurl{source_original_B26_SMALL_OLD_TRANSFER434}};
\href{../Green64/Library/Weights/Part091.lean}{\nolinkurl{Weights/Part091}}:
{\small\nolinkurl{source_original_unfiltered_ideal_LITERAL_SECTOR_TRANSFER434}};
\href{../Green64/Library/Correlation/Part063.lean}{\nolinkurl{Correlation/Part063}}:
{\small\nolinkurl{eventually_source_original_physical_FULL_SEEN_ERASED_CLEAN434}} \\
(M11)--(M13), exceptional transitions; (T18)--(T19), full local estimate
&
\href{../Green64/Library/Correlation/Part062.lean}{\nolinkurl{Correlation/Part062}}:
{\small\nolinkurl{source_original_edge_LITERAL_DIRTY_FINITE434}};
\href{../Green64/Library/Weights/Part094.lean}{\nolinkurl{Weights/Part094}}:
{\small\nolinkurl{eventually_source_original_edge_SOURCE_PHASED_BOUNDS434}};
\href{../Green64/Library/Correlation/Part065.lean}{\nolinkurl{Correlation/Part065}}:
{\small\nolinkurl{eventually_source_original_edge_FULL_NEAR_CONTRACTION434}},
{\small\nolinkurl{eventually_source_original_root_SOURCE_LOCAL_BANKS434}} \\
(T20), background births, signs and factorials &
\href{../Green64/Library/Weights/Part002.lean}{\nolinkurl{Weights/Part002}}:
{\small\nolinkurl{sourceGhostTerm}};
\href{../Green64/Library/Correlation/Part005.lean}{\nolinkurl{Correlation/Part005}}:
{\small\nolinkurl{physicalVisitGKernel}};
\href{../Green64/Library/Weights/Part094.lean}{\nolinkurl{Weights/Part094}}:
{\small\nolinkurl{source_original_evolution_FINAL_G_ENERGY434}} \\
(T21)--(T23), all moment constants &
\href{../Green64/Library/Correlation/Part005.lean}{\nolinkurl{Correlation/Part005}}:
{\small\nolinkurl{memorySmallRowConstant}},
{\small\nolinkurl{memoryPromotionDampingConstant}};
\href{../Green64/Library/Correlation/Part062.lean}{\nolinkurl{Part062}}:
{\small\nolinkurl{sourceOriginalRootPinnedRowBudget434}};
\href{../Green64/Library/Weights/Part072.lean}{\nolinkurl{Weights/Part072}}:
{\small\nolinkurl{sourceOriginalActualSourceSignedInitialSmallMass434}};
\href{../Green64/Library/Weights/Part093.lean}{\nolinkurl{Part093}}:
{\small\nolinkurl{sourceOriginalUnfilteredEdgeExponent434}};
\href{../Green64/Library/Weights/Part094.lean}{\nolinkurl{Part094}}:
{\small\nolinkurl{source_original_global_root_PREFACTOR434}} \\
(T24)--(T25), prime lifespans and the exact residue model &
\href{../Green64/Library/Correlation/Part009.lean}{\nolinkurl{Correlation/Part009}}:
{\small\nolinkurl{memoryAbsentTerm}},
{\small\nolinkurl{memoryRepresentedTerm}};
\href{../Green64/Library/Correlation/Part010.lean}{\nolinkurl{Part010}}:
{\small\nolinkurl{memory_frozen_prime_all_lifespans}},
{\small\nolinkurl{memory_frozen_prime_supported_normal_form}};
\href{../Green64/Library/Correlation/Part025.lean}{\nolinkurl{Part025}}:
{\small\nolinkurl{memory_source_root_actual_prime_product_exact}};
\href{../Green64/Library/Correlation/Part027.lean}{\nolinkurl{Part027}}:
{\small\nolinkurl{memory_source_original_physical_root_eq_frozen_Haar}} \\
(M16a)--(M19), same-integer densities and signed transfer &
\href{../Green64/Library/Weights/WeightedActualPaths.lean}{\nolinkurl{WeightedActualPaths}}:
{\small\nolinkurl{b645ActualMarkedPathModel434}},
{\small\nolinkurl{b645_eventually_actual_weighted_paths_CRT434}};
\href{../Green64/Library/Weights/Part073.lean}{\nolinkurl{Weights/Part073}}:
{\small\nolinkurl{b645_eventually_actual_parity_physical_root_transfer434}},
{\small\nolinkurl{b654_actual_parity_moment_eq_return434}};
\href{../Green64/Library/Correlation/Part042.lean}{\nolinkurl{Correlation/Part042}}:
{\small\nolinkurl{b650_physical_root_moment_array434}};
\href{../Green64/Library/Weights/Part048.lean}{\nolinkurl{Weights/Part048}}:
{\small\nolinkurl{b654ParityIntervalWindow434}},
{\small\nolinkurl{b654ParitySeedPositions434}} \\
(M20), spectral interpolation and position blocks &
\href{../Green64/Library/Weights/Part005.lean}{\nolinkurl{Weights/Part005}}:
{\small\nolinkurl{b652_mark_independent_position_endpoint434}} \\
(M20a), the full path window and fixed endpoint bounds &
\href{../Green64/Library/Weights/ParityStepCap.lean}{\nolinkurl{ParityStepCap}}:
{\small\nolinkurl{b651ParityDConstant434}},
{\small\nolinkurl{b651ParityMaxStep434}};
\href{../Green64/Library/Weights/ParityLocalizationWindow.lean}{\nolinkurl{ParityLocalizationWindow}}:
{\small\nolinkurl{b651ParityWindowSize434}};
\href{../Green64/Library/Weights/Part109.lean}{\nolinkurl{Weights/Part109}}:
{\small\nolinkurl{b660_eventually_actual_parity_dyadic_saving434}} (uses
norm exponent 9 and supremum exponent 5) \\
(M23a)--(M23b), rank budgets and the coupled auxiliary sum &
\href{../Green64/Library/Correlation/Part072.lean}{\nolinkurl{Correlation/Part072}}:
{\small\nolinkurl{parity_canonical_rank_card434}},
{\small\nolinkurl{parity_eventually_actual_rank_family_budget434}};
\href{../Green64/Library/Weights/Part112.lean}{\nolinkurl{Weights/Part112}}:
{\small\nolinkurl{parity_actual_rank_aux_coupled_reindex434}};
\href{../Green64/Library/Weights/Part111.lean}{\nolinkurl{Weights/Part111}}:
{\small\nolinkurl{parity_fixed_rank_cutoff_pair_budget434}} \\
(M23c)--(M23f), exact endpoints, shared restoration and the
free-auxiliary error &
\href{../Green64/Library/Weights/Part046.lean}{\nolinkurl{Weights/Part046}}:
{\small\nolinkurl{parityFirstRankMark434}},
{\small\nolinkurl{paritySecondRankMark434}};
\href{../Green64/Library/Correlation/Part072.lean}{\nolinkurl{Correlation/Part072}}:
{\small\nolinkurl{parityActualRankEndpoint434}},
{\small\nolinkurl{parity_eventually_comparison_weighted_restored434}};
\href{../Green64/Library/Weights/WeightedStrippedRank.lean}{\nolinkurl{WeightedStrippedRank}}:
{\small\nolinkurl{parity_actual_weighted_stripped_rank_difference434}};
\href{../Green64/Library/Weights/Part113.lean}{\nolinkurl{Weights/Part113}}:
{\small\nolinkurl{parity_eventually_actual_free_collision_pair_bound434}};
\href{../Green64/Library/Weights/FreeAuxCollisionJoin.lean}{\nolinkurl{FreeAuxCollisionJoin}}:
{\small\nolinkurl{parity_actual_stripped_family_rank_main_difference434}} \\
(M23g)--(M23i), Mellin-circle pairing with the full auxiliary product &
\href{../Green64/Library/Analysis/Estimates/RankFourier.lean}{\nolinkurl{RankFourier}}:
{\small\nolinkurl{parity_actual_rank_pair_fourier434}};
\href{../Green64/Library/Analysis/Estimates/Part016.lean}{\nolinkurl{Estimates/Part016}}:
{\small\nolinkurl{parity_actual_weighted_cell_rank_main_mellin434}},
{\small\nolinkurl{parity_actual_cell_rank_family_circle434}},
{\small\nolinkurl{parity_actual_cell_coupled_kernel_norm434}} \\
\end{longtable}

The chapter follows the finite constructions and their estimates in the
following files. These references identify the proof being explained;
the mathematical argument above uses the displayed identities and
inequalities.

\begin{longtable}[]{@{}
  >{\raggedright\arraybackslash}p{(\columnwidth - 2\tabcolsep) * \real{0.5000}}
  >{\raggedright\arraybackslash}p{(\columnwidth - 2\tabcolsep) * \real{0.5000}}@{}}
\toprule\noalign{}
\begin{minipage}[b]{\linewidth}\raggedright
Step
\end{minipage} & \begin{minipage}[b]{\linewidth}\raggedright
Included Lean source
\end{minipage} \\
\midrule\noalign{}
\endhead
\bottomrule\noalign{}
\endlastfoot
Hoeffding components and (M1) &
\href{../Green64/Library/Weights/Part017.lean}{\nolinkurl{Weights/Part017}},
{\small\nolinkurl{source_original_hoeffding_physical_all_degree_saving434}} \\
Actual BIG block sizes &
\href{../Green64/Library/Weights/Part100.lean}{\nolinkurl{Weights/Part100}} \\
Alternating identity and its squared bound &
\href{../Green64/Library/Weights/Part092.lean}{\nolinkurl{Weights/Part092}},
{\small\nolinkurl{source_original_unfiltered_ideal_RAW_FORM_RESIDUAL434}} \\
Character Gram matrices &
\href{../Green64/Library/Correlation/Part037.lean}{\nolinkurl{Correlation/Part037}} \\
Global kernel norm &
\href{../Green64/Library/Weights/ActualGlobalKernelL2.lean}{\nolinkurl{ActualGlobalKernelL2}} \\
Rational phase, product fibres, and minor arcs &
\href{../Green64/Library/Correlation/Part039.lean}{\nolinkurl{Correlation/Part039}},
\href{../Green64/Library/Weights/Part047.lean}{\nolinkurl{Weights/Part047}},
\href{../Green64/Library/Weights/Part053.lean}{\nolinkurl{Part053}} \\
Arithmetic mesh choices &
\href{../Green64/Library/Correlation/Part041.lean}{\nolinkurl{Correlation/Part041}} \\
SMALL omission uniqueness &
\href{../Green64/Library/Correlation/Part043.lean}{\nolinkurl{Correlation/Part043}} \\
OLD, density, and collision corrections &
\href{../Green64/Library/Correlation/Part062.lean}{\nolinkurl{Correlation/Part062}},
\href{../Green64/Library/Weights/Part094.lean}{\nolinkurl{Weights/Part094}} \\
Forest rank split and restoration &
\href{../Green64/Library/Correlation/GlobalRootScalar.lean}{\nolinkurl{GlobalRootScalar}},
\href{../Green64/Library/Correlation/Part065.lean}{\nolinkurl{Correlation/Part065}} \\
Removal of the OLD cutoff &
\href{../Green64/Library/Weights/Part064.lean}{\nolinkurl{Weights/Part064}},
\href{../Green64/Library/Weights/Part090.lean}{\nolinkurl{Part090}},
\href{../Green64/Library/Correlation/Part050.lean}{\nolinkurl{Correlation/Part050}},
\href{../Green64/Library/Correlation/Part052.lean}{\nolinkurl{Part052}} \\
Bonferroni--CRT interval estimate &
\href{../Green64/Library/Correlation/Part035.lean}{\nolinkurl{Correlation/Part035}} \\
Signed path transfer &
\href{../Green64/Library/Weights/WeightedActualPaths.lean}{\nolinkurl{WeightedActualPaths}},
\href{../Green64/Library/Sieve/Foundations/Part005.lean}{\nolinkurl{Foundations/Part005}} \\
Physical states, matrix entries, and returning paths &
\href{../Green64/Library/Correlation/Part021.lean}{\nolinkurl{Correlation/Part021}},
\href{../Green64/Library/Correlation/Part041.lean}{\nolinkurl{Part041}},
\href{../Green64/Library/Correlation/Part042.lean}{\nolinkurl{Part042}},
\href{../Green64/Library/Weights/Part073.lean}{\nolinkurl{Weights/Part073}} \\
Spectral endpoint estimate &
\href{../Green64/Library/Analysis/Estimates/Part015.lean}{\nolinkurl{Analysis/Estimates/Part015}},
{\small\nolinkurl{b655_eventually_actual_parity_endpoint_saving434}} \\
Literal core and ordered-mark weights &
\href{../Green64/Library/Weights/Part037.lean}{\nolinkurl{Weights/Part037}},
\href{../Green64/Library/Weights/Part038.lean}{\nolinkurl{Part038}},
\href{../Green64/Library/Weights/Part044.lean}{\nolinkurl{Part044}} \\
Marked arc estimate (M25)--(M26) &
\href{../Green64/Library/Weights/Part045.lean}{\nolinkurl{Weights/Part045}},
{\small\nolinkurl{eventually_source_original_arc_conclusion434}} \\
Comparison main term and its subtraction &
\href{../Green64/Library/Sieve/Foundations/Part007.lean}{\nolinkurl{Foundations/Part007}},
\href{../Green64/Library/Sieve/Foundations/Part008.lean}{\nolinkurl{Part008}} \\
Cap and shared-divisor errors &
\href{../Green64/Library/Weights/Part111.lean}{\nolinkurl{Weights/Part111}},
\href{../Green64/Library/Weights/Part115.lean}{\nolinkurl{Part115}} \\
\end{longtable}

\Needspace{10\baselineskip}

\hypertarget{physical-lift-source-correspondence}{%
\subsection{From one mark to the physical
matrix}\label{physical-lift-source-correspondence}}

The raw slice is {\small\nolinkurl{parityRawPhysicalSlice434}} in
\href{../Green64/Library/Correlation/Part072.lean}{\nolinkurl{Correlation/Part072}}.
The clean identity and (L4) are
{\small\nolinkurl{parity_raw_clean_slice_exact434}} and
{\small\nolinkurl{parity_raw_shared_restoration_bound434}} in
\href{../Green64/Library/Weights/RawSharedRestoration.lean}{\nolinkurl{RawSharedRestoration}}.
The carrier reindexing is in
\href{../Green64/Library/Weights/RawSharedCarrier.lean}{\nolinkurl{RawSharedCarrier}},
and the arithmetic cancellation and original-source identification are
in
\href{../Green64/Library/Weights/Part115.lean}{\nolinkurl{Weights/Part115}}.
The endpoint vectors, removed-product fourth moment, and exact Fourier
pairing are in
\href{../Green64/Library/Weights/Part109.lean}{\nolinkurl{Weights/Part109}}.

\Needspace{10\baselineskip}

\hypertarget{forest-moments-source-correspondence}{%
\subsection{The forest expansion of the history
sum}\label{forest-moments-source-correspondence}}

The list identity is
{\small\nolinkurl{source_original_root_LIST_FOREST_EXPANSION434}} in
\href{../Green64/Library/Weights/Part093.lean}{\nolinkurl{Weights/Part093}}.
The actual linear functional, exact forest sum, sign-field average,
owner count and energy iteration are in
\href{../Green64/Library/Correlation/Part064.lean}{\nolinkurl{Correlation/Part064}},
especially {\small\nolinkurl{sourceOriginalRootFunctionalRoot434}},
{\small\nolinkurl{source_original_root_ORIGINAL_FOREST_SUM434}} and
{\small\nolinkurl{source_original_root_FOREST_FINITE_BOUND434}}. The
pinned row estimate is
{\small\nolinkurl{source_original_forest_E_PINNED_FULL_ROW434}} in
\href{../Green64/Library/Weights/Part090.lean}{\nolinkurl{Weights/Part090}}.
The full-rank suffix induction and the initial distinctness case are in
\href{../Green64/Library/Correlation/Part064.lean}{\nolinkurl{Correlation/Part064}};
the resulting bound (F11) is
{\small\nolinkurl{eventually_source_original_global_root_SOURCE_HIGH434}}
in
\href{../Green64/Library/Correlation/Part065.lean}{\nolinkurl{Correlation/Part065}}.

\Needspace{10\baselineskip}

\hypertarget{local-energy-source-correspondence}{%
\subsection{From complete-cell discrepancy to additive arc
energy}\label{local-energy-source-correspondence}}

The local-energy transfer is
{\small\nolinkurl{source_linear_to_log_arc434}} in
\href{../Green64/Library/Weights/Part012.lean}{\nolinkurl{Weights/Part012}},
using {\small\nolinkurl{source_physical_collar_replacement_error434}} in
\href{../Green64/Library/Analysis/PrimeNumberTheorem/Part003.lean}{\nolinkurl{Analysis/PrimeNumberTheorem/Part003}}.
The weighted Schur estimate is
{\small\nolinkurl{source_actual_physical_normalized_L2_schur434}} in
\href{../Green64/Library/Weights/Part023.lean}{\nolinkurl{Weights/Part023}}.
The arithmetic identities and (E13) are in
\href{../Green64/Library/Analysis/Estimates/Part009.lean}{\nolinkurl{Analysis/Estimates/Part009}},
especially
{\small\nolinkurl{source_original_global_physical_mellin_exact434}} and
{\small\nolinkurl{eventually_source_original_global_physical_from_literal_cells434}}.
The direct low-frequency estimate is in
\href{../Green64/Library/Analysis/Estimates/SourceLowFrequencyFourFactorIntegral.lean}{\nolinkurl{SourceLowFrequencyFourFactorIntegral}}.
The three-region argument and retained Fourier frequencies are in
\href{../Green64/Library/Weights/Part040.lean}{\nolinkurl{Weights/Part040}}.
The exceptional residual estimates are in
\href{../Green64/Library/Weights/Part035.lean}{\nolinkurl{Weights/Part035}}
and
\href{../Green64/Library/Weights/Part039.lean}{\nolinkurl{Weights/Part039}}.
The long-prime statement (U1) and its shift are
{\small\nolinkurl{sourceCorrelationLongPrime_shifted_U1}} in
\href{../Green64/Library/Correlation/Part002.lean}{\nolinkurl{Correlation/Part002}}.
The contour bound (U4), its closure (U3)--(U5), and the matching
zero-free interface are
{\small\nolinkurl{p1p2_canonical_source_prime_bound}},
{\small\nolinkurl{sourceLongPrimeU1_of_sourceKhaleZeroFree}}, and
{\small\nolinkurl{KhalePublished11}} in
\href{../Green64/Library/Analysis/Estimates/Part003.lean}{\nolinkurl{Estimates/Part003}}.
The unconditional declaration there,
{\small\nolinkurl{sourceLongPrimeU1_unconditional}}, obtains (U2) from
the included stronger theorem
{\small\nolinkurl{OAI.DirichletCharacter.LFunction_ne_zero_of_seven_eighths_lt_re}}
in
\href{../Green64/Library/Analysis/NumberTheory/DirichletL/Detector/Part016.lean}{\nolinkurl{Detector/Part016}}.
The text invokes the weaker, independently published statement (U2); no
stronger zero-free conclusion is needed for the contour proof. The use
of (U1) on occupied intervals is in
\href{../Green64/Library/Analysis/Estimates/Part006.lean}{\nolinkurl{Analysis/Estimates/Part006}}.
The final four-term estimate is
{\small\nolinkurl{eventually_source_original_full_additive434}} in
\href{../Green64/Library/Correlation/Part033.lean}{\nolinkurl{Correlation/Part033}}.

\Needspace{10\baselineskip}

\hypertarget{congruences-source-correspondence}{%
\subsection{Distribution in residue
classes}\label{congruences-source-correspondence}}

\begin{longtable}[]{@{}
  >{\raggedright\arraybackslash}p{(\columnwidth - 2\tabcolsep) * \real{0.5000}}
  >{\raggedright\arraybackslash}p{(\columnwidth - 2\tabcolsep) * \real{0.5000}}@{}}
\toprule\noalign{}
\begin{minipage}[b]{\linewidth}\raggedright
Argument
\end{minipage} & \begin{minipage}[b]{\linewidth}\raggedright
Lean source and declaration
\end{minipage} \\
\midrule\noalign{}
\endhead
\bottomrule\noalign{}
\endlastfoot
Orthogonality with the nonunit correction &
\href{../Green64/Library/Weights/Part108.lean}{\nolinkurl{Weights/Part108}}:
{\small\nolinkurl{source_unit_class_remainder_identity}},
{\small\nolinkurl{source_shift_two_genuine_remainder_bound}} \\
Squarefree totient sum &
\href{../Green64/Library/Weights/Part107.lean}{\nolinkurl{Weights/Part107}}:
{\small\nolinkurl{source_squarefree_totient_sum_harmonic}} \\
Rough-support correction &
\href{../Green64/Library/Weights/Part108.lean}{\nolinkurl{Weights/Part108}}:
{\small\nolinkurl{shift_two_nonunit_reciprocal_harmonic}},
{\small\nolinkurl{shift_two_prime_reciprocal_rough_bound}} \\
Deleted-slot convolution and both masks &
\href{../Green64/Library/Sieve/Foundations/Part005.lean}{\nolinkurl{Foundations/Part005}}:
{\small\nolinkurl{source_original_moving_deleted_prime_convolution}},
{\small\nolinkurl{source_original_conductor_masked_convolution}} \\
Small conductors and final parameter choice &
\href{../Green64/Library/Weights/Part108.lean}{\nolinkurl{Weights/Part108}}:
{\small\nolinkurl{source_original_SW_finite_Echi_eventually}},
{\small\nolinkurl{source_shift_two_congruence_input_unconditional}} \\
Large-conductor bilinear estimate &
\href{../Green64/Library/Weights/Part107.lean}{\nolinkurl{Weights/Part107}},
with its dyadic expansion in
\href{../Green64/Library/Weights/Part108.lean}{\nolinkurl{Part108}} \\
Prime Siegel--Walfisz interface &
\href{../Green64/Library/Analysis/Estimates/Part005.lean}{\nolinkurl{Analysis/Estimates/Part005}}:
{\small\nolinkurl{sourcePrimeSiegelWalfisz_unconditional}} \\
Multiplicative large sieve &
\href{../Green64/Library/Sieve/Foundation/LargeSieve/Part001.lean}{\nolinkurl{Foundation/LargeSieve/Part001}} \\
Specialization used in the final argument &
\href{../Green64/Library/Sieve/Foundations/Part053.lean}{\nolinkurl{Foundations/Part053}}:
{\small\nolinkurl{parity_original_negative_congruence_D230_bound434}} \\
\end{longtable}

\Needspace{10\baselineskip}

\hypertarget{block-sieve-source-correspondence}{%
\subsection{The quantitative finite block
sieve}\label{block-sieve-source-correspondence}}

The statement and proof are {\small\nolinkurl{weighted_block_sieve}} in
\href{../Green64/Library/Sieve/Foundations/Part001.lean}{\nolinkurl{Foundations/Part001}}.
The shifted constant and depth estimates are in
\href{../Green64/Library/Sieve/Foundations/Part053.lean}{\nolinkurl{Foundations/Part053}}
and
\href{../Green64/Library/Sieve/Foundations/Part054.lean}{\nolinkurl{Part054}}.
The ordinary-rough constant is {\small\nolinkurl{sieveC}} in
\href{../Green64/Library/Analysis/Estimates/Part003.lean}{\nolinkurl{Analysis/Estimates/Part003}}.

\Needspace{10\baselineskip}

\hypertarget{rectangles-source-correspondence}{%
\subsection{A positive estimate on factor
rectangles}\label{rectangles-source-correspondence}}

\begin{longtable}[]{@{}
  >{\raggedright\arraybackslash}p{(\columnwidth - 2\tabcolsep) * \real{0.5000}}
  >{\raggedright\arraybackslash}p{(\columnwidth - 2\tabcolsep) * \real{0.5000}}@{}}
\toprule\noalign{}
\begin{minipage}[b]{\linewidth}\raggedright
Argument
\end{minipage} & \begin{minipage}[b]{\linewidth}\raggedright
Lean source and declaration
\end{minipage} \\
\midrule\noalign{}
\endhead
\bottomrule\noalign{}
\endlastfoot
Assignment-mask majorant &
\href{../Green64/Library/Sieve/Foundations/Part056.lean}{\nolinkurl{Foundations/Part056}}:
{\small\nolinkurl{parity_original_A0_assignment_mask_majorant434}} \\
Assignment moments, weighted-fibre reduction, sieve application &
\href{../Green64/Library/Sieve/Foundations/Part056.lean}{\nolinkurl{Foundations/Part056}} \\
CRT density and boundary count &
\href{../Green64/Library/Sieve/Foundations/Part002.lean}{\nolinkurl{Foundations/Part002}},
\href{../Green64/Library/Weights/Part035.lean}{\nolinkurl{Weights/Part035}} \\
Finite Brun--Hooley block sieve &
\href{../Green64/Library/Sieve/Foundations/Part001.lean}{\nolinkurl{Foundations/Part001}} \\
Euler-product comparison and mask average &
\href{../Green64/Library/Sieve/Foundations/Part057.lean}{\nolinkurl{Foundations/Part057}} \\
Complete finite bound (R1) &
\href{../Green64/Library/Sieve/Foundations/Part057.lean}{\nolinkurl{Foundations/Part057}}:
{\small\nolinkurl{parity_original_literal_weighted_box_normalized434}} \\
Error exponents and uniform normalization &
\href{../Green64/Sieve/BoxRemainder.lean}{\nolinkurl{BoxRemainder}},
\href{../Green64/Sieve/BoxError.lean}{\nolinkurl{BoxError}},
\href{../Green64/Sieve/BoxNormalization.lean}{\nolinkurl{BoxNormalization}} \\
Uniform final statement &
\href{../Green64/Sieve/WeightedBoxes.lean}{\nolinkurl{WeightedBoxes}}:
{\small\nolinkurl{weighted_box_bound_uniform}} \\
\end{longtable}

\Needspace{10\baselineskip}

\hypertarget{sieve-source-correspondence}{%
\subsection{Sieve estimates and the remaining
composites}\label{sieve-source-correspondence}}

\begin{longtable}[]{@{}
  >{\raggedright\arraybackslash}p{(\columnwidth - 2\tabcolsep) * \real{0.5000}}
  >{\raggedright\arraybackslash}p{(\columnwidth - 2\tabcolsep) * \real{0.5000}}@{}}
\toprule\noalign{}
\begin{minipage}[b]{\linewidth}\raggedright
Text
\end{minipage} & \begin{minipage}[b]{\linewidth}\raggedright
Definitions and proved identities
\end{minipage} \\
\midrule\noalign{}
\endhead
\bottomrule\noalign{}
\endlastfoot
Singular product and (S1) &
\href{../Green64/Library/Sieve/Foundations/Part006.lean}{\nolinkurl{Foundations/Part006}}:
{\small\nolinkurl{singularSeries_converges_and_ge_one}}; the congruence
source map is in
\protect\hyperlink{congruences-source-correspondence}{Section~\ref{congruences-source-correspondence}} \\
(S6), balanced rectangle count &
\href{../Green64/Sieve/BalancedBoxes.lean}{\nolinkurl{BalancedBoxes}}:
{\small\nolinkurl{balanced_rough_pair_bound}} \\
(S7)--(S8d), rough density and the full-cell estimate &
\href{../Green64/Library/Analysis/Estimates/Part003.lean}{\nolinkurl{Estimates/Part003}}:
{\small\nolinkurl{sourceDensity}} (the displayed index is \(k=n+2\));
\href{../Green64/Library/Weights/Part037.lean}{\nolinkurl{Weights/Part037}}:
{\small\nolinkurl{sourceD29_original_full_proxy_pointwise_compact}},
{\small\nolinkurl{literal_density_rough_unconditional}},
{\small\nolinkurl{literal_small_b_exp_bound_unconditional}} \\
Conditional Euler limit and (S9) &
\href{../Green64/Library/Sieve/Foundations/Part055.lean}{\nolinkurl{Foundations/Part055}}:
{\small\nolinkurl{parity_original_negative_conditional_natural_euler_limit434}},
{\small\nolinkurl{negative_prime_or_composite_mass_lower_bound}} \\
(S10), composite cover and two proxy errors &
\href{../Green64/Library/Sieve/Foundations/Part052.lean}{\nolinkurl{Foundations/Part052}}:
{\small\nolinkurl{parity_eventually_original_balanced_composite_prime_pair434}};
\href{../Green64/Sieve/ApproximationError.lean}{\nolinkurl{ApproximationError}}:
{\small\nolinkurl{balancedTwoStepError}};
\href{../Green64/Sieve/BalancedMass.lean}{\nolinkurl{BalancedMass}} \\
\end{longtable}

The one-bin specialization for (S10) is in
\href{../Green64/Sieve/CompositeMass.lean}{\nolinkurl{CompositeMass}}.


\clearpage
\begin{thebibliography}{9}
\bibitem{khale}
Tanmay Khale, \emph{An explicit Vinogradov--Korobov zero-free region for
Dirichlet $L$-functions}, Theorem~1.1,
\url{https://arxiv.org/abs/2210.06457}.
Published in \emph{The Quarterly Journal of Mathematics} 75 (2024), no.~1,
\url{https://doi.org/10.1093/qmath/haae010}.
\bibitem{green}
Ben Green, \emph{100 Open Problems}, December 2025 revision,
Problem 64, p.~31.
\url{https://people.maths.ox.ac.uk/greenbj/papers/open-problems.pdf}.
\bibitem{guo}
Jizhou Guo, \emph{Shifted primes with a prescribed parity of the number of
prime factors}, October 2026. Retained TeX source from
\url{https://github.com/aster2024/openai-math-followups}, revision
\nolinkurl{563cece082b628dfe3ba0290c9e4e30b92ffeca6},
\nolinkurl{shifted-prime-parity/source/main.tex}.
The source identity and availability are recorded in the accompanying
\nolinkurl{ACKNOWLEDGEMENTS.md}.
\bibitem{openai-e}
OpenAI, \emph{Prime Predecessors with an Even Number of Prime Factors},
manuscript dated 17 September 2026, \url{https://github.com/openai/math},
revision \nolinkurl{adc7f1241b42e322a6451854ab7e4b4c146bf78a}.
\bibitem{openai-s}
OpenAI, \emph{Weighted dilation graphs, smooth shifted primes and totient
fibers}, manuscript dated 24 September 2026,
\url{https://github.com/openai/math},
revision \nolinkurl{adc7f1241b42e322a6451854ab7e4b4c146bf78a}.
\end{thebibliography}
\end{document}
